\chapter[The Deligne series I: presentations, carriers, divisor comparison]{The Deligne series I: presentations, carriers and the divisor comparison}\label{ch:L24}
\begingroup\setlength{\emergencystretch}{3em}\sloppy

\providecommand{\ltwentyfourg}{\mathfrak g}
\providecommand{\ltwentyfourJ}{J_\infty}
\providecommand{\ltwentyfourgr}{\operatorname{gr}^{\mathrm{Li}}}
\providecommand{\ltwentyfourIns}{I_2^{\mathrm{ns}}}
\providecommand{\ltwentyfourRt}{\widetilde R}
\providecommand{\ltwentyfourHilb}{\operatorname{Hilb}}
\providecommand{\ltwentyfourch}{\operatorname{ch}}
\providecommand{\ltwentyfourvac}{\mathbf 1}

This chapter was written against an earlier state of the LeanVOA library, and every
Lean quotation below was checked (with \texttt{sync.py check}) to occur verbatim in the source files of that state. For this edition, \texttt{scripts/textbook/sync.py} checks that every Lean quotation below occurs verbatim (up to whitespace and comments) in the source files of that release. The
machine record is older than that release, and we state it exactly. The umbrella target \lean{VOA.Deligne} was last
built green at an earlier commit (11099 jobs), with every \texttt{\#print axioms} in the campaign files a subset of
\lean{[propext, Classical.choice, Quot.sound]}. The Lean sources have changed since then: certificate files were split
to fit the memory budget (among them the $A_2$ divisor certificate behind \lean{A2.divisorComparison}), and the
unconditional $A_2$ theorem of Chapter~\ref{ch:L27} was added. The project records that the exported statements of the
split files were preserved, and the quotations here confirm it for every statement this chapter prints. The changed files
were built green in targeted builds, and the aggregate build \texttt{bash scripts/campaign\_build.sh VOA} (the full \texttt{lake build VOA} under the project's memory-capped wrapper) was green on 2026-10-28 (a pre-release build): 18{,}901 jobs, and again 18{,}907 jobs after merging the then-current \texttt{origin/main}; the aggregate target contains the umbrella \lean{VOA.Deligne}. Section~\ref{l24:sec:verification} gives the details.

The chapter covers the 1195 Lean files that the book assigns to it
(\texttt{Notes/Textbook/assignment.txt}): roughly 170 files of genuine mathematics, which we present
declaration by declaration, and roughly 1020 files of certificate data and certificate checks (the
$E_8$ divisor comparison alone is 619 files), which we describe collectively in
Section~\ref{l24:sec:certs}: what each family checks, in what format, by which kernel mechanism, and at
what size.

\paragraph{Conventions for quotations and status boxes.}
\begin{itemize}
\item \emph{Omitted proofs.} A Lean quotation shows a statement, not its proof. A block that ends in
  \lean{:=} or \lean{:= by} with nothing after it is a statement whose proof is omitted; it is not a truncated
  quotation. A few definitions are shown with their bodies omitted in the same way; the dictionary says so
  each time.
\item \emph{``Axioms standard''.} This means that \texttt{\#print axioms} reports a subset of
  \lean{[propext, Classical.choice, Quot.sound]}, the three standard axioms of Lean's logic (propositional
  extensionality, the axiom of choice, quotient soundness); see Chapter~\ref{ch:01},
  Section~\ref{ch01:sec:trust}. It does \emph{not} mean ``unconditional'': \texttt{\#print axioms} does not
  list the explicit hypotheses of a theorem (Section~\ref{ch01:sec:uncond}). Every status box therefore names
  the hypotheses separately.
\item \emph{Where a status claim comes from.} Unless a status box says otherwise, an axiom or build claim
  is a \emph{campaign record}: it was recorded by the project's own builds and audits (the earlier umbrella
  build, the later targeted builds, and the \texttt{\#print axioms} sections at the end of each Lean file). A claim marked
  \emph{re-checked} was re-run by the author of this chapter on a scratch file; the list is in
  Section~\ref{l24:sec:verification}.
\end{itemize}

\paragraph{Prerequisites.} Three things are used without being re-developed here.
(1) Li's filtration $F^\bullet V$ of a VOA, its layers $\ltwentyfourgr_{p,N}V=F^pV_N/F^{p+1}V_N$ indexed by
the \emph{Li degree} $p$ and the conformal weight $N$, and the identification $\ltwentyfourgr_0V=R_V$ with
Zhu's $C_2$-algebra (Chapter~\ref{ch:L11}). (2) The presented jet algebra $\ltwentyfourJ(\C[x]/I)$ of a
polynomial quotient, with its bigrading by the number of derivatives and the conformal weight, and its universal
property as an arc algebra (Chapters~\ref{ch:L11} and \ref{ch:L20}). (3) The affine VOAs $V^k(\ltwentyfourg)$,
their simple quotients $L_k(\ltwentyfourg)$, and the concrete exceptional Lie algebras (Chapters~\ref{ch:L03},
\ref{ch:L04} and \ref{ch:L23}). For precise pointers: the Lean objects \lean{liWeightLayer W p N},
\lean{liAssociatedGraded W} and \lean{liPresentedJetMap} are in Chapter~\ref{ch:L11},
Section~\ref{l11:sec:jets} (Definitions~\ref{l11:def:jets} and \ref{l11:def:jetmap});
\lean{InfiniteJet.PresentedJet} and its arc universal property are Definition~\ref{l20:def:jet} and
Theorem~\ref{l20:thm:arc}; the Lie data \lean{LieData} are in Chapter~\ref{ch:L23}. Chapter~\ref{ch:L20} is the
closest model: it proves classical freeness for $L_k(\mathfrak{sl}_2)$ at the boundary admissible levels by a
squeeze of exactly the kind used here, and two of its collision-divisibility files
(\leanpath{VOA/Commutative/CollisionCommonFactors.lean},
\leanpath{VOA/Commutative/RootedGordonConductorDivisibility.lean}) are reused by the Deligne carrier.

\paragraph{What this chapter is about, in one paragraph.} The \emph{Deligne series} is the list of
simple Lie algebras $A_1\subset A_2\subset G_2\subset D_4\subset F_4\subset E_6\subset E_7\subset E_8$,
taken at the level $k=-1-h^\vee/6$. For $V=L_k(\ltwentyfourg)$ one wants to prove that Li's canonical
surjection $\ltwentyfourJ R_V\twoheadrightarrow\ltwentyfourgr V$ is injective (\emph{classical freeness}), and that
$R_V=\operatorname{Sym}\ltwentyfourg/(\ltwentyfourIns)$ for an explicit space $\ltwentyfourIns$ of
quadrics. The paper proof is a chain of coefficientwise inequalities between $q$-series, from the
character of $V$ through the Hilbert series of an explicit ``carrier'' down to a sum of affine Demazure
characters, closed by a bound of that sum by the character. This chapter formalizes (i) the reduction of classical
freeness to that chain (a conditional capstone with four packages), (ii) the explicit presentation
$\ltwentyfourIns\subseteq\ker$ for all seven higher-rank types, (iii) current-jet surjectivity, (iv)
the geometric carrier, and (v) the carrier half of the chain, steps B1--B8 and the resummation D3, with
type-generic theorems instantiated by kernel-checked certificates for $A_2$, $G_2$, $D_4$, $F_4$,
$E_6$, $E_7$, $E_8$. The character side is Chapter~\ref{ch:L25}, the fusion inputs U5 and U6 are
Chapter~\ref{ch:L26}, and the final per-type theorems (unconditional for $A_2$, conditional on named inputs for the other six types) are
Chapter~\ref{ch:L27}.

\paragraph{Why formalize this.} Classical freeness of $L_k(\ltwentyfourg)$ is a statement with direct physical
content: it says that the Schur index of the corresponding four-dimensional $\mathcal N=2$ theory equals the Hilbert
series of the arc space of a (non-reduced) Higgs branch, so that every Schur operator is a classical field
configuration. Its paper proof is long: about fifteen labelled steps, most of them commutative algebra, a few of
them citations of the literature, and one of them (the divisor comparison, below) a finite but very large
computation in $\ltwentyfourg^*\otimes\ltwentyfourg^*$. Each step is a place where a gap could hide. The Lean
development separates the steps that can be \emph{proved} (all the commutative algebra, and the divisor comparison
by certificates) from the genuine literature inputs, which survive as named hypotheses (U5, U6) or are deferred to
another chapter (the character bound $(\star)$). The divisor comparison gives the chapter its title because it is the
one step whose paper proof is a computation rather than an argument; the reason it is the only such step is explained
at the start of Section~\ref{l24:sec:sector}.

\begin{remark}[The logical arc and the ledger]\label{l24:rem:arc}
Write $\mathrm{carrierDim}(N)$ for the $q^N$-coefficient of $\ltwentyfourHilb\ltwentyfourJ\ltwentyfourRt$, the
Hilbert series of the arc algebra of $\ltwentyfourRt=\operatorname{Sym}\ltwentyfourg/(\ltwentyfourIns)$ (by
Theorem~\ref{l24:thm:B4} this is also the total dimension of the weight-$N$ carrier), and $\mathcal D_N$ for the
$q^N$-coefficient of the \emph{Demazure sum} $\sum_{n\ge0}q^n\,\ltwentyfourHilb_qD(1,n\theta)/(q)_n$ (notation below).
The whole Deligne argument is three statements, for every conformal weight $N$:
\begin{align*}
\text{(A)}\quad & \bigl(\mathrm{carrierDim}(N)\le\dim V_N\ \text{for all }N\bigr)\ \Longrightarrow\ V\text{ classically free, and }
\ker(\operatorname{Sym}\ltwentyfourg\to R_V)=(\ltwentyfourIns);\\
\text{(B)}\quad & \mathrm{carrierDim}(N)\ \le\ \mathcal D_N;\\
\text{(C)}\quad & \mathcal D_N\ \le\ \dim V_N .
\end{align*}
(A) is Theorem~\ref{l24:thm:upper}, together with current-jet surjectivity (Theorem~\ref{l24:thm:jetsurj}); it is
proved here with no further hypothesis. (B) is the carrier half of the chain, proved here conditionally on U5 and U6
(Theorems~\ref{l24:thm:carrierside} and \ref{l24:thm:seven}). (C) is the character side $(\star)$, treated in
Chapter~\ref{ch:L25}. (B) and (C) together give the hypothesis of (A). Since the opposite inequality
$\dim V_N\le\mathrm{carrierDim}(N)$ is automatic (Li's map is onto), all three quantities then coincide.

The ledger, step by step (``form'' says whether the step is used as an equality, an inequality or a containment):
\begin{center}\small
\begin{tabular}{p{0.30\textwidth}p{0.17\textwidth}p{0.42\textwidth}}
\toprule
step & form & status in this chapter\\
\midrule
U2: $\ltwentyfourIns\subseteq\ker$ & containment & proved, all seven types, no hypothesis\\
exactness $\ker=(\ltwentyfourIns)$ & equality & a \emph{conclusion} of (A), never assumed\\
current-jet surjectivity & surjection & proved, all seven types\\
B4: carrier $=$ dual of the jet algebra & isomorphism & proved, any quadrics\\
U3: $X_aQ\in(\ltwentyfourIns)$ & containment & certified per type\\
B1, B2--B3 & $\le$ and $=$ & proved, any type data\\
B5, generic comparison & containments & proved, any sector data\\
divisor comparison & containment & proved per type by certificates\\
B6, B8 & freeness, equality & proved from U5 and the divisor comparison\\
D3 & equality & proved (pure $q$-series)\\
U5 (fusion bound) & $\le$ & \textbf{named hypothesis}; reduced in Chapter~\ref{ch:L26}; discharged for $A_2$
  (Chapter~\ref{ch:L27})\\
U6 (mixed recursion) & equality of series & \textbf{named hypothesis} here; proved by construction in
  Chapter~\ref{ch:L26} (\lean{fusionOf})\\
$(\star)$: (C) above & $\le$ & Chapter~\ref{ch:L25}\\
final per-type theorems; $A_2$ unconditional & & Chapter~\ref{ch:L27}\\
\bottomrule
\end{tabular}
\end{center}
The paper states $(\star)$ as an equality (``Demazure sum $=$ character''); only the inequality (C) is used, and that
is the form in which the Lean development states it (\lean{DemazureBoundFor}, Chapter~\ref{ch:L25}).
Chapter~\ref{ch:L27} renames the remaining inputs: Chapter~\ref{ch:L26} reduces U5 to the total-dimension bound,
which is Input~T, and the Demazure realizations U7, which are Input~R; the star bound $\mathcal D_N\le c_N$ of
Chapter~\ref{ch:L25} is Input~S.
\end{remark}

\paragraph{The findings of this chapter, in brief.}
\begin{enumerate}
\item \textbf{Exactness of the $C_2$ presentation is a conclusion, not a hypothesis.} The capstone
  \lean{caseClassicallyFree_of_explicit_inputs} assumes only that the displayed ideal is
  \emph{contained} in the kernel of $\operatorname{Sym}\ltwentyfourg\to R_V$, and concludes that it
  \emph{equals} the kernel, through a short retraction argument
  (\lean{VOA.ker_eq_of_liPresentationClassicallyFree}, Theorem~\ref{l24:thm:kereq}).
\item \textbf{The explicit presentation is unconditional for all seven types.} For each type, a
  weight-two singular vector dies in $L_k(\ltwentyfourg)$, and coadjoint propagation carries its
  symbol to every relation of $\ltwentyfourIns$ (8, 27, 105, 324, 650, 1539 and 3875 relations). The
  resulting presentations \lean{explicitPresentation*} and \lean{I2ns_le_ker*} have no hypotheses
  (Theorem~\ref{l24:thm:sevenpres}).
\item \textbf{Three of the capstone's four packages are proved.} Package~1 (the presentation) is the explicit
  presentation above. Package~2, current-jet surjectivity, holds for every Deligne type
  (Theorem~\ref{l24:thm:jetsurj}). Package~3, the geometric carrier embedding, holds for every presentation by
  quadrics, and is in fact a linear \emph{equivalence} with a literal polynomial carrier
  (Theorem~\ref{l24:thm:geom}); on its own it is weak (Warning in Section~\ref{l24:sec:carrier}). Package~4, the
  character identity, is not proved in this chapter: it is equivalent to the inequality (A) above, which is what (B)
  and (C) supply, and (B) rests on U5 and U6. For $A_2$ the remaining inputs are proved, and Chapter~\ref{ch:L27} states
  the resulting unconditional theorem.
\item \textbf{The carrier chain is formalized type-generically.} B1 (the Casimir layer bound) and D3
  (the Gaussian resummation) are proved with no hypothesis; B5 ($M_m\subseteq W_m$), B6 (freeness of
  $M_m$, from U5), the generic comparison ($\Delta^L W_m\subseteq M_m$) and B8 (normality) are proved
  for arbitrary sector data; the divisor comparison (B7 at $\Delta=0$) is proved for each of the seven
  types by integer certificates. The output is, for each type, the sector fusion bound from U5 alone
  (\lean{X_sectorFusionBound_of_U5}; for $A_2$ see Theorem~\ref{l24:thm:seven}), and then
  $\sum_p\dim\mathrm{Carrier}_{p,N}\le\mathcal D_N$ from U5 and U6.
\item \textbf{What this chapter does not prove.} It proves no character equality, and no classical
  freeness theorem for any higher-rank Deligne VOA. The fusion bound U5 (total dimension: Kus--Littelmann, Pacific J.\ Math.\ 278, Thm.~4.2, with one
  flagged step; Feigin's earlier proof has a gap; see Chapter~\ref{ch:L27}) and U6
  (Ravinder's mixed recursion) remain named hypotheses here; Chapter~\ref{ch:L26} reduces them further.
  For $A_2$ every input is now a theorem, and the unconditional statement that
  $L_{-3/2}(\mathfrak{sl}_3)$ is classically free with $R_V=\operatorname{Sym}\mathfrak{sl}_3/(\ltwentyfourIns)$,
  \lean{VOA.DeligneAudit.A2.classicallyFree}, is presented in Chapter~\ref{ch:L27}. For the other six types the
  final theorems remain conditional (Chapter~\ref{ch:L27}).
\item \textbf{Size.} The certificate layer is large ($E_8$: 61\,504 two-slot tables, 9291 kernel theorems in 556 files,
  10.8\,MB of literal data); its sizes and measured costs are in Section~\ref{l24:sec:certs}.
\end{enumerate}

\paragraph{Notation used in this chapter.} Several letters are used by the Lean sources for more than one object.
In the mathematical text we use the following, and we keep the Lean names only inside quotations and
dictionaries.
\begin{center}\small
\begin{longtable}{p{0.22\textwidth}p{0.70\textwidth}}
\toprule
symbol & meaning\\
\midrule\endhead
$q$ & the formal variable of all $q$-series (conformal weight, or $t$-degree). In the $A_2$ file the Lean
  abbreviation \lean{q a b} means the monomial $X_aX_b$; it never denotes the variable.\\
$[q^N]f$ & the coefficient of $q^N$ in the power series $f$.\\
$(q)_n$, $\begin{bmatrix}n\\d\end{bmatrix}_q$ & $(q)_n=\prod_{i=1}^n(1-q^i)$; the Gaussian binomial
  $(q)_n/((q)_d(q)_{n-d})$.\\
$\ltwentyfourch V$ & the character (graded dimension) $\sum_N\dim V_N\,q^N$ of a VOA graded by conformal weight.\\
$\ltwentyfourHilb A$ & the Hilbert series $\sum_N\dim A_N\,q^N$ of a graded algebra or space; for jet algebras a
  current has weight $1$ and each derivative adds $1$.\\
$(p,N)$ & Li bidegree: $p$ is the Li degree (number of derivatives), $N$ the conformal weight.\\
$\ltwentyfourg$, $n$, $\theta$, $h^\vee$, $k$ & the simple Lie algebra, $n=\dim\ltwentyfourg$, the highest root, the dual Coxeter number,
  the level $k=-1-h^\vee/6$; $V(\lambda)$ is the irreducible $\ltwentyfourg$-module of highest weight $\lambda$.\\
$Q$, $\ltwentyfourIns$, $I_2$ & the Casimir quadric in $\operatorname{Sym}^2\ltwentyfourg$; the non-singlet weight-two relations
  (the superscript ``ns'' reads ``non-singlet'': $I_2$ with its singlet $\C Q$ removed); $I_2=\ltwentyfourIns\oplus\C Q$.\\
$\ltwentyfourRt$, $R_V$, $\ltwentyfourJ$ & $\operatorname{Sym}\ltwentyfourg/(\ltwentyfourIns)$; Zhu's $C_2$-algebra; the arc
  (infinite jet) algebra functor.\\
$\mathcal R=(q_r)_{r\in\rho}$ & a family of quadric coefficient matrices (Lean's variable \lean{Q}); $\mathrm{quad}(q_r)$ the quadric.\\
$W$; $W_m$ & a VOA (no subscript); the current carrier with $m$ particles (Definition~\ref{l24:def:carrier}).\\
$P_m$, $B_m$, $B_+$ & $\C[z_1,\dots,z_m]$; its symmetric invariants $P_m^{S_m}$; the ideal of $B_m$ of symmetric polynomials
  without constant term.\\
$M_m$, $\delta_{c_0}$, $\delta_t$ & the cyclic module generated by $\delta_{c_0}^{\otimes m}$ (Definition~\ref{l24:def:Mm}); the
  highest-weight vector of $\ltwentyfourg^*$; the basis tensor at a labelling $t$.\\
$\Delta$ & the discriminant $\prod_{i\neq j}(z_i-z_j)\in B_m$ (not to be confused with $\delta_{c_0}$).\\
$L$ & the cyclicity depth of the sector data (Definition~\ref{l24:def:sectordata}).\\
$\mathcal F_m$; $F=(F_m)_m$ & the fusion module (Definition~\ref{l24:def:U5}); a family of power series (Lean \lean{F}),
  whose intended value is the graded character of $\mathcal F_m$, for instance $F_1(q)=n$ and
  $F_2(q)=\dim V(2\theta)+q\dim\Lambda^2\ltwentyfourg+q^2\dim I_2$ (for $A_2$: $27+28q+9q^2$, total $64=8^2$).\\
$W_{\mathrm{loc}}(m\theta)$ & the local Weyl module of the current algebra $\ltwentyfourg[t]$ (Definition~\ref{l24:def:U5}).\\
$D(1,n\theta)$, $\mathrm{Dem}_n$ & the level-one affine Demazure module of highest weight $n\theta$, a $\ltwentyfourg[t]$-module;
  $\mathrm{Dem}_n=\ltwentyfourHilb_qD(1,n\theta)$ its $t$-graded dimension (Chapter~\ref{ch:L25}).\\
$\mathcal D_N$ & the $q^N$-coefficient of the Demazure sum $\sum_nq^n\mathrm{Dem}_n/(q)_n$.\\
$\mathscr C$, $S$ & carrier data (Lean \lean{D : DeligneCarrierData}); sector data (Lean \lean{S : DeligneSectorData}).\\
$\lambda_i(t)$, $\lambda_i$ & $\prod_{j\neq i}(t-z_j)\in P_m[t]$ and its value $\lambda_i(z_i)=\prod_{j\neq i}(z_i-z_j)$ (Lean
  \lean{Lpoly i}, \lean{Dslot i}).\\
$G_i$, $H$, $\Lambda_{ab}$, $E_{ab}$ & for slots $a\neq b$: $G_i=\prod_{k\notin\{a,b\}}(z_i-z_k)$, $H=G_aG_b$,
  $\Lambda_{ab}=\prod_{j\notin\{a,b\}}\lambda_j^L$ (Lean \lean{Fspec}), $E_{ab}=\prod(z_i-z_j)$ over ordered pairs $\{i,j\}\neq\{a,b\}$.\\
$\nu$, $\ell$ & in Section~\ref{l24:sec:divisor}: the number of pair-operator steps (Lean \lean{D}) and the collision order,
  the power of $z_a-z_b$ (the same letter in Lean, where it is called ``level'' in the name \lean{LevOK}; unrelated to the affine level $k$
  and to the Li degree).\\
$a_r(c,d)$ & integer coefficient matrices of the relations of $I_2$ (Lean \lean{qI}).\\
$\pi$; $\varpi$ & Li's surjection $\ltwentyfourJ R_V\to\ltwentyfourgr V$; a prime element of $B_m$.\\
\bottomrule
\end{longtable}
\end{center}

%======================================================================
\section{The mathematics}\label{l24:sec:math}
%======================================================================

\subsection{The Deligne series}

Deligne observed (1996) that the exceptional Lie algebras, together with a few classical ones, behave
as a one-parameter family: many representation-theoretic quantities, for instance the decomposition of
$\operatorname{Sym}^2\ltwentyfourg$ and of $\Lambda^2\ltwentyfourg$, are given by rational functions of
the dual Coxeter number $h^\vee$. In the long-root-square-two normalization the family and its levels
are:
\begin{center}\small
\begin{tabular}{lcccccl}
\toprule
type & $\dim\ltwentyfourg$ & $h^\vee$ & $k=-1-h^\vee/6$ & $\dim V(2\theta)$ & $\dim\ltwentyfourIns$ & $\ltwentyfourIns$ as a $\ltwentyfourg$-module\\
\midrule
$A_1$ & 3 & 2 & $-4/3$ & 5 & -- & (see Chapter~\ref{ch:L20}) \\
$A_2$ & 8 & 3 & $-3/2$ & 27 & 8 & $V(\theta)$ (the adjoint)\\
$G_2$ & 14 & 4 & $-5/3$ & 77 & 27 & $V(2\omega_1)$\\
$D_4$ & 28 & 6 & $-2$ & 300 & 105 & $V(2\omega_1)\oplus V(2\omega_3)\oplus V(2\omega_4)$\\
$F_4$ & 52 & 9 & $-5/2$ & 1053 & 324 & $V(2\omega_4)$\\
$E_6$ & 78 & 12 & $-3$ & 2430 & 650 & $V(\omega_1+\omega_6)$\\
$E_7$ & 133 & 18 & $-4$ & 7371 & 1539 & $V(\omega_6)$\\
$E_8$ & 248 & 30 & $-6$ & 27000 & 3875 & $V(\omega_1)$\\
\bottomrule
\end{tabular}
\end{center}
In every row $\operatorname{Sym}^2\ltwentyfourg=V(2\theta)\oplus\ltwentyfourIns\oplus\C Q$, where $\theta$
is the highest root and $Q$ the Casimir quadric; for example $406=300+105+1$ for $D_4$ and
$30876=27000+3875+1$ for $E_8$. We write $I_2=\ltwentyfourIns\oplus\C Q$ for the complement of the top
summand $V(2\theta)$. The ideal $(I_2)\subset\operatorname{Sym}\ltwentyfourg$ is the classical shadow of the
\emph{Joseph ideal}: for $\ltwentyfourg$ not of type $A$, the Joseph ideal is the unique completely prime primitive ideal of
the enveloping algebra $U(\ltwentyfourg)$ whose associated variety is the closure $\overline{\mathbb O_{\min}}$ of the
minimal nilpotent orbit (the orbit of a highest root vector), and its associated graded ideal in
$\operatorname{Sym}\ltwentyfourg=\operatorname{gr}U(\ltwentyfourg)$ is generated by $I_2$. In particular
$\operatorname{Sym}\ltwentyfourg/(I_2)$ is the coordinate ring of $\overline{\mathbb O_{\min}}$ (Kostant; Garfinkle; see
the Notes). The ring that matters here is
the slightly larger
\begin{equation*}
\ltwentyfourRt=\operatorname{Sym}\ltwentyfourg/(\ltwentyfourIns),
\end{equation*}
in which $Q$ is not killed but is nilpotent: $XQ\in(\ltwentyfourIns)$ for every $X\in\ltwentyfourg$, hence
$Q^2\in(\ltwentyfourIns)$ (input U3 below). The member $A_1$ is treated separately (its relation module
is $\mathfrak{sl}_2Q$ in degree three); it is the $m=1$ case of Chapter~\ref{ch:L20}.

\begin{physics}
The Deligne series at $k=-1-h^\vee/6$ is the list of chiral algebras of the rank-one four-dimensional
$\mathcal N=2$ superconformal field theories with a flavour symmetry in the series, via the 4d/2d
correspondence of Beem, Lemos, Liendo, Peelaers, Rastelli and van Rees: the Argyres--Douglas theories
$H_1$, $H_2$ (flavour $A_1$, $A_2$), $SU(2)$ with four flavours (flavour $D_4$), and the
Minahan--Nemeschansky theories (flavour $E_6$, $E_7$, $E_8$). The 2d level is $k_{2d}=-k_{4d}/2$, and
$h^\vee+6$ appears because these theories saturate a unitarity bound on the flavour central charge. The
Higgs branch of each theory is $\overline{\mathbb O_{\min}}$, the one-instanton moduli space, which is the
associated variety of $L_k(\ltwentyfourg)$ (Arakawa--Moreau). $G_2$ and $F_4$ have no known
four-dimensional origin; mathematically they belong to the same family.
\end{physics}

\subsection{The associated scheme, nilpotence of the Casimir, and the target statement}

For a VOA $V$, Zhu's $C_2$-algebra is $R_V=V/C_2(V)$, a commutative Poisson algebra; Li's
filtration has $\ltwentyfourgr_0V=R_V$, and there is a canonical surjection of differential algebras
\begin{equation}\label{l24:eq:li}
\pi:\ \ltwentyfourJ R_V\twoheadrightarrow\ltwentyfourgr V
\end{equation}
(Li 2005; Arakawa 2012; Chapter~\ref{ch:L11}). $V$ is \emph{classically free} when $\pi$ is injective.
For an affine VOA $L_k(\ltwentyfourg)$ the weight-one currents $x_{(-1)}\ltwentyfourvac$, $x\in\ltwentyfourg$,
generate $R_V$, so $R_V=\operatorname{Sym}\ltwentyfourg/\ker$ for the kernel of the current-coordinate
map. Let $\ltwentyfourIns$ be the image in $\operatorname{Sym}^2\ltwentyfourg$ of the $C_2$-symbols of the
weight-two singular vectors of $V^k(\ltwentyfourg)$ (they exist at $k=-1-h^\vee/6$, by
Arakawa--Moreau and by Perše for $A_2$; we use only their existence at this level, not any claim that this is the
only level at which they occur). The target theorem, for the seven higher-rank types, is
\begin{equation}\label{l24:eq:target}
\ltwentyfourJ\bigl(\operatorname{Sym}\ltwentyfourg/(\ltwentyfourIns)\bigr)\xrightarrow{\ \sim\ }\ltwentyfourgr L_k(\ltwentyfourg),
\qquad\text{and}\qquad \ker\bigl(\operatorname{Sym}\ltwentyfourg\to R_{L_k(\ltwentyfourg)}\bigr)=(\ltwentyfourIns).
\end{equation}
The second statement says that $R_V$ is the non-reduced scheme $\ltwentyfourRt$, with $Q\neq0$ and $Q^2=0$.

\begin{physics}
The Casimir $Q$ is, up to the factor $2(k+h^\vee)$, the $C_2$-shadow of the Sugawara stress tensor
$T$. ``$Q$ nilpotent in $R_V$'' is the statement, conjectured by Beem and Rastelli for every chiral
algebra of a 4d $\mathcal N=2$ theory, that $T$ is nilpotent in the $C_2$-algebra: the Higgs branch is
the reduced associated variety, but the chiral algebra remembers a nilpotent thickening. Classical
freeness says that the Schur index of the 4d theory equals the Hilbert series of the arc space of this
thickened scheme: every Schur operator is a classical field configuration on the Higgs branch, with no
extra quantum relations.
\end{physics}

\subsection{The paper proof as a chain of inequalities}\label{l24:sec:chain-paper}

The proof audited in \texttt{Notes/deligne/audit\_2026\_10\_25/DeligneCampaignAudit.tex} runs, coefficient
by coefficient in $q$ (the conformal weight; notation as in the table of the introduction),
\begin{equation}\label{l24:eq:chain}
\begin{aligned}
\ltwentyfourch V\ &\overset{\mathrm{U1}}{\le}\ \ltwentyfourHilb\ltwentyfourJ R_V\ \overset{\mathrm{U2}}{\le}\ \ltwentyfourHilb\ltwentyfourJ\ltwentyfourRt
\ \overset{\mathrm{B1}}{\le}\ \ltwentyfourHilb U_{\ltwentyfourg}
\ \overset{\mathrm{B2\text{--}B8}}{\le}\ \sum_{m,d\ge0}\frac{q^{m+d+d(m+d)}\,F_m(q)}{(q)_m\,(q)_d}\\
&\overset{\mathrm{D3,\,U6}}{=}\ \sum_{n\ge0}\frac{q^n\,\ltwentyfourHilb_qD(1,n\theta)}{(q)_n}
\ \overset{(\star)}{\le}\ \ltwentyfourch V .
\end{aligned}
\end{equation}
Here $F_m(q)$ is the graded character of the fusion module $\mathcal F_m$ (B6 below), $\ltwentyfourHilb_qD(1,n\theta)$
the $t$-graded dimension of the level-one Demazure module of highest weight $n\theta$, and $U_{\ltwentyfourg}$ an
auxiliary algebra described in B1. The steps, each with a gloss (the precise statements come later):
\begin{itemize}
\item \textbf{U1} (Li): $\pi$ is surjective, and $\ltwentyfourHilb\ltwentyfourgr V=\ltwentyfourch V$ because the layers of Li's
  filtration of $V_N$ add up to $\dim V_N$. So the quantum count is at most the classical count.
\item \textbf{U2} (containment): $\ltwentyfourIns\subseteq\ker$, so $R_V$ is a quotient of $\ltwentyfourRt$ and $\ltwentyfourJ R_V$ a
  quotient of $\ltwentyfourJ\ltwentyfourRt$.
\item \textbf{U3}: $X\cdot Q\in(\ltwentyfourIns)$ in $\operatorname{Sym}^3\ltwentyfourg$ for every current $X$; hence $Q^2\in(\ltwentyfourIns)$.
\item \textbf{B1}: filter $\ltwentyfourJ\ltwentyfourRt$ by the powers of the ideal generated by $Q$ and its derivatives. The $d$-th
  layer is bounded by the states of an \emph{auxiliary algebra} $U_{\ltwentyfourg}$, a jet algebra with $m$ \emph{current particles}
  (weight one) and $d$ \emph{Casimir particles} (weight two) obeying collision rules. Its Hilbert series is a sum over
  \emph{sectors} $(m,d)$.
\item \textbf{B2--B3}: dually, a state of sector $(m,d)$ is a polynomial ``wave function'' of $m+d$ positions. It is divisible
  by a fixed polynomial $\delta_{m,d}$, the \emph{conductor}, vanishing to order one at every current--Casimir collision and
  to order two at every Casimir--Casimir collision, of degree $md+d(d-1)$; dividing by it (the \emph{conductor shift}) leaves
  $W_m\otimes\C[v_1,\dots,v_d]^{S_d}$. The weights add up to $(m+2d)+(md+d(d-1))=m+d+d(m+d)$, and the symmetric
  polynomials in $d$ variables contribute $1/(q)_d$.
\item \textbf{B4}: $W_m$ is the space of symmetric $\ltwentyfourg^{*\otimes m}$-valued polynomials in
  $z_1,\dots,z_m$ satisfying, at every collision $z_i=z_j$, the quadratic equations of $I_2$ (the
  ``current carrier''); it is exactly the dual of the $m$-particle part of the jet algebra.
\item \textbf{B5}: $M_m:=U(\ltwentyfourg[t])\cdot\delta_{c_0}^{\otimes m}\subseteq W_m$, where $\delta_{c_0}$
  is the highest-weight vector of $\ltwentyfourg^*$ (weight $\theta$), the functional dual to $f_\theta$.
\item \textbf{B6}: $M_m$ is a free module over $B_m=\C[z_1,\dots,z_m]^{S_m}$, with fibre $M_m/B_+M_m$ the fusion module
  $\mathcal F_m=W_{\mathrm{loc}}(m\theta)/\langle(f_\alpha\otimes t^m)w\rangle$. Since $\ltwentyfourHilb B_m=1/(q)_m$, this gives
  $\ltwentyfourHilb M_m=F_m(q)/(q)_m$.
\item \textbf{B7}: two-point equality $M_2=W_2$, $\ltwentyfourch\mathcal F_2=V(2\theta)+q\Lambda^2\ltwentyfourg+q^2I_2$.
\item \textbf{B8}: codimension-one comparison plus normality of $B_m$ give $M_m=W_m$, hence
  $\ltwentyfourHilb W_m=F_m(q)/(q)_m$.
\item \textbf{U4, U5}: Chari--Fourier--Khandai (global Weyl modules) and the fusion bound
  $\dim\mathcal F_m\le(\dim\ltwentyfourg)^m$, with graded character $F_m(q)$ (total dimension: Kus--Littelmann,
  Pacific J.\ Math.\ 278, Thm.~4.2, with one flagged step; Feigin's earlier proof has a gap; see Chapter~\ref{ch:L27}).
\item \textbf{U6, D3}: Ravinder's mixed recursion, which expresses $\ltwentyfourHilb_qD(1,n\theta)$ through the $F_m$, and the
  Gaussian resummation of the double sum over $(m,d)$ into the single sum over $n=m+d$.
\item \textbf{$(\star)$}: the paper's character identity ``Demazure sum $=$ character'' (Chapter~\ref{ch:L25}). Only the
  inequality ``Demazure sum $\le$ character'' is needed, and that is the form stated in Lean (\lean{DemazureBoundFor},
  Chapter~\ref{ch:L25}).
\end{itemize}
The chain begins and ends with $\ltwentyfourch V$. Every term is finite in each conformal weight, so once the chain closes,
every inequality is an equality weight by weight; equality in the second inequality, read in Li degree zero, is
$R_V=\ltwentyfourRt$. This is the squeeze of Remark~\ref{l24:rem:arc}: the Lean development proves the inequalities from
$\ltwentyfourHilb\ltwentyfourJ\ltwentyfourRt$ to the Demazure sum (conditionally on U5 and U6), and
Theorem~\ref{l24:thm:upper} turns ``$\ltwentyfourHilb\ltwentyfourJ\ltwentyfourRt\le\ltwentyfourch V$'' into classical freeness.

\begin{physics}
Each side of \eqref{l24:eq:chain} counts states of weight $N$ in a different description of the same
theory: the quantum chiral algebra, its classical limit, a ``free-particle'' description in which
$m$ current quanta and $d$ Casimir quanta move on the line with prescribed collision behaviour
(a Feigin--Stoyanovsky functional model), and finally a sum over $n$ of Demazure modules, i.e.\ the
fusion of $n$ adjoint representations. That the counts agree is a statement that the operator content
is captured exactly by particles with the right exclusion statistics.
\end{physics}

\subsection{Where each step lives}\label{l24:sec:map}

\begin{center}\small
\begin{tabular}{p{0.16\textwidth}p{0.47\textwidth}p{0.27\textwidth}}
\toprule
step & Lean (this chapter unless noted) & status\\
\midrule
U1, squeeze & \lean{classicallyFree_of_carrier_weightCharacter} & proved \\
U2 & \lean{I2ns_le_ker*}, \lean{explicitPresentation*} & proved, all seven types \\
exactness & \lean{ker_eq_of_liPresentationClassicallyFree} & proved (a conclusion) \\
generation & \lean{currentJetSurjectivity_of_genWordSpan} & proved, all types \\
B4 & \lean{restrictedDualEquivCurrentCarrier}, \lean{geometricCarrierEmbedding_of_quadraticIdeal} & proved \\
U3 & field \lean{DeligneCarrierData.u3} & certified per type \\
B1 & \lean{casimirLayerBound}, \lean{casimirAuxiliaryBound} & proved, generic \\
B2--B3 & \lean{auxDim_conductor_add}, \lean{finrank_currentCarrier_le} & proved, generic \\
B5 & \lean{Mspace_le} & proved, generic \\
B6 & \lean{mFreeness} & proved from U5 \\
generic locus & \lean{genericComparisonG} & proved, generic \\
B7 at $\Delta=0$ & \lean{pairComparisonG_of_cert}, \lean{X_pairComparison} & proved, certificates \\
B8 & \lean{carrierGenerationG}, \lean{X_divisorComparison} & proved from B6 \\
assembly & \lean{sectorFusionBoundG}, \lean{X_sectorFusionBound_of_U5} & proved from U5 ($A_2$: see Thm.~\ref{l24:thm:seven}) \\
D3 & \lean{sectorCoeff_eq_of_mixed} & proved \\
U5, U6 & \lean{FusionModuleBoundG}, \lean{MixedFusionRecursion} & hypotheses here; Chapter~\ref{ch:L26}\\
$(\star)$ & \lean{DemazureBoundFor} (only ``Demazure sum $\le$ character'') & Chapter~\ref{ch:L25}\\
final theorems & \lean{X_classicallyFree_final} & Chapter~\ref{ch:L27}\\
\bottomrule
\end{tabular}
\end{center}
U4 (Chari--Fourier--Khandai) does not appear: the campaign found that the Lean chain needs only the
\emph{upper bound} U5, because B6 can be proved from a dimension count (Theorem~\ref{l24:thm:B6}).

%======================================================================
\section{The conditional capstone}\label{l24:sec:capstone}
%======================================================================

The capstone states, for any VOA carrying $\dim\ltwentyfourg$ weight-one currents, what must be supplied to
obtain classical freeness, and it does so in a form that cannot be satisfied by a trivial choice of
ideal. We first need the infinite jet map and the definition of classical freeness used throughout.

\begin{definition}[Degree-zero classes; the infinite Li jet map; presentation-level classical freeness]\label{l24:def:lpcf}
Let $W$ be a VOA. For a state $x$ of weight $N$, its \emph{degree-zero class} is the image of $x$ in
$\ltwentyfourgr_{0,N}W=F^0W_N/F^1W_N$. For a family $x:\sigma\to\ltwentyfourgr W$ and an ideal
$I\subseteq\C[x_i:i\in\sigma]$ that vanishes on the $x_i$, the \emph{infinite Li jet map} is the unique
differential-algebra map $\ltwentyfourJ(\C[x]/I)\to\ltwentyfourgr W$ sending $x_i^{(r)}$ to
$D^{(r)}x_i=D^rx_i/r!$, where $D$ is the derivation of $\ltwentyfourgr W$ induced by $L_{-1}$. The pair
$(x,I)$ is \emph{classically free} if this map is bijective. Finally, for a Li bidegree $(p,N)$ and a jet order $m$, the
\emph{bidegree jet map at order $m$} is the restriction of the finite-jet map (jets $x_i^{(r)}$ with $r\le m$ only) to the
classes of polynomials with $p$ derivatives in total and conformal weight $N$, composed with the projection to the layer
$\ltwentyfourgr_{p,N}W$. Order $m=p$ suffices for Li degree $p$, since such a monomial has no jet of order above $p$.
\end{definition}
\noindent File \leanpath{VOA/\allowbreak Filtered/\allowbreak ClassicalFreeness.lean}.
\begin{leancode}
noncomputable def liDegreeZeroClass (W : VertexOperatorAlgebra V) (N : ℕ)
    (x : W.grading (N : ℤ)) : liAssociatedGraded W :=
  DirectSum.of (fun d : ℕ × ℕ => liWeightLayer W d.1 d.2) (0, N)
    (((liWeightPiece W 1 N).comap (liWeightPiece W 0 N).subtype).mkQ
      ⟨x.1, by
        rw [liWeightPiece, Submodule.mem_inf]
        exact ⟨by rw [liFiltration_zero_eq_top]; exact Submodule.mem_top, x.2⟩⟩)
\end{leancode}
\noindent File \leanpath{VOA/\allowbreak Filtered/\allowbreak InfiniteClassicalFreeness.lean}.
\begin{leancode}
noncomputable def liInfinitePresentedJetMap (W : VertexOperatorAlgebra V)
    (x : σ → liAssociatedGraded W) (I : Ideal (MvPolynomial σ ℂ))
    (hI : I ≤ RingHom.ker (MvPolynomial.aeval x).toRingHom) :
    InfiniteJet.PresentedJet (R := ℂ) I →ₐ[ℂ] liAssociatedGraded W :=
  InfiniteJet.presentedJetToDifferentialAlgebra (liAssociatedGradedTranslation W)
    (liAssociatedGradedTranslation_leibniz W) (liAssociatedGradedTranslation_one W) x I hI

def LiPresentationClassicallyFree (W : VertexOperatorAlgebra V)
    (x : σ → liAssociatedGraded W) (I : Ideal (MvPolynomial σ ℂ))
    (hI : I ≤ RingHom.ker (MvPolynomial.aeval x).toRingHom) : Prop :=
  Function.Bijective (liInfinitePresentedJetMap W x I hI)
\end{leancode}
\noindent File \leanpath{VOA/\allowbreak Filtered/\allowbreak ClassicalFreeness.lean}.
\begin{leancode}
noncomputable def liPresentedJetBidegreeMap
    (W : VertexOperatorAlgebra V) (m : ℕ)
    (x : σ → liAssociatedGraded W) (I : Ideal (MvPolynomial σ ℂ))
    (hI : I ≤ RingHom.ker (MvPolynomial.aeval x).toRingHom)
    (baseWeight : σ → ℕ) (p N : ℕ) :
    FiniteJet.PresentedJetLiBidegreeImage baseWeight m p N I →ₗ[ℂ]
      liWeightLayer W p N :=
\end{leancode}

\begin{dictionary}
\lean{liAssociatedGraded W} is the external direct sum $\bigoplus_{(p,N)}\ltwentyfourgr_{p,N}W$ with
components \lean{liWeightLayer W p N}$=F^pW_N/F^{p+1}W_N$ (Chapter~\ref{ch:L11}). \lean{DirectSum.of … (0, N)}
inserts an element in the $(0,N)$ component; \lean{mkQ} is the quotient map $F^0W_N\to F^0W_N/F^1W_N$; the
proof term in the angle brackets only certifies that $x\in F^0W_N$ (every state is). So
\lean{liDegreeZeroClass W N x} is exactly the degree-zero class.
\lean{InfiniteJet.PresentedJet I} is $\C[x_i^{(r)}:i\in\sigma,r\ge0]/\ltwentyfourJ I$ (Definition~\ref{l20:def:jet});
\lean{presentedJetToDifferentialAlgebra} is the map given by its universal property for the derivation
\lean{liAssociatedGradedTranslation W} (which satisfies Leibniz and kills $1$, the two proof arguments).
The hypothesis \lean{hI} is ``the relations vanish at order zero''. \lean{LiPresentationClassicallyFree} is
bijectivity, as in the textbook. \lean{PresentedJetLiBidegreeImage baseWeight m p N I} is the image in the
order-$m$ presented jet algebra of the polynomials of Li degree $p$ (number of derivatives) and conformal
weight $N$; \lean{liPresentedJetBidegreeMap} is the jet map restricted to that image and projected to the
layer $(p,N)$ (Chapter~\ref{ch:L11}, Theorem~\ref{l11:thm:freeness}).

\textbf{Same as the literature?} The docstring stresses that the source is the \emph{presented} arc
algebra of $\C[x]/I$, not an independently defined arc scheme. For the identification with
$\ltwentyfourJ R_V$ one needs $I=\ker$ exactly, which is what Theorem~\ref{l24:thm:kereq} supplies. With
$I=\ker$ and $x_i$ the classes of strong generators, \lean{LiPresentationClassicallyFree} is classical
freeness in the sense of Li, Arakawa--Moreau and van Ekeren--Heluani.
\end{dictionary}

\begin{physics}
The degree-zero class of a current $J^a_{(-1)}\ltwentyfourvac$ is the classical field $J^a$ on the Higgs
branch; $x^{(r)}$ is the classical shadow of $\partial^rJ/r!$. Classical freeness says that the composite
Schur operators are exactly the polynomials in the classical fields and their derivatives modulo the
derivatives of the classical (Higgs branch chiral ring) relations.
\end{physics}

\begin{status}
Definitions. They are reused unchanged from the Arakawa--Dai campaign (Chapter~\ref{ch:L20}).
\end{status}

\begin{definition}[The seven higher-rank cases]\label{l24:def:case}
The type \lean{HigherRankDeligneCase} has the seven elements $A_2,G_2,D_4,F_4,E_6,E_7,E_8$, with the current
dimension $\dim\ltwentyfourg$, the dual Coxeter number $h^\vee$ and the level $k$ of the table in
Section~\ref{l24:sec:math}; and $k=-(h^\vee+6)/6$ in every case.
\end{definition}
\noindent File \leanpath{VOA/\allowbreak Examples/\allowbreak DeligneHigherRankConditionalCapstone.lean}.
\begin{leancode}
inductive HigherRankDeligneCase
  | A2
  | G2
  | D4
  | F4
  | E6
  | E7
  | E8
  deriving DecidableEq, Fintype

def currentDimension : HigherRankDeligneCase → ℕ
  | A2 => 8
  | G2 => 14
  | D4 => 28
  | F4 => 52
  | E6 => 78
  | E7 => 133
  | E8 => 248

def level : HigherRankDeligneCase → ℚ
  | A2 => -3 / 2
  | G2 => -5 / 3
  | D4 => -2
  | F4 => -5 / 2
  | E6 => -3
  | E7 => -4
  | E8 => -6

theorem level_eq_neg_dualCoxeter_add_six_div_six
    (case : HigherRankDeligneCase) :
    level case = -((dualCoxeterNumber case : ℚ) + 6) / 6 := by
\end{leancode}

\begin{dictionary}
\lean{dualCoxeterNumber} gives $3,4,6,9,12,18,30$ (same file). $-(h^\vee+6)/6=-1-h^\vee/6$. The level
is used only as a label in the capstone; the actual VOAs of Section~\ref{l24:sec:presentation} are built at
the literal levels $-3/2,\dots,-6$ in $\C$, and their non-criticality $k+h^\vee\neq0$ is proved there.
\end{dictionary}

\begin{physics}
$k+h^\vee=5h^\vee/6-1>0$: these levels are negative but not critical, so the Sugawara construction exists
and the central charge $c=k\dim\ltwentyfourg/(k+h^\vee)$ is negative ($-8$ for $A_2$, $-14$ for $D_4$,
$-62$ for $E_8$, from $-6\cdot248/24$), as is typical of chiral algebras of 4d theories ($c_{2d}=-12c_{4d}$).
\end{physics}

\begin{status}
Definitions; the level formula is proved by \lean{norm_num} on each case.
\end{status}

\begin{definition}[The four packages]\label{l24:def:packages}
Fix a case and a VOA $W$.
\begin{enumerate}
\item A \emph{current $C_2$ presentation} $P$ consists of weight-one states $\mathrm{gen}_i$,
  $i<\dim\ltwentyfourg$, and an ideal $I\subseteq\C[X_0,\dots,X_{n-1}]$ contained in the kernel of
  $X_i\mapsto[\mathrm{gen}_i]\in\ltwentyfourgr_0W$.
\item \emph{Current-jet surjectivity}: for every $(p,N)$ the bidegree jet map at order $p$ is surjective,
  and the infinite jet map is surjective.
\item A \emph{geometric carrier embedding} into a family of spaces $\mathrm{Carrier}_{p,N}$: injective linear
  maps from the dual of each presented-jet bidegree piece into $\mathrm{Carrier}_{p,N}$.
\item The \emph{ordinary carrier character} identity: $\sum_{p\le N}\dim\mathrm{Carrier}_{p,N}=\dim W_N$ for all $N$.
\end{enumerate}
The case is \emph{classically free} for $P$ if $(x,I)$ is classically free, $x_i=[\mathrm{gen}_i]$.
\end{definition}
\noindent File \leanpath{VOA/\allowbreak Examples/\allowbreak DeligneHigherRankConditionalCapstone.lean}.
\begin{leancode}
structure CurrentC2Presentation where
  gen : Fin case.currentDimension → V
  gen_weight_one : ∀ i, gen i ∈ W.grading (1 : ℤ)
  ideal : Ideal (MvPolynomial (Fin case.currentDimension) ℂ)
  relations : ideal ≤ RingHom.ker
      (MvPolynomial.aeval fun i =>
        liDegreeZeroClass W 1 ⟨gen i, gen_weight_one i⟩).toRingHom

structure CurrentJetSurjectivity (P : CurrentC2Presentation case W) : Prop where
  bidegree : ∀ p N, Function.Surjective
    (liPresentedJetBidegreeMap W p
      (fun i => liDegreeZeroClass W 1 ⟨P.gen i, P.gen_weight_one i⟩)
      P.ideal P.relations_vanish (fun _ => 1) p N)
  infinite : Function.Surjective
    (liInfinitePresentedJetMap W
      (fun i => liDegreeZeroClass W 1 ⟨P.gen i, P.gen_weight_one i⟩)
      P.ideal P.relations_vanish)

structure GeometricCarrierEmbedding
    (P : CurrentC2Presentation case W)
    (Carrier : ℕ → ℕ → Type uC)
    [∀ p N, AddCommGroup (Carrier p N)]
    [∀ p N, Module ℂ (Carrier p N)] where
  toCarrier : ∀ p N,
    Module.Dual ℂ
      (FiniteJet.PresentedJetLiBidegreeImage
        (fun _ : Fin case.currentDimension => 1) p p N P.ideal) →ₗ[ℂ]
      Carrier p N
  injective : ∀ p N, Function.Injective (toCarrier p N)

def OrdinaryCarrierCharacter
    (Carrier : ℕ → ℕ → Type uC)
    [∀ p N, AddCommGroup (Carrier p N)]
    [∀ p N, Module ℂ (Carrier p N)]
    [∀ p N, FiniteDimensional ℂ (Carrier p N)] : Prop :=
  ∀ N,
    (∑ p ∈ Finset.range (N + 1), Module.finrank ℂ (Carrier p N)) =
      Module.finrank ℂ (W.grading N)

def CaseClassicallyFree (P : CurrentC2Presentation case W) : Prop :=
  LiPresentationClassicallyFree W
    (fun i => liDegreeZeroClass W 1 ⟨P.gen i, P.gen_weight_one i⟩)
    P.ideal P.relations_vanish
\end{leancode}

\begin{dictionary}
\lean{Fin case.currentDimension} is $\{0,\dots,\dim\ltwentyfourg-1\}$. \lean{P.relations_vanish} is a
one-line renaming of the field \lean{relations}. In \lean{bidegree}, the jet order is $p$ (the first
argument of \lean{liPresentedJetBidegreeMap W p …}) and the base weights are all $1$ (\lean{fun _ => 1}):
the currents have conformal weight one. Order $p$ suffices for Li degree $p$, since a monomial of Li degree $p$
contains at most $p$ derivatives. \lean{Module.Dual ℂ M} is $M^*=\operatorname{Hom}(M,\C)$. The sum in
\lean{OrdinaryCarrierCharacter} runs over $p\le N$ (\lean{Finset.range (N + 1)}), and \lean{W.grading N} is
$W_N$; the carriers are assumed finite-dimensional. No Li-bigraded identity is assumed, only the
ordinary one-variable character.

\textbf{Hypotheses and their meaning.} The \lean{relations} field is \emph{containment only}. The docstring
explains why this is safe: a theorem about a strictly smaller ideal could not pass for classical
freeness, because classical freeness forces equality (Theorem~\ref{l24:thm:kereq}). There is no
hypothesis that $W$ is an affine VOA or that the \lean{gen} are currents in any Lie-theoretic sense; the
capstone is a statement about any VOA with that many weight-one states. The specific VOAs come in
Section~\ref{l24:sec:presentation}.
\end{dictionary}

\begin{physics}
Package 1 lists classical null relations among the currents (the Higgs branch chiral-ring relations).
Package 2 says every operator is a normal-ordered polynomial in currents and their derivatives (strong
generation by the currents). Package 3 provides a ``particle'' model large enough to see every classical
operator. Package 4 says the particle model counts exactly the operators of the chiral algebra: the Schur
index equals the particle partition function.
\end{physics}

\begin{status}
Definitions (structures and predicates; nothing assumed globally).
\end{status}

\begin{theorem}[Weightwise squeeze]\label{l24:thm:squeeze}
Let $W$ be a VOA, $\mathrm{gen}:\mathrm{Fin}\,g\to W_1$, $I$ an ideal contained in the kernel, and $\mathrm{Carrier}$
finite-dimensional. If every bidegree jet map and the infinite jet map are surjective, if there are
injective maps $(\text{bidegree piece})^*\hookrightarrow\mathrm{Carrier}_{p,N}$, and if
$\sum_{p\le N}\dim\mathrm{Carrier}_{p,N}=\dim W_N$ for every $N$, then $(x,I)$ is classically free.
\end{theorem}
\noindent File \leanpath{VOA/\allowbreak Examples/\allowbreak DeligneWeightwiseHilbertCapstone.lean}.
\begin{leancode}
theorem classicallyFree_of_carrier_weightCharacter
    (W : VertexOperatorAlgebra V)
    (gen : Fin g → V)
    (hgen : ∀ i, gen i ∈ W.grading (1 : ℤ))
    (I : Ideal (MvPolynomial (Fin g) ℂ))
    (hI : I ≤ RingHom.ker
      (MvPolynomial.aeval fun i =>
        liDegreeZeroClass W 1 ⟨gen i, hgen i⟩).toRingHom)
    (Carrier : ℕ → ℕ → Type uC)
    [∀ p N, AddCommGroup (Carrier p N)]
    [∀ p N, Module ℂ (Carrier p N)]
    [∀ p N, FiniteDimensional ℂ (Carrier p N)]
    (hsurjBidegree : ∀ p N, Function.Surjective
      (liPresentedJetBidegreeMap W p
        (fun i => liDegreeZeroClass W 1 ⟨gen i, hgen i⟩)
        I hI (fun _ => 1) p N))
    (hsurjInfinite : Function.Surjective
      (liInfinitePresentedJetMap W
        (fun i => liDegreeZeroClass W 1 ⟨gen i, hgen i⟩) I hI))
    (sourceDualToCarrier : ∀ p N,
      Module.Dual ℂ
        (FiniteJet.PresentedJetLiBidegreeImage (fun _ : Fin g => 1) p p N I) →ₗ[ℂ]
          Carrier p N)
    (hcarrier : ∀ p N, Function.Injective (sourceDualToCarrier p N))
    (hcharacter : ∀ N,
      (∑ p ∈ Finset.range (N + 1), Module.finrank ℂ (Carrier p N)) =
        Module.finrank ℂ (W.grading N)) :
    LiPresentationClassicallyFree W
      (fun i => liDegreeZeroClass W 1 ⟨gen i, hgen i⟩) I hI := by
\end{leancode}

\begin{dictionary}
The hypotheses are the four packages, unbundled. The conclusion is bijectivity of the infinite jet map.
\end{dictionary}

\textbf{Proof on paper.} Fix $N$ and write, for $p\le N$, $c_p=\dim\ltwentyfourgr_{p,N}W$, $a_p$ the
dimension of the jet bidegree piece and $b_p=\dim\mathrm{Carrier}_{p,N}$. Surjectivity gives $c_p\le a_p$;
the injection of the dual gives $a_p=\dim(\text{piece})^*\le b_p$; and
$\sum_pc_p=\dim W_N=\sum_pb_p$ (the first equality is Li's telescoping identity, Chapter~\ref{ch:L11},
Theorem~\ref{l11:thm:litele}). By the finite squeeze (Theorem~\ref{l11:thm:squeeze}) $a_p=c_p$ for every
$p$. Bidegrees with $p>N$ are zero on both sides. So in every bidegree the surjective jet map is between
spaces of equal finite dimension, hence bijective; since the jet ideal of a homogeneous ideal is
bihomogeneous, the infinite jet map is injective, and it is surjective by hypothesis. \qed

\begin{physics}
Counting operators in three ways (quantum, classical jets, particle model) at fixed weight $N$, but only
the total over Li degree for the particle model: the squeeze recovers the Li-graded counts for free. This is
why a one-variable character identity, which is what modular/MLDE methods give, suffices.
\end{physics}

\begin{status}
Proved; \texttt{\#print axioms}: \lean{[propext, Classical.choice, Quot.sound]} (campaign record).
\end{status}

\begin{theorem}[Exact presentation from classical freeness]\label{l24:thm:kereq}
Let $I\subseteq\ker(\operatorname{aeval}x)$. If $(x,I)$ is classically free, then $\ker(\operatorname{aeval}x)=I$.
\end{theorem}
\noindent File \leanpath{VOA/\allowbreak DeligneAudit/\allowbreak C2PresentationOfFreeness.lean}.
\begin{leancode}
noncomputable def presentedJetOrderZero (I : Ideal (MvPolynomial σ ℂ)) :
    MvPolynomial σ ℂ →ₐ[ℂ] PresentedJet (R := ℂ) I :=
  (presentedJetMk (R := ℂ) I).comp (MvPolynomial.rename fun i : σ => (i, 0))

noncomputable def presentedJetRetraction (I : Ideal (MvPolynomial σ ℂ)) :
    PresentedJet (R := ℂ) I →ₐ[ℂ] (MvPolynomial σ ℂ ⧸ I) :=

theorem ker_eq_of_liPresentationClassicallyFree (W : VertexOperatorAlgebra V)
    (x : σ → liAssociatedGraded W) (I : Ideal (MvPolynomial σ ℂ))
    (hI : I ≤ RingHom.ker (MvPolynomial.aeval x).toRingHom)
    (hfree : LiPresentationClassicallyFree W x I hI) :
    RingHom.ker (MvPolynomial.aeval x).toRingHom = I := by
\end{leancode}

\begin{dictionary}
\lean{presentedJetOrderZero} sends $x_i\mapsto x_i^{(0)}$ (\lean{rename fun i => (i, 0)}) and then to the
quotient. \lean{presentedJetRetraction} is the map $\ltwentyfourJ(\C[x]/I)\to\C[x]/I$ given by the arc universal
property with the \emph{zero} derivation on $\C[x]/I$: it sends $x_i^{(0)}\mapsto\bar x_i$ and every
$x_i^{(r)}$, $r\ge1$, to $0$. The theorem uses only \lean{hfree.1}, injectivity, as the docstring says. The
statement is exactly the mathematical one: an ideal \emph{contained} in the kernel equals it.

\textbf{Could it be vacuous?} No: it is not a statement about any specific $I$; it says that the
classical-freeness conclusion is strong enough to pin down the ideal. Consequently a capstone with a
deliberately small ideal (say $I=0$) can only be proved if $0$ is the whole kernel.
\end{dictionary}

\textbf{Proof on paper.} Let $f\in\ker(\operatorname{aeval}x)$. The jet map composed with the order-zero
inclusion is $\operatorname{aeval}x$ (both are algebra maps agreeing on the $x_i$), so the order-zero image
$\iota(f)$ of $f$ maps to $0$; by injectivity $\iota(f)=0$ in $\ltwentyfourJ(\C[x]/I)$. Apply the retraction:
$r(\iota(f))=\bar f\in\C[x]/I$, and $r(0)=0$, so $f\in I$. The argument uses only injectivity of the infinite jet map,
not surjectivity; the Lean statement nevertheless takes the full (bijective) hypothesis
\lean{LiPresentationClassicallyFree}, and uses only \lean{hfree.1}. \qed

\begin{physics}
If the classical limit has no relations beyond the derivatives of those listed, then in particular the
listed chiral-ring relations are \emph{all} the chiral-ring relations: the Higgs branch (with its nilpotent
thickening) is cut out exactly by them. The retraction ``set all derivatives to zero'' is the restriction
from field configurations to constant (zero-mode) configurations.
\end{physics}

\begin{status}
Proved; \texttt{\#print axioms VOA.ker\_eq\_of\_liPresentationClassicallyFree}:
\lean{[propext, Classical.choice, Quot.sound]}. Re-checked by the author of this chapter with
\texttt{lake env lean -DmaxSynthPendingDepth=3} on a scratch file (Section~\ref{l24:sec:verification}).
\end{status}

\begin{theorem}[The conditional capstone]\label{l24:thm:capstone}
For any of the seven cases, any VOA $W$, any current $C_2$ presentation $P$ and any finite-dimensional carrier
family: current-jet surjectivity, a geometric carrier embedding and the ordinary carrier character identity
imply that the case is classically free for $P$ and that $P$'s ideal is the full kernel of the
current-coordinate map.
\end{theorem}
\noindent File \leanpath{VOA/\allowbreak Examples/\allowbreak DeligneHigherRankConditionalCapstone.lean}.
\begin{leancode}
theorem CurrentC2Presentation.exact_kernel_of_caseClassicallyFree
    {P : CurrentC2Presentation case W} (hfree : CaseClassicallyFree case W P) :
    RingHom.ker
        (MvPolynomial.aeval fun i =>
          liDegreeZeroClass W 1 ⟨P.gen i, P.gen_weight_one i⟩).toRingHom =
      P.ideal :=
  ker_eq_of_liPresentationClassicallyFree W _ P.ideal P.relations_vanish hfree

theorem caseClassicallyFree_of_explicit_inputs
    (P : CurrentC2Presentation case W)
    (Carrier : ℕ → ℕ → Type uC)
    [∀ p N, AddCommGroup (Carrier p N)]
    [∀ p N, Module ℂ (Carrier p N)]
    [∀ p N, FiniteDimensional ℂ (Carrier p N)]
    (hsurj : CurrentJetSurjectivity case W P)
    (hgeom : GeometricCarrierEmbedding case W P Carrier)
    (hchar : OrdinaryCarrierCharacter W Carrier) :
    CaseClassicallyFree case W P ∧
      RingHom.ker
          (MvPolynomial.aeval fun i =>
            liDegreeZeroClass W 1 ⟨P.gen i, P.gen_weight_one i⟩).toRingHom =
        P.ideal := by
\end{leancode}

\begin{dictionary}
The conclusion is a conjunction: classical freeness \emph{and} $\ker=$ \lean{P.ideal}. The hypotheses are
exactly packages 2, 3, 4; package 1 is $P$ itself. The proof is the composition of
Theorems~\ref{l24:thm:squeeze} and \ref{l24:thm:kereq}. As explained in the warning below, the strength of the
statement lies in package 4 for a carrier that is \emph{not} chosen ad hoc.
\end{dictionary}

\textbf{Proof on paper.} Theorem~\ref{l24:thm:squeeze} with the unbundled packages, then
Theorem~\ref{l24:thm:kereq}. \qed

\begin{physics}
``If the particle model is large enough (package 3) and counts the Schur index exactly (package 4), the
chiral algebra is classically free and its Higgs branch relations are the listed ones.''
\end{physics}

\begin{status}
Proved; \texttt{\#print axioms}: \lean{[propext, Classical.choice, Quot.sound]} (re-checked by the author,
Section~\ref{l24:sec:verification}). It is a reduction: no Deligne VOA is classically free in this chapter.
\end{status}

The module docstring of the capstone file and the docstring of \lean{caseClassicallyFree_of_explicit_inputs} are out of
date (they predate the general affine VOAs and the containment field); the statements are correct. Both documentation
defects are listed, once, in Section~\ref{l24:sec:defects}.

%======================================================================
\section[The explicit presentation]{The explicit presentation: singular vectors and coadjoint propagation}\label{l24:sec:presentation}
%======================================================================

The paper input U2 is: the symbols of the weight-two singular vectors lie in the kernel. The Lean proof is
uniform: (1) a vector killed by all positive current modes lies in the maximal ideal, so it is $0$ in
$L_k$; (2) the $C_2$-symbol of the explicit singular vector is the highest relation $P_0$ (up to a
scalar and a $C_2$ correction); (3) the kernel of the current-coordinate map is stable under the coadjoint
derivations $\operatorname{ad}_b$, because in $\ltwentyfourgr W$ the zero mode of a current acts as a
derivation; (4) a breadth-first list of single-term identities $\operatorname{ad}_bP_j=c\,P_k$, $j<k$,
reaches every relation.

\subsection{The $A_2$ relations}

\begin{definition}[The $A_2$ relation ideal]\label{l24:def:A2rel}
In the coordinates $X_0,\dots,X_7$ of $e_1,e_2,e_{12},f_1,f_2,f_{12},h_1,h_2\in\mathfrak{sl}_3$, put
$Z=3\sum_aX_ag^a$ (with $g^a$ the trace-dual basis) and let $P_a$ be one third of the entry of $Z^2$ at the
position of the $a$-th generator (and of the two diagonal differences for $a=6,7$). Then
$\ltwentyfourIns=\operatorname{span}(P_0,\dots,P_7)\cong V(\theta)$, and $I_2^{\mathrm{ns}}$ (the ideal)
is generated by these eight quadrics.
\end{definition}
\noindent File \leanpath{VOA/\allowbreak DeligneAudit/\allowbreak A2Relations.lean}.
\begin{leancode}
abbrev q (a b : Fin 8) : MvPolynomial (Fin 8) ℂ := X a * X b

def relationPoly : Fin 8 → MvPolynomial (Fin 8) ℂ :=
  ![ -- weight α₁:  e₁ (h₁ + 2h₂) + 3 e₁₂ f₂
     q 0 6 + 2 * q 0 7 + 3 * q 2 4,
     -- weight α₂:  -e₂ (2h₁ + h₂) + 3 e₁₂ f₁
     -(2 * q 1 6) - q 1 7 + 3 * q 2 3,
     -- weight θ (highest):  e₁₂ (h₁ - h₂) + 3 e₁ e₂
     q 2 6 - q 2 7 + 3 * q 0 1,
     -- weight -α₁:  f₁ (h₁ + 2h₂) + 3 e₂ f₁₂
     q 3 6 + 2 * q 3 7 + 3 * q 1 5,
     -- weight -α₂:  -f₂ (2h₁ + h₂) + 3 e₁ f₁₂
     -(2 * q 4 6) - q 4 7 + 3 * q 0 5,
     -- weight -θ:  f₁₂ (h₁ - h₂) + 3 f₁ f₂
     q 5 6 - q 5 7 + 3 * q 3 4,
     -- weight 0:  h₁² + 2 h₁h₂ + 3 e₁₂ f₁₂ - 3 e₂ f₂
     q 6 6 + 2 * q 6 7 + 3 * q 2 5 - 3 * q 1 4,
     -- weight 0:  -h₂² - 2 h₁h₂ - 3 e₁₂ f₁₂ + 3 e₁ f₁
     -q 7 7 - 2 * q 6 7 - 3 * q 2 5 + 3 * q 0 3 ]

def I2ns : Ideal (MvPolynomial (Fin 8) ℂ) := Ideal.span (Set.range relationPoly)
\end{leancode}

\begin{dictionary}
Namespace \lean{VOA.DeligneAudit.A2}. \lean{![…]} is a vector literal (entry $a$ is \lean{relationPoly a});
\lean{X a * X b} is the monomial $X_aX_b$. The comments give each polynomial in Lie notation; for
instance the highest one is $P_2=e_{12}(h_1-h_2)+3e_1e_2$. The Lean side proves homogeneity of degree $2$
(\lean{relationPoly_isHomogeneous}) and linear independence (\lean{relationPoly_linearIndependent}).
\textbf{Not Lean-checked}: that the span is $\operatorname{ad}$-stable and equals the adjoint summand of
$\operatorname{Sym}^2\mathfrak{sl}_3$; this was checked outside Lean, by exact rational computation in the generator
scripts. It is
not needed for soundness: the presentation theorem below proves each $P_a$ is in the kernel, and whether
the eight span all of $\ltwentyfourIns$ only matters for the \emph{later} character equality, where a wrong
list would make the carrier too large and the chain impossible to close. Note also that the Casimir is not
among the relations.
\end{dictionary}

\begin{physics}
These eight quadrics are the $C_2$-shadows of the eight null states at weight two of $L_{-3/2}(\mathfrak{sl}_3)$,
the chiral algebra of the $H_2$ Argyres--Douglas theory: Higgs branch chiral-ring relations transforming in
the adjoint of the $SU(3)$ flavour symmetry. \emph{Heuristic, not Lean:} the scheme they cut out in
$\mathfrak{sl}_3^*$ is non-reduced (the Casimir is nilpotent modulo the eight quadrics, not zero), and its reduced
variety is expected to be the minimal nilpotent orbit closure, the one-$SU(3)$-instanton moduli space. The library
proves neither statement, and it does not identify $(\ltwentyfourIns)$ with any ideal from the literature: for
$\mathfrak{sl}_3$ the ring $R_V$ is $\operatorname{Sym}\mathfrak{sl}_3$ modulo these eight quadrics and nothing more.
\end{physics}

\begin{status}
Definitions; homogeneity and independence proved, axioms standard.
\end{status}

\subsection{Singular vectors die in the simple quotient}

\begin{theorem}[Positive-mode singularity implies membership in the maximal ideal]\label{l24:thm:singp}
Let $\ltwentyfourg$ be given by Lie data $D$, $k+h^\vee\neq0$, and $t\in V^k(\ltwentyfourg)$ homogeneous of
weight $w>0$. If $J^a_{(i)}t=0$ for every current $a$ and every $i\ge1$, then $t$ lies in the maximal graded
ideal, hence vanishes in $L_k(\ltwentyfourg)$.
\end{theorem}
\noindent File \leanpath{VOA/\allowbreak DeligneAudit/\allowbreak D4LieSingularGeneric.lean} (a re-export, under the
\lean{D4Lie} namespace, of the type-generic lemmas of \leanpath{VOA/AffineGeneral/Singular.lean}).
\begin{leancode}
def SingP (t : Vacuum D k) : Prop := ∀ (a : ι) (i : ℤ), 1 ≤ i → J D k a i t = 0

theorem mem_maxIdeal_of_SingP {t : Vacuum D k} (ht : SingP D k t)
    {w : ℤ} (hw : 0 < w) (hmem : t ∈ gradingZA D k w) : t ∈ maxIdeal D k hk :=
  VOA.AffineGeneral.mem_maxIdeal_of_SingP D k hk ht hw hmem
\end{leancode}

\begin{dictionary}
Namespace \lean{VOA.DeligneAudit.D4Lie}, but the section is generic: \lean{D : LieData κ ι} for any finite
index type $\iota$ (Chapter~\ref{ch:L23}). \lean{Vacuum D k} is the vacuum module $V^k(\ltwentyfourg)$,
\lean{J D k a i} the mode $J^a_i=J^a_{(i)}$ of the current (for weight-one currents the two mode labels
coincide), \lean{gradingZA D k w} the weight-$w$ subspace and \lean{maxIdeal D k hk} the maximal graded
ideal, whose quotient is \lean{simpleVOA D k hk} $=L_k(\ltwentyfourg)$. The $A_2$ file
\leanpath{VOA/DeligneAudit/A2RelationsSingular.lean} has the same theorem for the bespoke
$\mathfrak{sl}_3$ construction (\lean{mem_maxIdeal_of_SingP3}). The zero vector is included (trivially).
\end{dictionary}

\textbf{Proof on paper.} If $t=0$ there is nothing to prove, so let $t\neq0$. Two observations.
(a) Zero modes preserve the hypothesis: if all $J^b_{(i)}$, $i\ge1$, kill $t$, they also kill $J^a_{(0)}t$, because
$[J^b_{(i)},J^a_{(0)}]=[b,a]_{(i)}$ (the central term $i\,k(b,a)\delta_{i,0}$ vanishes for $i\ge1$).
(b) \emph{Dong-type lemma}: for every state $u$ of weight $\wt u$ and every $N\ge\wt u$, $u_{(N)}t=0$. It suffices to take
$u$ a creation word $J^a_{(-n-1)}u'$, since such words span $V^k(\ltwentyfourg)$; induct on the length of the word. The
Borcherds iterate formula gives
\begin{equation*}
\bigl(J^a_{(-n-1)}u'\bigr)_{(N)}t=\sum_{i\ge0}(-1)^i\binom{-n-1}{i}\Bigl(J^a_{(-n-1-i)}\,u'_{(N+i)}t-(-1)^{-n-1}\,u'_{(-n-1+N-i)}\,J^a_{(i)}t\Bigr).
\end{equation*}
In the first term $N+i\ge\wt u'$, so $u'_{(N+i)}t=0$ by induction. In the second, $J^a_{(i)}t=0$ for $i\ge1$; for $i=0$,
$J^a_{(0)}t$ again satisfies the hypothesis by (a), and $-n-1+N\ge\wt u'$ because $\wt u=\wt u'+n+1$, so the induction
hypothesis applies. Since $u_{(N)}$ shifts the weight by $\wt u-N-1$, (b) says that every mode that lowers the weight
kills $t$. Weight-one states of $V^k(\ltwentyfourg)$ are quasi-primary ($L_1x_{(-1)}\ltwentyfourvac=0$). Hence $t$ is
\emph{strongly singular} of positive weight (Chapter~\ref{ch:L04}, Definition~\ref{l04:def:singular}: nonzero,
homogeneous of positive weight, and killed by every weight-lowering mode), so the ideal it generates lives in positive
weight, does not contain the vacuum, and lies in the maximal graded ideal (Definition~\ref{l04:def:maxideal}). \qed

\begin{physics}
A singular (null) vector generates a submodule of zero norm; in the physical (simple) theory it is set to
zero. The theorem says ``primary under the currents and of positive weight'' is enough to be null.
\end{physics}

\begin{status}
Proved; axioms standard. The $A_2$ analogue likewise.
\end{status}

\begin{theorem}[The $A_2$ singular vector]\label{l24:thm:A2sing}
At $k=-3/2$, the vector $v=2h_1(-1)e_{12}(-1)\ltwentyfourvac-2h_2(-1)e_{12}(-1)\ltwentyfourvac+6e_2(-1)e_1(-1)\ltwentyfourvac+3e_{12}(-2)\ltwentyfourvac$ of
$V^{-3/2}(\mathfrak{sl}_3)$ is killed by all positive current modes; hence its image in
$L_{-3/2}(\mathfrak{sl}_3)$ is zero.
\end{theorem}
\noindent File \leanpath{VOA/\allowbreak DeligneAudit/\allowbreak A2RelationsSingular.lean}.
\begin{leancode}
def singVec : Vacuum (-3 / 2) :=
  (2 : ℂ) • J (-3 / 2) 6 (-1) (J (-3 / 2) 2 (-1) (vacuumA (-3 / 2)))
    - (2 : ℂ) • J (-3 / 2) 7 (-1) (J (-3 / 2) 2 (-1) (vacuumA (-3 / 2)))
    + (6 : ℂ) • J (-3 / 2) 1 (-1) (J (-3 / 2) 0 (-1) (vacuumA (-3 / 2)))
    + (3 : ℂ) • J (-3 / 2) 2 (-2) (vacuumA (-3 / 2))

theorem singVec_singP3 : SingP3 (-3 / 2) singVec := by

theorem simpleProjection_singVec :
    simpleProjection (-3 / 2) level_add_three_ne_zero singVec = 0 := by
\end{leancode}

\begin{dictionary}
\lean{J (-3 / 2) a m} is the mode $m$ of the $a$-th generator on the bespoke $V^{-3/2}(\mathfrak{sl}_3)$
(the $\mathfrak{sl}_3$ files \leanpath{VOA/Examples/AffineSl3*.lean}); index $6,7$ are $h_1,h_2$, $2$ is $e_{12}$, $0,1$ are
$e_1,e_2$. \lean{SingP3 k t} is \lean{SingP} for this construction: $J^a_i t=0$ for all $a$ and $i\ge1$. \lean{simpleProjection} is the quotient map to $L_{-3/2}(\mathfrak{sl}_3)$. The vector is Perše's,
translated from the certificate basis (map recorded in \leanpath{A2Relations.lean}). Its $C_2$-symbol is
$2P_2$ minus the $C_2$-term $3e_{12}(-2)\ltwentyfourvac$. Whether $v\neq0$ is \emph{not} proved and not
needed (it would need PBW independence; the argument only uses membership in the ideal).
\end{dictionary}

\textbf{Proof on paper.} Modes $x(1)$ and $x(2)$ are computed from the affine commutator
$[x_m,y_n]=[x,y]_{m+n}+mk(x,y)\delta_{m+n,0}$: for each of the eight generators $x$ and each of the two modes $x(1)$,
$x(2)$, one expansion, sixteen in all, each giving zero exactly because $k=-3/2$. Modes $x(m)$, $m\ge3$, vanish because they lower the weight below zero. Then
Theorem~\ref{l24:thm:singp}. \qed

\begin{physics}
This is the weight-two null state of the $H_2$ chiral algebra; its existence at precisely $k=-3/2$ is the
2d reflection of the 4d Higgs branch relation $(\mu\otimes\mu)|_{\mathbf 8}=0$ for the $SU(3)$ moment map $\mu$.
\end{physics}

\begin{status}
Proved; \texttt{\#print axioms} of \lean{singVec_singP3}, \lean{simpleProjection_singVec}: standard three.
\end{status}

\subsection{Coadjoint propagation}

\begin{theorem}[Li zero modes are the coadjoint derivations]\label{l24:thm:adprop}
Let $W$ be any VOA with eight weight-one states $g_i$, and $T$ a table with $(g_b)_{(0)}g_i=\sum_cT_{bic}g_c$ in
weight one, as an equality of degree-zero classes (in $F^0W_1/F^1W_1$), not of states. Let $\operatorname{ad}_b$ be the derivation of $\C[X]$ with
$X_i\mapsto\sum_cT_{bic}X_c$. Then $\mathrm{cls}(\operatorname{ad}_bp)=(g_b)_{(0)}\cdot\mathrm{cls}(p)$ for every
polynomial $p$, where $\mathrm{cls}$ is the current-coordinate map into $\ltwentyfourgr W$. In particular
$\ker(\mathrm{cls})$ is $\operatorname{ad}$-stable.
\end{theorem}
\noindent File \leanpath{VOA/\allowbreak DeligneAudit/\allowbreak A2RelationsPresentation.lean}.
\begin{leancode}
def adPolyT (T : Fin 8 → Fin 8 → Fin 8 → ℂ) (b : Fin 8) :
    Derivation ℂ (MvPolynomial (Fin 8) ℂ) (MvPolynomial (Fin 8) ℂ) :=
  mkDerivation ℂ fun i => ∑ c, T b i c • X c

theorem clsG_adPolyT (T : Fin 8 → Fin 8 → Fin 8 → ℂ)
    (hT : ∀ b i, ∃ h : W'.smul (g b) 0 (g i) ∈ W'.grading ((1 : ℕ) : ℤ),
      liDegreeZeroClass W' 1 ⟨W'.smul (g b) 0 (g i), h⟩ = clsG W' g hg (∑ c, T b i c • X c))
    (b : Fin 8) (p : MvPolynomial (Fin 8) ℂ) :
    clsG W' g hg (adPolyT T b p) = zeroModeG W' g hg b (clsG W' g hg p) := by
\end{leancode}

\begin{dictionary}
\lean{clsG W' g hg} is \lean{aeval} of the degree-zero classes; \lean{zeroModeG … b} is the Li-graded mode
$0$ of the class of $g_b$. On $\ltwentyfourgr W$ the products $a_{(n)}b$ with $n\ge0$ lower the Li filtration and induce a
bracket; for $n=0$ and $a$ of weight one this is the Poisson bracket $\{g_b,-\}$ of the classical limit, and it is a
derivation of the commutative product (\lean{liAssociatedGradedMode_leibniz}). \lean{W'.smul (g b) 0 (g i)} is $(g_b)_{(0)}g_i$. The
generalization to any finite index set is \lean{clsG_adPolyT'} in
\leanpath{VOA/DeligneAudit/D4LiePresentationBase.lean}, used by all other types.
\end{dictionary}

\textbf{Proof on paper.} Both sides are derivations in $p$ (the right side by Leibniz for the Li zero mode),
and they agree on generators by $hT$. \qed

\begin{physics}
The flavour symmetry acts on the Higgs branch chiral ring by the coadjoint action; the space of relations
is a $G$-representation. This is the classical (Poisson) shadow of the current zero modes generating the
global symmetry.
\end{physics}

\begin{status}
Proved; axioms standard.
\end{status}

\begin{theorem}[The explicit presentations for all seven types]\label{l24:thm:sevenpres}
For each $X\in\{A_2,G_2,D_4,F_4,E_6,E_7,E_8\}$ and the simple affine VOA $W_X=L_k(\ltwentyfourg_X)$ at the Deligne
level, the ideal generated by the explicit relations of $\ltwentyfourIns$ is contained in the kernel of the
current-coordinate map, so $(\mathrm{currents},(\ltwentyfourIns))$ is a current $C_2$ presentation.
\end{theorem}
\noindent File \leanpath{VOA/\allowbreak DeligneAudit/\allowbreak A2RelationsPresentation.lean}.
\begin{leancode}
theorem I2ns_le_ker : I2ns ≤ RingHom.ker
    (MvPolynomial.aeval fun i =>
      liDegreeZeroClass W 1 ⟨gen i, gen_weight_one i⟩).toRingHom :=
  I2ns_le_ker_of_two hT_A2 clsG_relationPoly_two

def explicitPresentation : CurrentC2Presentation .A2 W where
  gen := gen
  gen_weight_one := gen_weight_one
  ideal := I2ns
  relations := I2ns_le_ker
\end{leancode}
\noindent File \leanpath{VOA/\allowbreak DeligneAudit/\allowbreak G2LiePresentation.lean}.
\begin{leancode}
theorem I2ns_le_ker : G2data.I ≤ RingHom.ker
    (MvPolynomial.aeval fun i =>
      liDegreeZeroClass g2W 1 ⟨genS g2CertData (-5/3) g2_level_ne i,
        genS_weight_one g2CertData (-5/3) g2_level_ne i⟩).toRingHom := by

def explicitPresentation : CurrentC2Presentation .G2 g2W where
  gen := genS g2CertData (-5/3) g2_level_ne
  gen_weight_one := genS_weight_one g2CertData (-5/3) g2_level_ne
  ideal := G2data.I
  relations := I2ns_le_ker
\end{leancode}
\noindent File \leanpath{VOA/\allowbreak DeligneAudit/\allowbreak D4LiePresentation.lean}.
\begin{leancode}
theorem I2ns_le_ker4 : DeligneCarrier.D4.D4data.I ≤ RingHom.ker
    (MvPolynomial.aeval fun i => liDegreeZeroClass W4 1 ⟨gen4 i, gen4_weight_one i⟩).toRingHom := by

def explicitPresentation4 : CurrentC2Presentation .D4 W4 where
  gen := gen4
  gen_weight_one := gen4_weight_one
  ideal := DeligneCarrier.D4.D4data.I
  relations := I2ns_le_ker4
\end{leancode}
\noindent File \leanpath{VOA/\allowbreak DeligneAudit/\allowbreak F4LiePresentation.lean}.
\begin{leancode}
theorem I2ns_le_ker : F4data.I ≤ RingHom.ker
    (MvPolynomial.aeval fun i =>
      liDegreeZeroClass f4W 1 ⟨genS f4CertData (-5/2) f4_level_ne i,
        genS_weight_one f4CertData (-5/2) f4_level_ne i⟩).toRingHom := by

def explicitPresentation : CurrentC2Presentation .F4 f4W where
  gen := genS f4CertData (-5/2) f4_level_ne
  gen_weight_one := genS_weight_one f4CertData (-5/2) f4_level_ne
  ideal := F4data.I
  relations := I2ns_le_ker
\end{leancode}
\noindent File \leanpath{VOA/\allowbreak DeligneAudit/\allowbreak E6Presentation.lean}.
\begin{leancode}
theorem I2ns_le_ker : E6data.I ≤ RingHom.ker
    (MvPolynomial.aeval fun i =>
      liDegreeZeroClass W6 1 ⟨genS e6LieData k6 hk6 i, genS_weight_one e6LieData k6 hk6 i⟩).toRingHom := by

def explicitPresentation : CurrentC2Presentation .E6 W6 where
  gen := genS e6LieData k6 hk6
  gen_weight_one := genS_weight_one e6LieData k6 hk6
  ideal := E6data.I
  relations := I2ns_le_ker
\end{leancode}
\noindent File \leanpath{VOA/\allowbreak DeligneAudit/\allowbreak E7LiePresentation.lean}.
\begin{leancode}
theorem I2ns_le_ker7 : DeligneCarrier.E7.E7data.I ≤ RingHom.ker
    (MvPolynomial.aeval fun i => liDegreeZeroClass W7 1 ⟨gen7 i, gen7_weight_one i⟩).toRingHom := by

def explicitPresentation7 : CurrentC2Presentation .E7 W7 where
  gen := gen7
  gen_weight_one := gen7_weight_one
  ideal := DeligneCarrier.E7.E7data.I
  relations := I2ns_le_ker7
\end{leancode}
\noindent File \leanpath{VOA/\allowbreak DeligneAudit/\allowbreak E8Presentation.lean}.
\begin{leancode}
theorem I2ns_le_ker8 : E8data.I ≤ RingHom.ker
    (MvPolynomial.aeval fun i => liDegreeZeroClass W8 1 ⟨gen8 i, gen8_weight_one i⟩).toRingHom := by

def explicitPresentation8 : CurrentC2Presentation .E8 W8 where
  gen := gen8
  gen_weight_one := gen8_weight_one
  ideal := E8data.I
  relations := I2ns_le_ker8
\end{leancode}

\begin{dictionary}
\textbf{Names.} The short name \lean{I2ns_le_ker} denotes four different theorems; always use the full
names: \lean{VOA.DeligneAudit.A2.I2ns_le_ker}, \lean{VOA.DeligneAudit.G2Lie.I2ns_le_ker},
\lean{VOA.DeligneAudit.F4Lie.I2ns_le_ker}, \lean{VOA.DeligneAudit.}\allowbreak\lean{DeligneCarrier.E6.I2ns_le_ker}, and
\lean{VOA.DeligneAudit.D4Lie.I2ns_le_ker4}, \lean{VOA.DeligneAudit.E7Lie.I2ns_le_ker7},
\lean{VOA.DeligneAudit.E8Lie.I2ns_le_ker8}. Likewise \lean{explicitPresentation} exists in the namespaces
\lean{A2}, \lean{G2Lie}, \lean{F4Lie}, \lean{DeligneCarrier.E6}.

\textbf{The VOAs.} $W$ (for $A_2$) is \lean{AffineSl3.simpleVOA (-3/2) _}, the bespoke $\mathfrak{sl}_3$
construction. \lean{g2W}, \lean{W4}, \lean{f4W}, \lean{W6}, \lean{W7}, \lean{W8} are
\lean{simpleVOA D k _} of Chapter~\ref{ch:L23} for the Lie data \lean{g2CertData}, \lean{d4CertData},
\lean{f4CertData}, \lean{e6LieData}, \lean{e7CertData}, \lean{e8CertData} at $k=-5/3,-2,-5/2,-3,-4,-6$
(\lean{k6} is $-3$). The currents \lean{genS …}, \lean{gen4}, … are $J^a_{(-1)}\ltwentyfourvac$ in the
``certificate basis'' of the Lie algebra, the basis in which the relation lists were computed.

\textbf{The ideals.} \lean{X.I} is \lean{quadraticIdeal X.rel} (Definition~\ref{l24:def:carrierdata}): the ideal
spanned by the listed quadrics (27, 105, 324, 650, 1539, 3875 of them). As for $A_2$, Lean proves each listed
quadric is in the kernel; that the list spans exactly $\ltwentyfourIns$ is checked outside Lean and is not needed for
soundness (see the dictionary of Definition~\ref{l24:def:A2rel}).

\textbf{Hypotheses.} None. These are unconditional theorems about specific VOAs.
\end{dictionary}

\textbf{Proof on paper.} For each type: (1) one weight-two singular vector (three for $D_4$, one per
summand) is killed by positive modes, checked by integer computations of $J^c(1)$ and $J^c(2)$ on
$x(-1)y(-1)\ltwentyfourvac$ (\lean{J_one_word2}, \lean{J_two_word2}) against the table; hence it is $0$ in
$L_k$ (Theorem~\ref{l24:thm:singp}). (2) Its $C_2$-symbol is the highest relation $P_0$; for the exceptional
types no correction is needed because the quadratic part has no $x(-2)$ term (for $A_2$ there is one, which
lies in $C_2$). (3) The kernel is $\operatorname{ad}$-stable (Theorem~\ref{l24:thm:adprop}). (4) A
breadth-first search from $P_0$ under the lowering operators $\operatorname{ad}F_i$ produced, for every
$k\ge1$, a parent $j<k$, a generator $b$ and nonzero scalars with $q\cdot\operatorname{ad}_bP_j=p\cdot P_k$;
each identity is a sparse polynomial identity checked by the kernel ($A_2$: 7 identities by hand;
$G_2$: 26; $D_4$: 102; $F_4$: 323; $E_6$: 649; $E_7$: 1538; $E_8$: 3874). Strong induction on $k$ gives every
$P_k$ in the kernel: if $P_j$ is in the kernel, so is $\operatorname{ad}_bP_j$, by the $\operatorname{ad}$-stability of (3), and
therefore so is $P_k=(q/p)\operatorname{ad}_bP_j$. \qed

\begin{physics}
One null state of weight two per irreducible summand of the relation module, plus the flavour symmetry,
generates all Higgs-branch relations at quadratic order. For $E_8$ this is the statement that the $3875$
quadratic relations of the one-$E_8$-instanton moduli space (besides the Casimir) follow from one null
vector of the Minahan--Nemeschansky chiral algebra.
\end{physics}

\begin{status}
Proved for all seven types; every \texttt{\#print axioms} in the presentation files is
\lean{[propext, Classical.choice, Quot.sound]} (campaign record: earlier umbrella build; later targeted builds; aggregate build green on 2026-10-28,
Section~\ref{l24:sec:verification}). Unconditional, as containments $\text{listed quadrics}\subseteq\ker$; equality of
the ideal with the full kernel is a theorem only for $A_2$ (Chapter~\ref{ch:L27}).
\end{status}

\begin{warning}[title={Which VOA? Certificate-basis Lie data}]
For $D_4$, $E_7$, $E_8$ the VOA of the presentation is \lean{simpleVOA d4CertData (-2) hk4} (resp.\
\lean{e7CertData}, \lean{e8CertData}): Lie data on the lattice Lie algebra (Chapter~\ref{ch:L23}) with the
certificate basis, the form Killing$/12$ (resp.\ $/36$, $/60$) and the Casimir transferred from the
lattice data. The library also has \lean{d4SimpleVOA} etc.\ built from \lean{Module.finBasis}. The two are
built from the same Lie algebra, form and $h^\vee$, but from different \lean{LieData} (different bases), and
the equality of the resulting VOA structures (basis independence of the affine construction) is
\emph{not} proved. All $D_4$, $E_7$, $E_8$ results of the campaign are about the
certificate-basis VOAs. Similarly three models of $\mathfrak e_6$ coexist (structure-constant table, lattice
$V_1$, the \lean{E6Lat} cert basis) and are not identified in Lean; the $E_6$ presentation uses the
structure-constant model \lean{e6LieData}. This is a scope remark, not an error: each theorem is about a
genuine $L_k(\ltwentyfourg)$.
\end{warning}

\begin{warning}[title={Dead end: the ``span plus $C_2$'' generation criterion}]
The older surjectivity lemmas \lean{liPresentedJetBidegreeMap_surjective_of_generators} and
\lean{liPresentedJetMap_hits_homogeneous_class_of_generators} (\leanpath{VOA/Filtered/LiGeneratorWords.lean},
Chapter~\ref{ch:L11}, Theorem~\ref{l11:thm:hits}) assume
$\top=\operatorname{span}(\mathrm{gen}\cup\{\ltwentyfourvac\})+C_2(V)$. For every Deligne VOA this is
\emph{false}: in weight two the right side is $C_2(V)_2$, but $(R_V)_2=\operatorname{Sym}^2\ltwentyfourg/\ltwentyfourIns\supseteq
V(2\theta)\oplus\C Q\neq0$; for any finite family it fails because $R_V$ is infinite-dimensional: these VOAs are
\emph{quasi-lisse} (the associated variety $\operatorname{Specm}R_V$ has finitely many symplectic leaves; here it is
$\overline{\mathbb O_{\min}}$) but not \emph{lisse} ($R_V$ finite-dimensional). Generation must come from ordered current
words (next section).
\end{warning}

%======================================================================
\section{Current-jet surjectivity}\label{l24:sec:jets}
%======================================================================

\begin{definition}[Generator words span Li's filtration]\label{l24:def:genword}
For states $a^1,\dots,a^g$ of weights $w_i$, a \emph{generator word} is
$a^{i_1}_{(-r_1-1)}\cdots a^{i_s}_{(-r_s-1)}\ltwentyfourvac$, of \emph{level} $r_1+\cdots+r_s$. Let $C_p$ be the span of
the generator words of level $\ge p$. The generators \emph{span Li's filtration} if $F^pV\subseteq C_p$ for
all $p$.
\end{definition}
\noindent File \leanpath{VOA/\allowbreak Filtered/\allowbreak GeneratorWordLiSurjectivity.lean}.
\begin{leancode}
abbrev genWordState (l : List (Fin g × ℕ)) : V :=
  wordAct (adjointWeakModule W) (genWord gen wt l) W.vacuum

def genWordPiece (p : ℕ) : Submodule ℂ V :=
  Submodule.span ℂ {u | ∃ l : List (Fin g × ℕ), p ≤ pairsLevel l ∧ genWordState W gen wt l = u}

def GeneratorWordsSpanLiFiltration : Prop :=
  ∀ p : ℕ, liFiltration W p ≤ genWordPiece W gen wt p
\end{leancode}

\begin{dictionary}
A list of pairs $(i,r)$ encodes a word, \lean{pairsLevel} is $\sum r$, and \lean{wordAct} applies the
letters as modes (Chapter~\ref{ch:L11}, Theorem~\ref{l11:thm:hits}). Since generator words of level $\ge p$
always lie in $F^p$ (\lean{genWordPiece_le_liFiltration}), the condition is equivalent to $F^p=C_p$. It is a
\emph{linear spanning} condition on all of $V$, which holds for strongly generated VOAs; it is not the
false ``$\operatorname{span}+C_2$'' condition of the warning above.
\end{dictionary}

\begin{physics}
Every operator of Li level $\ge p$ is a sum of normal-ordered products of the generators with at least $p$
derivatives in total. For affine VOAs this is the PBW basis of current modes.
\end{physics}

\begin{status}
Definitions (Chapter~\ref{ch:L11} extension: \leanpath{GeneratorCurrentLiSpanning.lean},
\leanpath{GeneratorWordLiSurjectivity.lean}).
\end{status}

\begin{theorem}[Current-jet surjectivity for every Deligne type]\label{l24:thm:jetsurj}
(i) For any case, VOA and presentation, if the currents span Li's filtration then
\lean{CurrentJetSurjectivity} holds. (ii) For every Lie data $D$ and non-critical $k$, the currents of
$L_k(\ltwentyfourg)$ span Li's filtration. Hence current-jet surjectivity holds for every explicit
presentation of Theorem~\ref{l24:thm:sevenpres}.
\end{theorem}
\noindent File \leanpath{VOA/\allowbreak DeligneAudit/\allowbreak A2CurrentJetSurjectivity.lean}.
\begin{leancode}
theorem currentJetSurjectivity_of_genWordSpan {V : Type*} [AddCommGroup V] [Module ℂ V]
    {case : HigherRankDeligneCase} {W : VertexOperatorAlgebra V}
    (P : CurrentC2Presentation case W)
    (h : GeneratorWordsSpanLiFiltration W P.gen (fun _ => 1)) :
    CurrentJetSurjectivity case W P :=
\end{leancode}
\noindent File \leanpath{VOA/\allowbreak AffineGeneral/\allowbreak DeligneInstance.lean}.
\begin{leancode}
theorem genWordsSpan :
    GeneratorWordsSpanLiFiltration (simpleVOA D k hk) (presentation case D k hk).gen (fun _ => 1) :=
\end{leancode}
\noindent File \leanpath{VOA/\allowbreak Filtered/\allowbreak GeneratorCurrentLiSpanning.lean}.
\begin{leancode}
def GeneratorModeTracked : Prop :=
  ∀ (i : Fin g) (m : ℤ) (q : ℕ) (v : V), v ∈ genWordPiece W gen wt q →
    W.smul (gen i) m v ∈ genWordPiece W gen wt (currentModeTarget q m)

theorem generatorWordsSpanLiFiltration_of_tracked (htrack : GeneratorModeTracked W gen wt)
    (hspan : Submodule.span ℂ (Set.range (genWordState W gen wt)) = ⊤) :
    GeneratorWordsSpanLiFiltration W gen wt := by

theorem generatorModeTracked_of_current_comm (hJ : ∀ i m v, W.smul (gen i) m v = J i m v)
    (hcomm : ∀ i j m p, J i m * J j p - J j p * J i m =
      (∑ c, f i j c • J c (m + p)) +
        (if m + p = 0 then cen i j m • (1 : Module.End ℂ V) else 0))
    (hvac : ∀ i m, 0 ≤ m → J i m W.vacuum = 0) :
    GeneratorModeTracked W gen wt := by

theorem generatorWordsSpanLiFiltration_quotient (h : GeneratorWordsSpanLiFiltration W gen wt) :
    GeneratorWordsSpanLiFiltration (QuotVOA.quotientVOA W I hI hgr)
      (fun i => I.mkQ (gen i)) wt := by
\end{leancode}

\begin{dictionary}
\lean{currentModeTarget q m} is $\max(q-m-1,0)$ (\leanpath{Filtered/AffineCurrentLi.lean}): the tracked
estimate says $a^i_{(m)}C_q\subseteq C_{\max(q-m-1,0)}$. \lean{hcomm} is the affine commutator with structure
constants $f$ and a central term \lean{cen}; \lean{hvac} is $J^i_m\ltwentyfourvac=0$ for $m\ge0$. The quotient
theorem passes spanning to any graded quotient, in particular to $L_k=V^k/\text{max ideal}$. In (ii),
\lean{presentation case D k hk} is the full-kernel presentation of Chapter~\ref{ch:L23}'s $L_k(\ltwentyfourg)$;
its \lean{gen} are the same currents as those of the explicit presentations, which is why each
\lean{explicitCurrentJetSurjectivity*} is a one-line application. For $A_2$ the bespoke version is
\lean{VOA.DeligneAudit.A2.currentJetSurjectivity}.
\end{dictionary}

\textbf{Proof on paper.} (i) The jet monomial $\prod x_{i_k}^{(r_k)}$ maps to the class of the corresponding
generator word (Chapter~\ref{ch:L11}); if $F^p=C_p$, every class in $\ltwentyfourgr_{p,N}$ is the class of a
combination of words of level exactly $p$, hence in the image. (ii) The affine commutator moves a mode
$J^i_m$ through a word: commuting past $J^j_{(-r-1)}$ produces $J^c_{m-r-1}$ (level drops by at most $m+1$)
and a central scalar; non-negative modes kill the vacuum. This gives the tracked estimate; with PBW
spanning of $V^k$ by current words, the Borcherds iterate formula identifies $F^p$ with $C_p$, and the
property descends to the simple quotient. \qed

\begin{physics}
The classical shadow of every operator is a differential polynomial in the currents: currents strongly
generate, and Li's filtration counts derivatives.
\end{physics}

\begin{status}
Proved, unconditional, all types; axioms standard (\lean{A2.currentJetSurjectivity},
\lean{currentJetSurjectivity_of_genWordSpan}, the three generic spanning theorems).
\end{status}

%======================================================================
\section{The geometric carrier}\label{l24:sec:carrier}
%======================================================================

\subsection{Quadratic ideals and the current carrier}

\begin{definition}[Quadratic relations and ideals]\label{l24:def:quad}
For a coefficient matrix $q:\sigma\times\sigma\to K$ put $\mathrm{quad}(q)=\sum_{x,y}q(x,y)X_xX_y$; for a family
$Q:\rho\to(\sigma\times\sigma\to K)$, $\mathrm{quadIdeal}(Q)$ is the ideal generated by the $\mathrm{quad}(Q_r)$.
\end{definition}
\noindent File \leanpath{VOA/\allowbreak Commutative/\allowbreak QuadraticJetCoefficientPairing.lean}.
\begin{leancode}
def quadraticRelation (q : σ × σ → K) : MvPolynomial σ K :=
  ∑ xy, q xy • (MvPolynomial.X xy.1 * MvPolynomial.X xy.2)

def quadraticIdeal (Q : ρ → σ × σ → K) : Ideal (MvPolynomial σ K) :=
  Ideal.span (Set.range (fun r => quadraticRelation (Q r)))
\end{leancode}

\begin{dictionary}
Exactly as stated. The matrix need not be symmetric; only $q(x,y)+q(y,x)$ matters.
\end{dictionary}

\begin{physics}
A quadratic Higgs-branch relation, written in a basis of the flavour currents: a matrix $q$ gives the
chiral-ring relation $\sum_{x,y}q(x,y)\,J^xJ^y=0$ among the moment-map components, and a family of such matrices is a
list of these relations. The ideal they generate is the ideal of the (possibly non-reduced) Higgs branch.
\end{physics}

\begin{status}
Definitions.
\end{status}

\begin{definition}[The current carrier $W_m$ and the auxiliary carrier]\label{l24:def:carrier}
Fix a basis $\sigma$ of $\ltwentyfourg^*$ and quadrics $Q$. A \emph{tensor polynomial} with $m$ current and $d$
auxiliary particles is a function $t\mapsto p_t$ from labellings $t:\{1..m\}\to\sigma$ to polynomials in
$z_1,\dots,z_{m+d}$. It is \emph{symmetric} if simultaneous permutation of roots and labels fixes it. For a
pair $i<j$ of current slots, a relation $q$ and spectator labels $t$, the \emph{collision readback} is
$\sum_{x,y}q(x,y)\,p_{t[i\mapsto x,j\mapsto y]}|_{z_j=z_i}$. The \emph{current carrier}
$\mathrm{currentCarrier}(Q,m,d)$ consists of the symmetric tensor polynomials all of whose readbacks vanish.
The \emph{auxiliary carrier} additionally requires, for every pair involving an auxiliary particle, that the
polynomial vanish on that collision. $\mathrm{degreePiece}(C,r)$ is the subspace of $C$ of polynomials
homogeneous of degree $r$. $W_m:=\mathrm{currentCarrier}(I_2,m,0)$.
\end{definition}
\noindent File \leanpath{VOA/\allowbreak Commutative/\allowbreak QuadraticCurrentFunctionalCarrier.lean}.
\begin{leancode}
abbrev TensorPolynomial (m d : ℕ) :=
  (Fin m → σ) → MvPolynomial (Particle m d) K

def Symmetric {m d : ℕ} (p : TensorPolynomial (K := K) (σ := σ) m d) : Prop :=
  ∀ (e : Equiv.Perm (Particle m d)) (g : Equiv.Perm (Fin m)),
    ExtendsCurrentPermutation e g → ∀ t,
      MvPolynomial.rename e (p t) = p (t ∘ g.symm)

def currentReadback {m d : ℕ} (q : σ × σ → K)
    (ij : IncreasingPair (Fin m)) (t : Fin m → σ)
    (p : TensorPolynomial (K := K) (σ := σ) m d) :=
  projectedDiagonal (ij.1.1.castAdd d) (ij.1.2.castAdd d)
    ((Fin.castAdd_injective m d).ne (ne_of_gt ij.2)) q
    (fun xy => p (insertPair ij t xy))

def CurrentEquations {m d : ℕ} (Q : ρ → σ × σ → K)
    (p : TensorPolynomial (K := K) (σ := σ) m d) : Prop :=
  ∀ r ij t, currentReadback (Q r) ij t p = 0

def currentCarrier (Q : ρ → σ × σ → K) (m d : ℕ) :
    Submodule K (TensorPolynomial (K := K) (σ := σ) m d) where
  carrier := {p | Symmetric p ∧ CurrentEquations Q p}

def auxiliaryCarrier (Q : ρ → σ × σ → K) (m d : ℕ) :
    Submodule K (TensorPolynomial (K := K) (σ := σ) m d) where
  carrier := {p | p ∈ currentCarrier Q m d ∧
    ∀ t (ij : Pair m d), m ≤ ij.1.2.val →
      collisionMap ij.1.2 ij.1.1 (ne_of_lt ij.2) (p t) = 0}
\end{leancode}
\noindent File \leanpath{VOA/\allowbreak Commutative/\allowbreak DeligneAuxiliaryCollision.lean}.
\begin{leancode}
abbrev Particle (m d : ℕ) := Fin (m + d)

def conductor (m d : ℕ) : MvPolynomial (Particle m d) K :=
  ∏ ij : Pair m d, factor (K := K) ij ^ order ij

def conductorDegree (m d : ℕ) : ℕ := ∑ ij : Pair m d, order ij
\end{leancode}
\noindent File \leanpath{VOA/\allowbreak Commutative/\allowbreak QuadraticCurrentGradedCarrier.lean}.
\begin{leancode}
def degreePiece {m d : ℕ}
    (C : Submodule K (TensorPolynomial (K := K) (σ := σ) m d)) (r : ℕ) :
    Submodule K C where
  carrier := {p | ∀ t, (p.1 t).IsWeightedHomogeneous (fun _ => 1) r}
\end{leancode}

\begin{dictionary}
\lean{Particle m d} $=\{0,\dots,m+d-1\}$: currents first, then auxiliary particles. \lean{ExtendsCurrentPermutation e g}
says the root permutation $e$ restricts to $g$ on current slots (auxiliary roots may be permuted freely).
\lean{insertPair ij t (x,y)} replaces the labels in slots $i,j$ by $x,y$. \lean{projectedDiagonal a b _ c p}
is $\sum_tC(c_t)\,\mathrm{collisionMap}_{a\to b}(p_t)$, where \lean{collisionMap} substitutes $z_a:=z_b$
(Chapter~\ref{ch:L20}). The auxiliary condition is vanishing of $p_t$ on every collision $z_i=z_j$ with $j$
auxiliary (\lean{m ≤ ij.1.2.val}). The \lean{conductor} is $\prod_{i<j}(z_j-z_i)^{o(i,j)}$ with $o=0$ for two
currents, $1$ for a current and an auxiliary particle, $2$ for two auxiliary ones; its degree is
$md+d(d-1)$ (\lean{conductorDegree_eq}). The proof fields of the submodules are omitted.

\textbf{Same as the paper?} Yes: this is B4's $W_m$, literally as ``polynomial kernel'', with no fusion
module or representation theory in the definition (the docstring says so).
\end{dictionary}

\begin{physics}
A Feigin--Stoyanovsky ``functional model'': a state with $m$ current quanta is a symmetric wave function of
their positions (the $z_i$) with values in $\ltwentyfourg^{*\otimes m}$; the quadratic relations become
boundary conditions on two-particle collisions (an exclusion rule). Casimir quanta are extra particles with
a stronger exclusion (a double zero against each other).
\end{physics}

\begin{status}
Definitions; built in the umbrella target \lean{VOA.Deligne}.
\end{status}

\begin{theorem}[The restricted dual of a jet piece is a carrier piece (B4)]\label{l24:thm:B4}
For any field $K$, quadrics $Q$, and $r\le N$: the dual of the piece of the order-$N$ presented jet algebra of
$\mathrm{quadIdeal}(Q)$ with $m$ particles and $r$ derivatives is linearly isomorphic to
$\mathrm{degreePiece}(\mathrm{currentCarrier}(Q,m,0),r)$.
\end{theorem}
\noindent File \leanpath{VOA/\allowbreak Commutative/\allowbreak QuadraticJetRestrictedDual.lean}.
\begin{leancode}
def restrictedDualEquivCurrentCarrier (Q : ρ → σ × σ → K) (N m r : ℕ) (hr : r ≤ N) :
    (quotientPiece Q N m r →ₗ[K] K) ≃ₗ[K] degreePiece (currentCarrier Q m 0) r :=
  (pieceDualEquivCoefficients Q N m r).trans (fixedQuadraticReindexingEquiv Q N m r hr)
\end{leancode}

\begin{dictionary}
\lean{quotientPiece Q N m r} is the image in $J_N(K[\sigma]/\mathrm{quadIdeal}\,Q)$ of the jet polynomials with $m$
factors (``particles'') and total derivative order $r$. \lean{≃ₗ[K]} is a linear isomorphism. The isomorphism
reads a functional on jet monomials $x_{t_1}^{(a_1)}\cdots x_{t_m}^{(a_m)}$ as the coefficient of
$z_1^{a_1}\cdots z_m^{a_m}$ in the $t$-component (\lean{restrictedDualEquivCurrentCarrier_coeff}). No
characteristic or representation-theoretic hypothesis.
\end{dictionary}

\textbf{Proof on paper.} A functional on the jet algebra is a family of coefficients $F(e)$ on jet
monomials. It descends to the quotient iff it annihilates the ideal, i.e.\ iff for every relation, every
order $n$ and every spectator monomial, the Cauchy convolution $\sum_{a+b=n}q(x,y)F(\dots x^{(a)}y^{(b)}\dots)$
vanishes (spectators must be handled explicitly; annihilating the generators alone is not enough). Reindex
$F$ as a tensor polynomial: the coefficient of $z^{a}$ at labels $t$. The convolution of order $n$ in slots
$i,j$ is the $w^n$-coefficient of the readback at $z_i=z_j=w$; so the descent condition is exactly the
current equations; symmetry is commutativity of the jet monomials. \qed

\begin{physics}
The dual of the space of classical operators with $m$ currents and $r$ derivatives is the space of
$m$-particle wave functions of total degree $r$ obeying the collision rules: operators $\leftrightarrow$
correlation-function-like generating functions.
\end{physics}

\begin{status}
Proved; axioms standard.
\end{status}

\begin{theorem}[Geometric carrier embedding for every quadratic presentation]\label{l24:thm:geom}
If the ideal of a current $C_2$ presentation is $\mathrm{quadIdeal}(Q)$, then
$\mathrm{Carrier}_{p,N}:=\mathrm{degreePiece}(\mathrm{currentCarrier}(Q,N-p,0),p)$ receives an injective (for
$p\le N$ bijective) restricted-dual map in every bidegree; hence a \lean{GeometricCarrierEmbedding}.
\end{theorem}
\noindent File \leanpath{VOA/\allowbreak DeligneAudit/\allowbreak A2CarrierEmbedding.lean}.
\begin{leancode}
abbrev carrierPiece (Q : ρ → σ × σ → K) (p N : ℕ) : Submodule K (currentCarrier Q (N - p) 0) :=

abbrev Carrier (Q : ρ → σ × σ → K) (p N : ℕ) : Type _ := carrierPiece Q p N

def geometricCarrierEmbedding_of_quadraticIdeal {ρ : Type}
    (P : CurrentC2Presentation case W)
    (Q : ρ → Fin case.currentDimension × Fin case.currentDimension → ℂ)
    (hQ : P.ideal = quadraticIdeal Q) :
    GeometricCarrierEmbedding case W P (Carrier Q) where
  toCarrier := (carrierMaps Q P.ideal hQ).1
  injective := (carrierMaps Q P.ideal hQ).2
\end{leancode}

\begin{dictionary}
\lean{carrierPiece Q p N} is \lean{degreePiece (currentCarrier Q (N - p) 0) p} (the body is omitted in the
quote). With base weights one, Li bidegree $(p,N)$ means $N-p$ currents and $p$ derivatives
(\lean{jetLiBidegreePiece_eq_polynomialPiece}); for $N<p$ the piece is $0$. A companion
\lean{geometricCarrierEmbedding_of_span} accepts an ideal spanned by homogeneous quadrics, which it converts
to \lean{quadraticIdeal} form (\lean{span_homogeneousQuadratic_eq_quadraticIdeal}).
\end{dictionary}

\textbf{Proof on paper.} Theorem~\ref{l24:thm:B4} plus the bidegree dictionary. \qed

\begin{physics}
The particle model sees every classical operator, for free: it \emph{is} the dual of the classical operator
space.
\end{physics}

\begin{status}
Proved; axioms standard. Instantiated for all seven types
(\lean{A2.geometricCarrierEmbedding}, \lean{explicitGeometricCarrierEmbedding4}, \lean{E6.geometricCarrierEmbedding}, …).
\end{status}

\begin{warning}[title={The geometric package is weak on its own}]
For \emph{any} presentation one could take $\mathrm{Carrier}_{p,N}:=$ the dual of the jet piece and the
identity map; \lean{GeometricCarrierEmbedding} would hold trivially. Its content is only in combination with
\lean{OrdinaryCarrierCharacter} for the \emph{same} carrier. For the literal carrier above, the carrier is
(dual to) the jet algebra itself, so package 4 is the statement $\ltwentyfourHilb\ltwentyfourJ\ltwentyfourRt=\ltwentyfourch V$,
which is the whole content of classical freeness. All the hard mathematics (U3--U6, B1--B8, $(\star)$) is in
proving the inequality $\le$ of the next theorem; the geometric package is not where the difficulty lies.
\end{warning}

\begin{theorem}[The character package is one inequality]\label{l24:thm:upper}
Given current-jet surjectivity and a geometric carrier embedding, $\dim W_N\le\sum_{p\le N}\dim\mathrm{Carrier}_{p,N}$
holds automatically; hence \lean{OrdinaryCarrierCharacter} follows from the reverse inequality. For a
quadratic presentation, classical freeness and $\ker=\mathrm{quadIdeal}(Q)$ follow from current-jet surjectivity
and the named residual $\sum_p\dim\mathrm{Carrier}(Q)_{p,N}\le\dim W_N$.
\end{theorem}
\noindent File \leanpath{VOA/\allowbreak DeligneAudit/\allowbreak A2CarrierUpperBound.lean}.
\begin{leancode}
structure QuadraticCarrierUpperBound {ρ : Type} {n : ℕ} (Q : ρ → Fin n × Fin n → ℂ)
    (W : VertexOperatorAlgebra V) : Prop where
  le : ∀ N, (∑ p ∈ Finset.range (N + 1), Module.finrank ℂ (Carrier Q p N)) ≤
    Module.finrank ℂ (W.grading N)

theorem caseClassicallyFree_of_quadraticCarrierUpperBound {ρ : Type}
    (P : CurrentC2Presentation case W)
    (Q : ρ → Fin case.currentDimension × Fin case.currentDimension → ℂ)
    (hQ : P.ideal = quadraticIdeal Q)
    (hsurj : CurrentJetSurjectivity case W P)
    (hbound : QuadraticCarrierUpperBound Q W) :
    CaseClassicallyFree case W P ∧
      RingHom.ker
          (MvPolynomial.aeval fun i =>
            liDegreeZeroClass W 1 ⟨P.gen i, P.gen_weight_one i⟩).toRingHom =
        P.ideal :=
\end{leancode}

\begin{dictionary}
\lean{QuadraticCarrierUpperBound Q W} is a named \emph{hypothesis} structure (a \lean{Prop}), never an axiom; by
\lean{quadraticCarrierUpperBound_iff} it is equivalently ``$\sum_p\dim$(jet bidegree piece)$\le\dim W_N$'', i.e.\
$\ltwentyfourHilb\ltwentyfourJ(\C[X]/\mathrm{quadIdeal}\,Q)\le\ltwentyfourch W$ coefficientwise. Per type this gives one-input
theorems, e.g.\ \lean{D4Lie.explicit_classicallyFree_of_upperBound},
\lean{DeligneCarrier.E6.classicallyFree_of_upperBound}.
\end{dictionary}

\textbf{Proof on paper.} $\dim W_N=\sum_p\dim\ltwentyfourgr_{p,N}\le\sum_p\dim(\text{jet piece})\le\sum_p\dim\mathrm{Carrier}_{p,N}$
by surjectivity and injectivity; with the reverse inequality we get equality, and the capstone applies. \qed

\begin{physics}
One direction is free: the chiral algebra never has more operators of weight $N$ than its classical limit, because
every operator has a classical shadow (Li's map is onto); this is $\dim W_N\le\sum_p\dim\mathrm{Carrier}_{p,N}$. The
nontrivial direction is the reverse, that the classical limit has no more operators than the chiral algebra (no
classical operator fails to lift). Proving that upper bound on the classical count forces equality, and equality is
classical freeness. This is the same statement as the physics box of Theorem~\ref{l24:thm:carrierside}, seen from the
other end of the chain.
\end{physics}

\begin{status}
Proved; axioms standard. The upper bound itself is what the rest of the Deligne campaign
(this chapter's chain plus Chapters~\ref{ch:L25}--\ref{ch:L27}) supplies: unconditionally for $A_2$, conditionally on
named inputs for the other six types.
\end{status}

\begin{example}[$A_2$ in weights one and two]\label{l24:ex:A2weight2}
Take $\ltwentyfourg=\mathfrak{sl}_3$, $n=8$, the presentation ideal $(\ltwentyfourIns)$ with its eight quadrics, and
$\mathrm{Carrier}_{p,N}=\mathrm{degreePiece}(\mathrm{currentCarrier}(\mathcal R,N-p,0),p)$: $N-p$ current particles and
$p$ derivatives.
\begin{itemize}
\item $N=1$. For $p=0$: one particle, degree $0$: a function of one label with constant values, dimension $8$. For $p=1$:
  zero particles and degree one, which is $0$ (a polynomial in no variables of degree one vanishes). Total $8$, the
  number of currents $x_{(-1)}\ltwentyfourvac$.
\item $N=2$. For $p=0$: two particles, degree $0$, symmetric, with the collision equations of the eight relations; this
  is the dual of $\operatorname{Sym}^2\ltwentyfourg/\ltwentyfourIns$, of dimension $36-8=28=27+1$ (the summand $V(2\theta)$
  and the Casimir $Q$, which is \emph{not} killed). For $p=1$: one particle, degree one: the eight linear functions
  $c\,z_1$, one per label, dimension $8$ (the jets $x^{(1)}$, shadows of $x_{(-2)}\ltwentyfourvac$). For $p=2$: zero
  particles, $0$. Total $28+8=36$.
\end{itemize}
On the VOA side, $V^{-3/2}(\mathfrak{sl}_3)$ has $8+36=44$ states of weight two (the $x_{(-2)}\ltwentyfourvac$ and the
symmetric products $x_{(-1)}y_{(-1)}\ltwentyfourvac$). The singular vector of Theorem~\ref{l24:thm:A2sing} and its
$\mathfrak{sl}_3$-orbit, an eight-dimensional copy of the adjoint, die in $L_{-3/2}(\mathfrak{sl}_3)$, and the literature character
$\ltwentyfourch L_{-3/2}(\mathfrak{sl}_3)=1+8q+36q^2+\cdots$ says nothing else does in weight two: $44-8=36$. So in weights
one and two the carrier count equals $\dim W_N$, which is what the character package (and Theorem~\ref{l24:thm:upper})
asks in every weight. What Lean proves here, with no input, is the inequality $\dim(W_N)\le$ carrier count of the
theorem above, so $\dim L_{-3/2}(\mathfrak{sl}_3)_2\le36$; the reverse inequality, in all weights, is the open
content.
\end{example}

\subsection{The commutative layer, collectively}\label{l24:sec:commutative}

Twenty-one files of \leanpath{VOA/Commutative/} serve the carrier. They are pure commutative algebra over a
field, with no Lie theory. We list what each group proves.
\begin{itemize}
\item \emph{Coefficient duality} (\leanpath{QuadraticJetCoefficientPairing}, \leanpath{QuadraticCurrentCoefficientReindexing},
  \leanpath{QuadraticCurrentCollisionCoefficients}, \leanpath{QuadraticJetRestrictedDual}, \leanpath{QuadraticCurrentFunctionalCarrier}, the file of
  Definition~\ref{l24:def:carrier}): Theorem~\ref{l24:thm:B4}.
\item \emph{Auxiliary particles} (\leanpath{DeligneAuxiliaryCollision}, \leanpath{DeligneAuxiliaryDegree},
  \leanpath{QuadraticCurrentAuxiliaryCoefficients}, \leanpath{QuadraticCurrentAuxiliaryTensor},
  \leanpath{QuadraticCurrentAuxiliaryGrading}, \leanpath{QuadraticCurrentGradedCarrier},
  \leanpath{SymmetricPolynomialTensorCoefficients}): the conductor divides every auxiliary-carrier element and division is
  an isomorphism onto the current carrier with $d$ spectators shifted by $md+d(d-1)$ (\lean{gradedConductorEquiv},
  \lean{finrank_auxiliary_eq_current}); the current carrier with $d$ auxiliary roots is
  $\mathrm{currentCarrier}(Q,m,0)\otimes\C[v_1..v_d]^{S_d}$, degree by degree (\lean{currentAuxiliaryTensorEquiv},
  \lean{gradedCoefficientsEquiv}). These are B2--B3. They reuse \leanpath{RootedGordonConductorDivisibility} and
  \leanpath{P3ConductorDegree} from Chapter~\ref{ch:L20}.
\item \emph{Casimir filtration} (\leanpath{CasimirCostFiltration}, \leanpath{FiniteCostDegeneration},
  \leanpath{DeligneCasimirJetFiltration}, \leanpath{MvPolynomialMinimumWeight}): the filtration of a weight-$N$ quotient
  by powers of an ideal generated in weight $\ge2$ has length $\le N/2$, and its layers telescope.
\item \emph{Two points} (\leanpath{QuadraticTwoPointCarrier}): B7 at $m=2$ (historical; superseded by the certificates
  of Section~\ref{l24:sec:divisor}).
\item \emph{Globalization} (\leanpath{HeightOneGlobalization}, \leanpath{FreeLatticeGlobalization},
  \leanpath{DiscriminantLatticeGlobalization}): a noetherian UFD is the intersection of its localizations at height-one
  primes (\lean{iInf_heightOneLocalization_eq_bot}; Mathlib has this only for Dedekind domains), and the free-lattice
  version (B8's normality).
\item \emph{Support} (\leanpath{FiniteJetLiSupport}): jet bidegrees $(p,N)$ with $p>N$ are zero.
\end{itemize}
All are unconditional and \texttt{\#print axioms}-standard.

%======================================================================
\section{The carrier chain: B1 and the resummation D3}\label{l24:sec:chain}
%======================================================================

\begin{definition}[Carrier type data]\label{l24:def:carrierdata}
Carrier data $\mathscr C$ (Lean: \lean{D}) consist of $n=\dim\ltwentyfourg$, a label set $\rho$ (no finiteness is required here; sector data add
\lean{ρFin}), relations $\mathrm{rel}:\rho\to$ (matrices) spanning
$\ltwentyfourIns$, the Casimir matrix $\mathrm{cas}$, and a proof of U3: $X_a\cdot Q\in\mathrm{quadIdeal}(\mathrm{rel})$ for every
$a$. Derived: $I=\mathrm{quadIdeal}(\mathrm{rel})$, $Q=\mathrm{quad}(\mathrm{cas})$, the carrier family, and the full
matrix family $\mathrm{fullMatrix}:\mathrm{Option}\,\rho\to\cdots$ with \lean{none}$\mapsto\mathrm{cas}$ (i.e.\ $I_2$).
\end{definition}
\noindent File \leanpath{VOA/\allowbreak DeligneAudit/\allowbreak DeligneCarrierData.lean}.
\begin{leancode}
structure DeligneCarrierData where
  n : ℕ
  ρ : Type
  rel : ρ → Fin n × Fin n → ℂ
  cas : Fin n × Fin n → ℂ
  u3 : ∀ a : Fin n, X a * quadraticRelation cas ∈ quadraticIdeal rel

theorem Q_sq_mem : D.Q * D.Q ∈ D.I :=

def fullMatrix : Option D.ρ → Fin D.n × Fin D.n → ℂ
  | none => D.cas
  | some r => D.rel r
\end{leancode}

\begin{dictionary}
Namespace \lean{VOA.DeligneAudit.DeligneCarrier}. The only proof obligation is \lean{u3}. Nothing forces
\lean{cas} to be the Casimir or \lean{rel} to span $\ltwentyfourIns$; the instances use the certified tables
(Section~\ref{l24:sec:instances}), and the chain is sound for any data satisfying U3. The link to the VOA is
made through the ideal: the presentation's ideal is \lean{X.I} by \lean{rfl}. \lean{Q_sq_mem} ($Q^2\in I$)
follows from U3 since $Q$ lies in the ideal of the $X_a$.
\end{dictionary}

\begin{physics}
U3 says the Sugawara shadow $Q$ is annihilated by every current in the chiral ring modulo $\ltwentyfourIns$:
$T$ is nilpotent ($Q^2=0$) in $R_V$.
\end{physics}

\begin{status}
Definition; \lean{Q_sq_mem} proved (axioms standard). Instances: \lean{A2data}, \lean{D4data}, \lean{G2data},
\lean{F4data}, \lean{E6data}, \lean{E7data}, \lean{E8data}, each with U3 kernel-checked.
\end{status}

\begin{definition}[Carrier and sector dimensions; the named chain inputs]\label{l24:def:chaininputs}
For carrier data $\mathscr C$: $\mathrm{carrierDim}(N)=\sum_{p\le N}\dim\mathrm{Carrier}_{p,N}$;
$\mathrm{auxDim}(m,d,p)=\dim$ of degree-$p$ piece of the auxiliary carrier of $I_2$; $\mathrm{sectorDim}(N,m,d)=\mathrm{auxDim}(m,d,N-m-2d)$ if
$m+2d\le N$, else $0$; $\mathrm{currentDim}(m,d,r)$ the analogue for the current carrier. The \emph{sector fusion
bound} for a family $F_m$ of power series is $\mathrm{currentDim}(m,d,r)\le[q^r]F_m/((q)_m(q)_d)$; the \emph{mixed
fusion recursion} (U6) is $\mathrm{Dem}_n=M_{n,0}$ for the recursion $M_{0,b}=F_b$,
$M_{a+1,b}=M_{a,b+1}+q^{2a+b+1}M_{a,b}$. Here $\mathrm{Dem}_n$ stands for $\ltwentyfourHilb_qD(1,n\theta)$, the
$t$-graded dimension of the $n$-th level-one Demazure module, and the recursion is Ravinder's: it builds the Demazure
module of $n\theta$ out of fusion modules $\mathcal F_b$, one step $a\mapsto a+1$ at a time, the term $q^{2a+b+1}M_{a,b}$
recording the energy of the added piece (Chapter~\ref{ch:L26} discusses it).
\end{definition}
\noindent File \leanpath{VOA/\allowbreak DeligneAudit/\allowbreak DeligneCarrierChain.lean}.
\begin{leancode}
def auxDim (m d p : ℕ) : ℕ :=
  Module.finrank ℂ (degreePiece (auxiliaryCarrier D.fullMatrix m d) p)

def currentDim (m d r : ℕ) : ℕ :=
  Module.finrank ℂ (degreePiece (currentCarrier D.fullMatrix m d) r)

def carrierDim (N : ℕ) : ℕ := ∑ p ∈ Finset.range (N + 1), Module.finrank ℂ ((Carrier D) p N)

def sectorDim (N m d : ℕ) : ℕ := if m + 2 * d ≤ N then (auxDim D) m d (N - m - 2 * d) else 0

structure CasimirAuxiliaryBound : Prop where
  le : ∀ N, (carrierDim D) N ≤
    ∑ n ∈ Finset.range (N + 1), ∑ md ∈ Finset.antidiagonal n, (sectorDim D) N md.1 md.2

structure SectorFusionBound (F : ℕ → PowerSeries ℚ) : Prop where
  le : ∀ m d r, ((currentDim D) m d r : ℚ) ≤ coeff r (F m * (qpoch m)⁻¹ * (qpoch d)⁻¹)

structure MixedFusionRecursion (F Dem : ℕ → PowerSeries ℚ) : Prop where
  eq : ∀ n, Dem n = VOA.DeligneGaussian.mixed (X : PowerSeries ℚ) F n 0
\end{leancode}
\noindent File \leanpath{VOA/\allowbreak DeligneAudit/\allowbreak GaussianTransform.lean}.
\begin{leancode}
def mixed (f : ℕ → R) : ℕ → ℕ → R
  | 0, b => f b
  | a + 1, b => mixed f a (b + 1) + q ^ (2 * a + b + 1) * mixed f a b
\end{leancode}

\begin{dictionary}
Note the different relation families: \lean{Carrier D} uses \lean{D.rel} ($\ltwentyfourIns$, the presentation ring
$\ltwentyfourRt$), while the sector dimensions use \lean{D.fullMatrix} ($I_2$, Casimir included): in the auxiliary
algebra the Casimir is a separate particle and the current sector obeys all of $I_2$. \lean{qpoch m} is
$(q)_m$ in $\Q[[q]]$, \lean{coeff r} the $q^r$-coefficient. Sector $(m,d)$ contributes at conformal weight
$m+2d+p$ (currents weigh one, the Casimir two). The three structures are named \lean{Prop}s, never axioms;
\lean{CasimirAuxiliaryBound} is \emph{proved} below, the other two are discharged or reduced elsewhere.
\end{dictionary}

\begin{physics}
The weight-$N$ operator count is split by the number of current quanta $m$ and Casimir quanta $d$; each
sector is a free-particle Fock space with a fixed fusion fibre $\mathcal F_m$ and Bose statistics for $d$ identical
Casimir particles ($1/(q)_d$).
\end{physics}

\begin{status}
Definitions and named hypotheses.
\end{status}

\begin{theorem}[B1: the Casimir layer bound]\label{l24:thm:B1}
For any carrier data $\mathscr C$ and all $p,N,d$: the $d$-th layer of the Casimir-adic filtration of the bidegree
$(p,N)$ piece of $\ltwentyfourJ\ltwentyfourRt$ has dimension at most $\mathrm{auxDim}(N-p-2d,d,p)$ (and is $0$ if $p+2d>N$).
Consequently $\mathrm{carrierDim}(N)\le\sum_{m+2d\le N}\mathrm{sectorDim}(N,m,d)$ (\lean{CasimirAuxiliaryBound D}) holds.
\end{theorem}
\noindent File \leanpath{VOA/\allowbreak DeligneAudit/\allowbreak DeligneCarrierFiltration.lean}.
\begin{leancode}
structure CasimirLayerBound : Prop where
  le : ∀ p N d, Module.finrank ℂ ((casimirLayer D) p N d) ≤
    if p + 2 * d ≤ N then (auxDim D) (N - p - 2 * d) d p else 0
\end{leancode}
\noindent File \leanpath{VOA/\allowbreak DeligneAudit/\allowbreak DeligneCarrierLayerBound.lean}.
\begin{leancode}
theorem casimirLayerBound : (CasimirLayerBound D) := by

theorem casimirAuxiliaryBound : (CasimirAuxiliaryBound D) :=
  (casimirAuxiliaryBound_of_layerBound D) (casimirLayerBound D)
\end{leancode}

\begin{dictionary}
\lean{casimirLayer D p N d} is the $d$-th associated quotient of the filtration of the bidegree image by powers
of the jet ideal generated by the Hasse--Schmidt coefficients of $Q$ (\lean{casimirFiltration}, built from the
generic \lean{CasimirCostFiltration}). Its layers telescope to the bidegree dimension
(\lean{sum_casimirLayer_finrank}). No hypothesis other than the carrier data.

\textbf{Hasse--Schmidt coefficients.} For a polynomial $f\in\C[x_1,\dots,x_n]$, substitute $x_i\mapsto\sum_{r\ge0}x_i^{(r)}u^r$
and expand: $f\bigl(\sum_rx^{(r)}u^r\bigr)=\sum_{b\ge0}f_{(b)}u^b$. The coefficient $f_{(b)}$ is a polynomial in the jets,
$f_{(b)}=D^{(b)}f=D^bf/b!$ for the jet derivation $D$; the jet ideal of $(f)$ is generated by all $f_{(b)}$. For the
quadric $Q$, the coefficient $Q_{(b)}$ has conformal weight $b+2$.
\end{dictionary}

\textbf{Proof on paper.} Fix the bidegree $(p,N)$ and write $\mathcal J=\ltwentyfourJ\ltwentyfourRt$, $\mathcal J_{p,N}$ for
its bidegree piece, and $K\subseteq\ltwentyfourJ\operatorname{Sym}\ltwentyfourg$ for the jet ideal of $(Q)$, generated by the
Hasse--Schmidt coefficients $Q_{(b)}$, $b\ge0$.

\emph{Step 1 (the filtration).} Filter $\mathcal J_{p,N}$ by the images $K^d\mathcal J\cap\mathcal J_{p,N}$, $d=0,1,2,\dots$.
Every $Q_{(b)}$ has weight at least two, so $K^d$ vanishes in weight $N$ once $2d>N$; the filtration has finitely many
steps and the dimensions of its layers $\mathcal J^{(d)}_{p,N}=K^d/K^{d+1}$ add up to $\dim\mathcal J_{p,N}$. A layer is
spanned by classes of products of exactly $d$ Casimir jets $Q_{(b_1)}\cdots Q_{(b_d)}$ with a jet monomial in the currents.
Counting weights, the current part has $m=N-p-2d$ letters: the $d$ Casimir factors carry weight $2d$ and the $p$
derivatives are shared among all factors.

\emph{Step 2 (an auxiliary algebra that surjects onto the layer).} Add one variable $c$ of weight two (a ``Casimir
particle'') and impose three kinds of quadratic relations: all of $I_2$ (including $Q$) on the currents, the mixed
relations $x_ac$ for every $a$, and $c^2$. Map auxiliary jets to $\mathcal J$ by $c^{(b)}\mapsto Q_{(b)}$; this substitution
commutes with taking Hasse--Schmidt coefficients (because it is an algebra map compatible with the jet derivation). Restrict
to auxiliary jet monomials with exactly $d$ letters $c$ and pass to the $d$-th layer. Every auxiliary relation dies there:
\begin{itemize}
\item a relation of $\ltwentyfourIns$ on currents is already $0$ in $\ltwentyfourRt$;
\item the relation $Q$ on currents maps to some $Q_{(b')}$ times a monomial that already contains $d$ Casimir jets, which lies
  in $K^{d+1}$ and so is $0$ in the $d$-th layer;
\item the coefficient of order $b'$ of $x_ac$ maps to the coefficient of order $b'$ of $X_aQ$, which lies in the jet ideal of
  $(\ltwentyfourIns)$ by U3, hence is $0$;
\item likewise $c^2\mapsto Q^2\in(\ltwentyfourIns)$ (\lean{Q_sq_mem}, a consequence of U3).
\end{itemize}
So the map factors through the auxiliary quotient, and it is onto the layer because the layer is spanned by the images
of such monomials (\lean{auxToLayer_surjective}).

\emph{Step 3 (dualize).} A linear functional $\varphi$ on the layer gives coefficients
$\varphi(\text{image of an auxiliary monomial})$ on auxiliary jet monomials; they satisfy every auxiliary quadratic
coefficient equation (Step 2), and they determine $\varphi$ (surjectivity). By the coefficient duality of
Theorem~\ref{l24:thm:B4}, applied to the auxiliary relation family on $m+d$ particles, such coefficient families are the
symmetric tensor polynomials whose collision readbacks for all auxiliary relations vanish. Restrict to the \emph{label
pattern} ``currents in the first $m$ slots, $c$ in the last $d$''. By symmetry no information is lost (every labelling
with $d$ labels $c$ is a permutation of the pattern), and the readback conditions become: the equations of $I_2$ at
current--current collisions; vanishing at every current--Casimir collision (the relation $x_ac$ has a single nonzero entry,
so its readback is one component); vanishing at every Casimir--Casimir collision (from $c^2$). These are exactly the
conditions defining $\mathrm{auxiliaryCarrier}(I_2,m,d)$, in degree $p$. So the dual of the layer injects into
$\mathrm{degreePiece}(\mathrm{auxiliaryCarrier}(I_2,m,d),p)$ (\lean{layerToAux_injective}), whose dimension is
$\mathrm{auxDim}(m,d,p)$ with $m=N-p-2d$. Summing over $d$ and $p$ and reindexing $(p,d)\mapsto(m,d)$ gives
\lean{CasimirAuxiliaryBound}. \qed

\begin{physics}
Organize operators by how many Sugawara factors they contain. Since $T$ is nilpotent up to derivatives, each
extra $T$ behaves like a new particle with hard-core repulsion; the bound says there are no more operators than in
that particle model.
\end{physics}

\begin{status}
Proved, \textbf{unconditional}, for arbitrary carrier data; \texttt{\#print axioms casimirLayerBound}: standard three
(re-checked by the author, Section~\ref{l24:sec:verification}).
\end{status}

The docstrings of \lean{CasimirAuxiliaryBound} and \lean{SectorFusionBound} in \leanpath{DeligneCarrierChain.lean} are out of
date (the first is proved, the second does not use U4); see Section~\ref{l24:sec:defects}.

\begin{theorem}[D3: the Gaussian resummation]\label{l24:thm:D3}
For any family $F_m$ of power series over $\Q$ and every $N$,
\begin{equation*}
\sum_{m+d\le N}[q^N]\frac{q^{m+d+d(m+d)}F_m}{(q)_m(q)_d}=\sum_{n\le N}[q^N]\frac{q^n\,d_n}{(q)_n},\qquad
d_n=\sum_{d\le n}q^{nd}\begin{bmatrix}n\\d\end{bmatrix}_qF_{n-d};
\end{equation*}
and if $\mathrm{Dem}_n=M_{n,0}$ (U6), the right side equals $\sum_{n\le N}[q^N]q^n\mathrm{Dem}_n/(q)_n$.
\end{theorem}
\noindent File \leanpath{VOA/\allowbreak DeligneAudit/\allowbreak A2CarrierResummation.lean}.
\begin{leancode}
def sectorTerm (F : ℕ → PowerSeries ℚ) (m d : ℕ) : PowerSeries ℚ :=
  X ^ (m + d + d * (m + d)) * (F m * (qpoch m)⁻¹ * (qpoch d)⁻¹)

def sectorCoeff (F : ℕ → PowerSeries ℚ) (N : ℕ) : ℚ :=
  ∑ n ∈ Finset.range (N + 1), ∑ md ∈ Finset.antidiagonal n, coeff N (sectorTerm F md.1 md.2)

def demazureCoeffOf (Dem : ℕ → PowerSeries ℚ) (N : ℕ) : ℚ :=
  ∑ n ∈ Finset.range (N + 1), coeff N (X ^ n * Dem n * (qpoch n)⁻¹)

theorem sectorCoeff_eq_demazureCoeffOf (F : ℕ → PowerSeries ℚ) (N : ℕ) :
    sectorCoeff F N = demazureCoeffOf (demazureSeries F) N := by

theorem sectorCoeff_eq_of_mixed (F Dem : ℕ → PowerSeries ℚ)
    (hDem : ∀ n, Dem n = mixed (X : PowerSeries ℚ) F n 0) (N : ℕ) :
    sectorCoeff F N = demazureCoeffOf Dem N := by
\end{leancode}

\begin{dictionary}
Namespace \lean{VOA.DeligneAudit.CarrierChain}. The truncations $m+d\le N$ and $n\le N$ are harmless: the omitted
terms have zero $q^N$-coefficient since $q^{m+d+\dots}$ and $q^n$ already exceed $N$. \lean{demazureSeries F n} is
$d_n$. \lean{mixed X F n 0} is $M_{n,0}$; its closed form \lean{demazure_eq_gaussian} (Chapter~\ref{ch:L25}) is
the $d_n$ formula.
\end{dictionary}

\textbf{Proof on paper.} $[n,d](q)_d(q)_{n-d}=(q)_n$ gives
$q^nd_n/(q)_n=\sum_dq^{n+nd}F_{n-d}/((q)_{n-d}(q)_d)$; with $m=n-d$, $n+nd=m+d+d(m+d)$. Sum over the
antidiagonals. \qed

\begin{physics}
Regrouping the $(m,d)$ particle sectors by total particle number $n=m+d$ turns the free-particle count into a
sum over $n$ of fusion products of $n$ adjoint representations (Demazure modules) dressed by $1/(q)_n$.
\end{physics}

\begin{status}
Proved, unconditional (pure $q$-series); axioms standard.
\end{status}

\begin{theorem}[Carrier side of the chain, any type]\label{l24:thm:carrierside}
For carrier data $\mathscr C$, a table $\mathrm{hilb}$ of Demazure Hilbert series and a family $F$: if
\lean{SectorFusionBound D F} and $\mathrm{hilbSeries}(\mathrm{hilb})_n=M_{n,0}$ (U6), then
$\mathrm{carrierDim}(N)\le\sum_{n\le N}[q^N]q^n\,\mathrm{hilb}_n/(q)_n$ for every $N$.
\end{theorem}
\noindent File \leanpath{VOA/\allowbreak DeligneAudit/\allowbreak DeligneCarrierDemazure.lean}.
\begin{leancode}
def hilbSeries (hilb : ℕ → ℕ → ℕ) (n : ℕ) : PowerSeries ℚ :=
  PowerSeries.mk fun N => (hilb n N : ℚ)

theorem carrierDim_le_demazureSum (D : DeligneCarrierData) (hilb : ℕ → ℕ → ℕ)
    (F : ℕ → PowerSeries ℚ) (hF : SectorFusionBound D F)
    (hU6 : MixedFusionRecursion F (hilbSeries hilb)) (N : ℕ) :
    carrierDim D N ≤ demazureSumCoeff hilb N :=
\end{leancode}

\begin{dictionary}
\lean{demazureSumCoeff hilb N} is the natural-number Demazure sum of Chapter~\ref{ch:L25}
(\lean{DeligneCharacter.lean}); \lean{demazureSumCoeff_eq_demazureCoeffOf} identifies it with \lean{demazureCoeffOf}.
Hypotheses: the sector fusion bound (B2--B8 with U5; proved per type from U5 in Section~\ref{l24:sec:instances})
and U6. A sibling \lean{carrierCharacterLowerBound_of_sector} adds Chapter~\ref{ch:L25}'s
\lean{DemazureBoundFor W hilb} to produce the capstone's character input.
\end{dictionary}

\textbf{Proof on paper.} B1 (Theorem~\ref{l24:thm:B1}) bounds $\mathrm{carrierDim}(N)$ by the sum of sector dimensions;
the conductor shift (\lean{auxDim_conductor_add}) turns $\mathrm{auxDim}(m,d,\mathrm{cond}+r)$ into $\mathrm{currentDim}(m,d,r)$
and kills lower degrees; the exponent identity $m+d+d(m+d)=m+2d+(md+d(d-1))$ matches each sector with
$[q^N]\,\mathrm{sectorTerm}$; the sector fusion bound bounds each term; D3 with U6 resums. \qed

\begin{physics}
The classical operator count is at most the Demazure (fusion) count. The other half of the squeeze, that the Demazure count
is at most the quantum count (step (C) of Remark~\ref{l24:rem:arc}), is the subject of Chapter~\ref{ch:L25}.
\end{physics}

\begin{status}
Proved, conditional on \lean{SectorFusionBound D F} and \lean{MixedFusionRecursion}; axioms standard.
\end{status}

\begin{remark}[The shifted Mehler identity]\label{l24:rem:mehler}
The files \leanpath{A2CarrierMehlerFinite.lean} and \leanpath{A2CarrierMehler.lean} prove, with no hypothesis, the
$q$-series identity \lean{A2.shiftedMehler : ShiftedMehler}, i.e.\ $U_m(q)_\infty^3(1-q^{m+1})=m+1$ for
$U_m=\sum_sH_s(q)H_{s+m}(q)/(q)_s$ ($H$ the Rogers--Szeg\H{o} polynomials), via Rogers--Szeg\H{o} Vandermonde
expansions and two Euler sums; and \lean{starIdentity_of_fourTriangular}, which reduces the $A_2$ star identity
to Legendre's four-triangular-numbers identity. These belong logically to the character side and are discussed
in Chapter~\ref{ch:L25}.
\end{remark}

%======================================================================
\section{The sector chain: B5, B6, the generic comparison and B8}\label{l24:sec:sector}
%======================================================================

This section proves the sector fusion bound from U5 and a divisor comparison. The objects are the current
carrier $W_m$ of $I_2$ (Definition~\ref{l24:def:carrier}), the cyclic module $M_m\subseteq W_m$ generated by
$\delta_{c_0}^{\otimes m}$ under the loop operators, and the ring $B_m=\C[z_1,\dots,z_m]^{S_m}$ acting on both.

\paragraph{Why the divisor comparison is the crux.} The goal is a bound on $\dim(W_m)_r$ for each degree $r$. The
module $M_m$ is controlled: it is free over $B_m$ (B6), with a fibre whose size is bounded by U5, so its Hilbert series
is at most $F_m(q)/(q)_m$. What has to be shown is that $W_m$ is no bigger, i.e.\ $W_m=M_m$; this equality is what turns
a bound on the fibre into a bound on $\dim W_m$. The argument proceeds in four steps.
\begin{enumerate}
\item $M_m\subseteq W_m$ (B5), so only $W_m\subseteq M_m$ is at stake.
\item Away from collisions the two agree: $\Delta^LW_m\subseteq M_m$ (the generic comparison). So an element $x$ of $W_m$
  can fail to lie in $M_m$ only ``along'' the zero set of $\Delta$.
\item Since $B_m$ is a polynomial ring (a normal domain) and $M_m$ is free, membership of $x$ in $M_m$ can be tested at the
  height-one primes of $B_m$ alone (B8). Those not containing $\Delta$ are handled by step 2. The remaining ones contain
  $\Delta$; their zero sets are the images of the hyperplanes $z_i=z_j$, i.e.\ the loci where exactly two particles
  collide. This is the \emph{divisor comparison}.
\item At a generic point of $z_a=z_b$ the other particles are at distinct positions and act as spectators, so the
  divisor comparison reduces to a two-particle problem, a finite computation in $\ltwentyfourg^*\otimes\ltwentyfourg^*$.
  That computation is what the certificates of Section~\ref{l24:sec:divisor} do.
\end{enumerate}
Triple collisions and simultaneous pairs never need to be examined: they have codimension two in $\operatorname{Spec}B_m$
and are invisible to the height-one test. (In $z$-space a triple collision also has codimension two; the discriminant
$\Delta=0$, the union of all pair collisions, is a divisor both upstairs in $\C^m$ and downstairs in
$\operatorname{Spec}B_m=\C^m/S_m$.)

\begin{definition}[Loop operators, $M_m$, $B_m$, the discriminant]\label{l24:def:Mm}
For structure constants $f$ of $\ltwentyfourg$ in the chosen basis, the loop operator $x\otimes t^r$ acts on tensor
polynomials by the coadjoint action in each slot, weighted by $z_i^r$:
$(x\otimes t^r\cdot p)_t=-\sum_iz_i^r\sum_cf(x,t_i,c)\,p_{t[i\mapsto c]}$. Let $\delta_{c_0}^{\otimes m}$ be the
tensor polynomial equal to $1$ at the labelling $t\equiv c_0$ and $0$ elsewhere, $M_m$ the span of all loop words
applied to it, $B_m$ the symmetric polynomials and $\Delta=\prod_{i\neq j}(z_i-z_j)\in B_m$.
\end{definition}
\noindent File \leanpath{VOA/\allowbreak DeligneAudit/\allowbreak A2SectorLoop.lean}.
\begin{leancode}
def loopOp (m : ℕ) (x : σ) (r : ℕ) :
    TensorPolynomial (K := K) (σ := σ) m 0 →ₗ[K] TensorPolynomial (K := K) (σ := σ) m 0 where
  toFun p t := -∑ i, zr i ^ r * ∑ c, C (f x (t i) c) * p (Function.update t i c)

def deltaVec (m : ℕ) (c₀ : σ) : TensorPolynomial (K := K) (σ := σ) m 0 :=
  fun t => if t = fun _ => c₀ then 1 else 0

def Mspace (m : ℕ) (c₀ : σ) : Submodule K (TensorPolynomial (K := K) (σ := σ) m 0) :=
  Submodule.span K (Set.range fun l => applyWord f m l (deltaVec m c₀))
\end{leancode}
\noindent File \leanpath{VOA/\allowbreak DeligneAudit/\allowbreak A2SectorFiber.lean}.
\begin{leancode}
abbrev SymAlg := MvPolynomial.symmetricSubalgebra (Fin (m + 0)) K
\end{leancode}
\noindent File \leanpath{VOA/\allowbreak DeligneAudit/\allowbreak A2SectorB8.lean}.
\begin{leancode}
def discS (m : ℕ) : SymAlg (K := ℂ) m :=
  ⟨∏ i : Fin (m + 0), ∏ j : Fin (m + 0), if i = j then 1 else (X i - X j), by
\end{leancode}

\begin{dictionary}
\lean{zr i} is the root variable $z_i$; \lean{C} embeds scalars; \lean{Function.update t i c} relabels slot $i$
to $c$. \lean{applyWord f m l} applies a list of $(x,r)$ letters. Although the files are named \lean{A2…},
these definitions are generic in $K,\sigma,f$ and are used by all types. \lean{Fin (m + 0)} is $\{1..m\}$ written so
that it matches \lean{Particle m 0}. $\Delta$ is (up to sign) the square of the Vandermonde, hence symmetric (the
proof term). The coadjoint vector $\delta_{c_0}$ is the functional dual to $X_{c_0}=f_\theta$, which has weight
$-\theta$ in $\ltwentyfourg$; as an element of $\ltwentyfourg^*$ it is therefore a \emph{highest}-weight vector, of weight
$+\theta$. In the data of Definition~\ref{l24:def:sectordata} this is literal: with the coadjoint action
$(\lean{slotA}\,x\,v)_y=-\sum_df_{xyd}v_d$, \lean{raising_table} ($f_{e y c_0}=0$) says every raising vector kills
$\delta_{c_0}$, and \lean{cartan_table} ($f_{h y c_0}=-\theta(h)$ if $y=c_0$, else $0$) says $h$ acts by $+\theta(h)$. These
are the conditions \lean{LocalWeylRelG.raising} and \lean{cartan_zero} impose on the cyclic vector of
$W_{\mathrm{loc}}(m\theta)$.
\end{dictionary}

\begin{physics}
$M_m$ is the space of $m$-particle wave functions generated from the ``ground state'' where every particle sits in
the highest-weight coadjoint state $\delta_{c_0}$, by acting with the modes $x\otimes t^r$ of the current algebra at the particle positions:
an evaluation (fusion) representation of $\ltwentyfourg[t]$ at the points $z_1,\dots,z_m$.
\end{physics}

\begin{status}
Definitions.
\end{status}

\begin{definition}[Sector type data]\label{l24:def:sectordata}
Sector data extend carrier data by: structure constants $f$ satisfying the Jacobi identity; ad-stability of
$I_2=\ltwentyfourIns\oplus\C Q$; an index $c_0$ with no relation containing $X_{c_0}^2$; the local Weyl data for
$W_{\mathrm{loc}}(m\theta)$ (raising vectors killing $\delta_{c_0}$, Cartan vectors acting by $\theta(h)$, simple
lowering vectors with exponents $\theta(h_i)$, a weight function compatible with $f$ and a combinatorial
no-weight fact, a fusion set); and a \emph{depth} $L$ such that every basis vector of $\ltwentyfourg^*$ is a
combination of coadjoint words of length $\le L$ applied to $\delta_{c_0}$.
\end{definition}
\noindent File \leanpath{VOA/\allowbreak DeligneAudit/\allowbreak DeligneSectorData.lean}.
\begin{leancode}
def slotA {n : ℕ} (f : Fin n → Fin n → Fin n → ℂ) (x : Fin n) (v : Fin n → ℂ) : Fin n → ℂ :=
  fun y => -∑ d, f x y d * v d

def slotWords {n : ℕ} (f : Fin n → Fin n → Fin n → ℂ) (c₀ : Fin n) (L : ℕ) :
    Submodule ℂ (Fin n → ℂ) :=
  Submodule.span ℂ {v | ∃ l : List (Fin n), l.length ≤ L ∧ v = wordA f l (Pi.single c₀ 1)}

structure DeligneSectorData extends DeligneCarrierData where
  ρFin : Fintype ρ
  f : Fin n → Fin n → Fin n → ℂ
  jacobi : JacobiCoord f
  stable : RelationsStable f toDeligneCarrierData.fullMatrix
  c₀ : Fin n
  diag : ∀ r, toDeligneCarrierData.fullMatrix r (c₀, c₀) = 0
  raising : Fin n → Prop
  raising_table : ∀ e, raising e → ∀ y, f e y c₀ = 0
  cartan : Fin n → Prop
  thetaH : Fin n → ℕ
  cartan_table : ∀ h, cartan h → ∀ y, f h y c₀ = if y = c₀ then -(thetaH h : ℂ) else 0
  hθ : Fin n
  hθ_cartan : cartan hθ
  hθ_ne : thetaH hθ ≠ 0
  lowering : Fin n → Prop
  lowExp : Fin n → ℕ
  R : ℕ
  wt : Fin n → Fin R → ℤ
  table_wt : ∀ x y c, wt c ≠ wt x + wt y → f x y c = 0
  lowering_noweight : ∀ i, lowering i → ∀ (m : ℕ) (t : Fin m → Fin n),
    ∑ j, wt (t j) ≠ m • wt c₀ - (m * lowExp i + 1) • wt i
  fusion : Fin n → Prop
  depth : ℕ
  cyclic : ∀ c, (Pi.single c 1 : Fin n → ℂ) ∈ slotWords f c₀ depth
  cartanN : Fin n → Prop := fun _ => False
  thetaN : Fin n → ℕ := fun _ => 0
  cartanN_table : ∀ h, cartanN h → ∀ y, f h y c₀ = if y = c₀ then (thetaN h : ℂ) else 0 := by
    intro h hh; exact False.elim hh
\end{leancode}
\noindent File \leanpath{VOA/\allowbreak DeligneAudit/\allowbreak A2SectorRep.lean}.
\begin{leancode}
def JacobiCoord : Prop :=
  ∀ a b x c, (∑ d, f a x d * f b d c) - (∑ d, f b x d * f a d c) = -∑ e, f a b e * f e x c
\end{leancode}
\noindent File \leanpath{VOA/\allowbreak DeligneAudit/\allowbreak A2SectorLoop.lean}.
\begin{leancode}
def RelationsStable (Q : ρ → σ × σ → K) : Prop :=
  ∀ x r, quadraticRelation (adMatrix f x (Q r)) ∈
    Submodule.span K (Set.range fun s => quadraticRelation (Q s))
\end{leancode}

\begin{dictionary}
\lean{JacobiCoord f} is the Jacobi identity in coordinates, in the form needed for the loop operators to satisfy
the current-algebra relations (\lean{loopOp_comm}). \lean{RelationsStable} says the derivation $\operatorname{ad}_x$
maps each relation of $I_2$ into the \emph{linear span} of the relations. \lean{cartanN}/\lean{thetaN} (with
defaults) cover Cartan basis vectors with $\theta(h)<0$, which occur in the $G_2$ and $F_4$ bases.
\textbf{The weight fields.} \lean{R} and \lean{wt} assign to each basis vector of $\ltwentyfourg$ a weight in
$\Z^R$, a coarse version of the root lattice grading (for instance the coordinates of a root in the simple roots); the
same function labels the basis of $\ltwentyfourg^*$. \lean{table_wt} says the structure constants respect it:
$f(x,y,c)=0$ unless $\mathrm{wt}(c)=\mathrm{wt}(x)+\mathrm{wt}(y)$. Consequently the coadjoint action of $x$ in one slot
lowers the total weight $\sum_j\mathrm{wt}(t_j)$ of a labelling $t$ by $\mathrm{wt}(x)$, and $k$ applications of the
lowering vector $x=i$ to $\delta_{c_0}^{\otimes m}$ (total weight $m\,\mathrm{wt}(c_0)$) produce a tensor supported on labellings of
total weight $m\,\mathrm{wt}(c_0)-k\,\mathrm{wt}(i)$. \lean{lowering_noweight} says that for $k=m\,\theta(h_i)+1$ (with
$\theta(h_i)$ = \lean{lowExp i}) \emph{no} labelling of the $m$ slots has that total weight. Hence
$(f_i\otimes1)^{m\theta(h_i)+1}\delta_{c_0}^{\otimes m}=0$ (\lean{lowering_deltaVec}), which is the integrability relation of
$W_{\mathrm{loc}}(m\theta)$ and is used in B6. The \lean{fusion} set is unconstrained data; see the dictionary of
Definition~\ref{l24:def:U5}. All fields are proved per type by certificates (Section~\ref{l24:sec:instances}); the
structure has no unprovable or contradictory field (the seven instances exist).
\end{dictionary}

\begin{physics}
The Lie algebra in a weight basis, the highest-weight coadjoint state $\delta_{c_0}$ (the highest-weight vector of the
coadjoint representation; under the invariant form it is $e_\theta$, a point of the minimal orbit, the classical value of
the moment map in a highest-weight state),
the data of the local Weyl module $W_{\mathrm{loc}}(m\theta)$ (the fusion of $m$ adjoint representations), and the
fact that the coadjoint representation is reached from its highest-weight vector in few steps.
\end{physics}

\begin{status}
Definition; seven instances, each \texttt{\#print axioms}-standard.
\end{status}

\begin{definition}[Current modules, local Weyl relations, input U5]\label{l24:def:U5}
A \emph{current module} is a vector space with operators $\rho(x\otimes t^r)$, $r\ge0$, satisfying
$[x\otimes t^r,y\otimes t^s]=\sum_ef(x,y,e)\,e\otimes t^{r+s}$. For a vector $w$, $\mathrm{cycDeg}(w,b)$ is the span of
loop words of total $t$-degree $b$ applied to $w$. The \emph{local Weyl and fusion relations} for $m\theta$ are:
raising vectors $\otimes t^r$ kill $w$; Cartan $h\otimes t^r$ kills $w$ for $r\ge1$ and acts by $m\theta(h)$ for
$r=0$; $(f_i\otimes1)^{m\theta(h_i)+1}w=0$; $(x\otimes t^m)w=0$ for $x$ in the fusion set. \emph{Input U5} for a
family $F$: every such module has $\dim\mathrm{cycDeg}(w,b)\le[q^b]F_m$, and
$\sum_{b<B}\dim\mathrm{cycDeg}(w,b)\le(\dim\ltwentyfourg)^m$ for every $B$.
The cyclic module generated by $w$ subject only to the local Weyl relations (the first three) is the \emph{local Weyl
module} $W_{\mathrm{loc}}(m\theta)$ of Chari and Pressley, the largest graded cyclic $\ltwentyfourg[t]$-module of highest weight
$m\theta$ that is finite-dimensional; adding the fusion relations gives the \emph{fusion module} $\mathcal F_m$. Every
module satisfying the relations, with $w$ cyclic, is a quotient of $\mathcal F_m$, so U5 is a statement about $\mathcal F_m$.
\end{definition}
\noindent File \leanpath{VOA/\allowbreak DeligneAudit/\allowbreak A2SectorWeyl.lean}.
\begin{leancode}
structure CurrentModule [Fintype σ] (V : Type*) [AddCommGroup V] [Module K V] where
  op : σ → ℕ → V →ₗ[K] V
  comm : ∀ a b r s v, op a r (op b s v) - op b s (op a r v) = ∑ e, f a b e • op e (r + s) v

def cycDeg (ρ : CurrentModule f V) (w : V) (b : ℕ) : Submodule K V :=
  Submodule.span K {v | ∃ l, SectorGraded.wordDeg l = b ∧ v = wordOp ρ l w}
\end{leancode}
\noindent File \leanpath{VOA/\allowbreak DeligneAudit/\allowbreak DeligneSectorWeyl.lean}.
\begin{leancode}
structure LocalWeylRelG {V : Type*} [AddCommGroup V] [Module ℂ V]
    (ρ : CurrentModule S.f V) (w : V) (m : ℕ) : Prop where
  raising : ∀ e, S.raising e → ∀ r, ρ.op e r w = 0
  cartan_pos : ∀ h, S.cartan h → ∀ r, 1 ≤ r → ρ.op h r w = 0
  cartan_zero : ∀ h, S.cartan h → ρ.op h 0 w = ((m * S.thetaH h : ℕ) : ℂ) • w
  lowering : ∀ i, S.lowering i → (ρ.op i 0)^[m * S.lowExp i + 1] w = 0
  fusion : ∀ x, S.fusion x → ρ.op x m w = 0
  cartanN_pos : ∀ h, S.cartanN h → ∀ r, 1 ≤ r → ρ.op h r w = 0
  cartanN_zero : ∀ h, S.cartanN h → ρ.op h 0 w = -((m * S.thetaN h : ℕ) : ℂ) • w

structure FusionModuleBoundG (F : ℕ → PowerSeries ℚ) : Prop where
  le : ∀ (m : ℕ) (V : Type) [AddCommGroup V] [Module ℂ V] (ρ : CurrentModule S.f V) (w : V),
    LocalWeylRelG S ρ w m → ∀ b, (Module.finrank ℂ (cycDeg ρ w b) : ℚ) ≤ PowerSeries.coeff b (F m)
  total : ∀ (m : ℕ) (V : Type) [AddCommGroup V] [Module ℂ V] (ρ : CurrentModule S.f V) (w : V),
    LocalWeylRelG S ρ w m → ∀ B : ℕ,
      ∑ b ∈ Finset.range B, Module.finrank ℂ (cycDeg ρ w b) ≤ S.n ^ m
\end{leancode}

\begin{dictionary}
\lean{(ρ.op i 0)^[k]} is the $k$-fold iterate. The relations are Chari--Pressley's presentation of
$W_{\mathrm{loc}}(m\theta)$ ($\mathfrak n^+[t]w=0$, $(h\otimes t^s)w=\delta_{s0}m\theta(h)w$,
$(f_i)^{m\theta(h_i)+1}w=0$) plus the fusion relations of Feigin and Ravinder. Every module satisfying them is a
quotient of the fusion module $\mathcal F_m=W_{\mathrm{loc}}(m\theta)/\langle(f_\alpha\otimes t^m)w\rangle$, so U5 is exactly ``$\mathcal F_m$ has graded
character $\le F_m(q)$ and dimension $\le(\dim\ltwentyfourg)^m$''. The published source of the dimension half is
Kus--Littelmann, Pacific J.\ Math.\ 278, Theorem~4.2 (one step, on p.~437, flagged by the project's source audit);
Feigin's earlier proof has a gap, and Ravinder's Proposition~5 is Feigin's statement, so it is not independent
(Chapter~\ref{ch:L27}).

\textbf{Two remarks on the formulation.} (1) \lean{Module.finrank} of an infinite-dimensional space is $0$ in
Mathlib, so for a hypothetical module with an infinite-dimensional degree piece the inequalities say nothing. This
makes the hypothesis \emph{weaker} (easier to discharge), and is harmless for its use: it is applied only to the
finite-dimensional fibre of $M_m$. (2) The fusion set is free data. Ravinder's set is
$\{f_\alpha:\langle\alpha,\theta^\vee\rangle=1\text{ or }\alpha=\theta\}$; $E_6$, $E_7$, $E_8$ use exactly that, $A_2$, $D_4$,
$G_2$, $F_4$ use all negative root vectors. Every positive root $\alpha\neq\theta$ has
$\langle\alpha,\theta^\vee\rangle\in\{0,1\}$, so the short roots of $G_2$ and $F_4$ with $\langle\alpha,\theta^\vee\rangle=1$ are
already in Ravinder's set, and the extra vectors are the $f_\alpha$ with $\langle\alpha,\theta^\vee\rangle=0$. For these,
$f_\alpha(r)w=0$ holds for every $r$ in any module with the local Weyl relations, by the $h(r)$-bracket argument of
Lemma~\ref{l27:lem:truncweyl}. So the two readings impose the same relations, and give the same module, for every type
(Chapter~\ref{ch:L27}; Remark~\ref{l26:rem:fusionset}). The chain proves the fusion
relations for every negative root vector on the fibre, so it can use either reading.
\end{dictionary}

\begin{physics}
$W_{\mathrm{loc}}(m\theta)$ is the space of states obtained by fusing $m$ copies of the adjoint representation at
colliding points; $\mathcal F_m$ is the further quotient in which each particle has a bounded number of excitations. U5 says
the fused space is no larger than the unfused tensor product $\ltwentyfourg^{\otimes m}$, with a definite energy grading.
\end{physics}

\begin{status}
Definitions; \lean{FusionModuleBoundG} is a named hypothesis (literature input U5). Chapter~\ref{ch:L26} reduces it to a
total-dimension bound plus Demazure realizations.
\end{status}

\begin{theorem}[B5: $M_m\subseteq W_m$]\label{l24:thm:B5}
Let $K$ be a field of characteristic zero, $\sigma$ a finite basis, $f$ arbitrary structure constants and $\rho$ a
finite set of relations. If the span of the relations is stable under the coadjoint action of $f$
(\lean{RelationsStable}) and no relation contains $X_{c_0}^2$, then $M_m$ lies in the current carrier of the relations,
for every $m$. No Jacobi identity for $f$ is assumed.
\end{theorem}
\noindent File \leanpath{VOA/\allowbreak DeligneAudit/\allowbreak A2SectorLoop.lean}.
\begin{leancode}
theorem Mspace_le [CharZero K] {Q : ρ → σ × σ → K} [Fintype ρ] (hQ : RelationsStable f Q)
    (c₀ : σ) (hdiag : ∀ r, Q r (c₀, c₀) = 0) : Mspace f m c₀ ≤ currentCarrier Q m 0 := by
\end{leancode}

\begin{dictionary}
Generic in the field (characteristic zero), basis, structure constants and relations. In the sector chain it is
applied as \lean{Mspace_le S.f S.stable S.c₀ S.diag}. The hypotheses are exactly \lean{[CharZero K]},
\lean{[Fintype ρ]}, \lean{hQ : RelationsStable f Q} and \lean{hdiag}: the section variables of the file are only
\lean{{K σ ρ : Type*} [Field K] [Fintype σ] [DecidableEq σ]} and \lean{(f : σ → σ → σ → K)}, with no Jacobi
hypothesis, and the stability step \lean{loopOp_mem} uses only \lean{RelationsStable}. So the Lean theorem is
stronger than a statement assuming Jacobi. The Jacobi identity (\lean{JacobiCoord}) enters the chain at B6, through
\lean{loopOp_comm} in \leanpath{A2SectorRep.lean} (commutation of loop operators).
\end{dictionary}

\textbf{Proof on paper.} $\delta_{c_0}^{\otimes m}$ is in $W_m$: it is symmetric, and a collision readback picks the
coefficient of $X_{c_0}^2$ in a relation, which is $0$. $W_m$ is stable under loop operators: on the diagonal
$z_i=z_j$ the pair action is the diagonal coadjoint action, and the readback of $\operatorname{ad}$ applied to a
relation is a combination of readbacks of relations, by ad-stability (\lean{loopOp_mem}). \qed

\begin{physics}
Fusing particles in the highest-weight coadjoint state and exciting them by current modes never violates the collision rules: the
collision rules are flavour-covariant.
\end{physics}

\begin{status}
Proved; axioms standard.
\end{status}

\begin{lemma}[Symmetrization]\label{l24:lem:reynolds}
Let $P_mM_m$ be the $P_m$-span of $M_m$. A symmetric tensor polynomial that lies in $P_mM_m$ lies in $M_m$.
\end{lemma}
\noindent File \leanpath{VOA/\allowbreak DeligneAudit/\allowbreak DeligneSectorGeneric.lean}.
\begin{leancode}
theorem mem_Mspace_of_symmetric {v : TPs S m} (hv : v ∈ MP S m) (hs : Symmetric v) :
    v ∈ Mspace S.f m S.c₀ := by
\end{leancode}

\begin{dictionary}
\lean{MP S m} is $P_mM_m$ (\lean{Submodule.span (Pm m) (Mspace …)}) and \lean{Symmetric} is the symmetry of
Definition~\ref{l24:def:carrier}. The statement is exactly the lemma, for the sector data $S$; no further hypothesis.
\end{dictionary}

\textbf{Proof on paper.} Let $S_m$ act on tensor polynomials by permuting slots and roots simultaneously,
$(g\cdot v)_t=g(v_{t\circ g})$; symmetric means $g\cdot v=v$ for all $g$, and on a product $g\cdot(pv)=g(p)\,(g\cdot v)$ for
$p\in P_m$. Write $v=\sum_ip_iy_i$ with $p_i\in P_m$ and $y_i\in M_m$, and average over $S_m$:
\begin{equation*}
v=\frac1{m!}\sum_{g\in S_m}g\cdot v=\frac1{m!}\sum_{g}\sum_ig(p_i)\,(g\cdot y_i)=\sum_i\Bigl(\frac1{m!}\sum_gg(p_i)\Bigr)y_i .
\end{equation*}
The last equality is the point: each $y_i\in M_m$ is itself symmetric, because $M_m\subseteq W_m$ (Theorem~\ref{l24:thm:B5})
and elements of $W_m$ are symmetric; so $g\cdot y_i=y_i$. The averages $\frac1{m!}\sum_gg(p_i)$ are symmetric polynomials, i.e.\
lie in $B_m$, and $M_m$ is a $B_m$-module (\lean{smul_mem_Mspace}). Hence $v\in M_m$. (It is not enough that averaging
commutes with the loop operators: averaging a general element of $P_mM_m$ only returns an element of $P_mM_m$. What is used is
that the generators $y_i$ are already invariant.) \qed

\begin{physics}
A wave function that is symmetric under exchange of identical particles, and is built from the ground state with
position-dependent coefficients, can be rebuilt with symmetric coefficients: Bose symmetry of the building blocks lets one
average the coefficients over particle exchanges.
\end{physics}

\begin{status}
Proved; axioms standard (re-checked). Used three times: in the generic comparison, in B6 and in the reduction of the divisor
comparison to pairs.
\end{status}

\begin{theorem}[The generic comparison: $\Delta^L W_m\subseteq M_m$]\label{l24:thm:generic}
For sector data of depth $L$ and $x\in W_m$, $\Delta^L\cdot x\in M_m$.
\end{theorem}
\noindent File \leanpath{VOA/\allowbreak DeligneAudit/\allowbreak DeligneSectorGeneric.lean}.
\begin{leancode}
theorem genericComparisonG (m : ℕ) {x : TPs S m} (hx : x ∈ currentCarrier S.fullMatrix m 0) :
    (discS m : SymAlg (K := ℂ) m) ^ S.depth • x ∈ Mspace S.f m S.c₀ := by
\end{leancode}

\begin{dictionary}
\lean{TPs S m} is the tensor polynomials for the type. The hypothesis \lean{hx} is used only through the symmetry
of $x$; the inclusion holds for every symmetric tensor (and $\Delta^Lx\in P_mM_m$ for every tensor, symmetric or not).
\end{dictionary}

\textbf{Proof on paper.} Let $P_mM_m$ be the $P_m$-span of $M_m$ (polynomial, not only symmetric, coefficients). Four steps.

\emph{Step 1 (interpolation isolates one slot).} Fix a slot $i$ and put $\lambda_i(t)=\prod_{j\neq i}(t-z_j)=\sum_rc_rt^r$, with
coefficients $c_r\in P_m$. Since $x\otimes t^r$ acts by $\sum_kz_k^r\,x^{(k)}$ ($x^{(k)}$ the coadjoint action in slot $k$),
\begin{equation*}
\sum_rc_r\,(x\otimes t^r)=\sum_k\lambda_i(z_k)\,x^{(k)}=\lambda_i\,x^{(i)},\qquad \lambda_i=\lambda_i(z_i)=\prod_{j\neq i}(z_i-z_j),
\end{equation*}
because $\lambda_i(z_k)=0$ for $k\neq i$ (\lean{interp}). The coefficients $c_r$ are not symmetric, so this operator preserves
$P_mM_m$ but not $M_m$; that is why the argument runs in $P_mM_m$ first.

\emph{Step 2 (one slot, by cyclicity).} By the depth-$L$ cyclicity of the sector data, every basis vector $e_c$ of
$\ltwentyfourg^*$ is a combination of words of length $\le L$ in the coadjoint action applied to $\delta_{c_0}$. Each letter costs
one factor $\lambda_i$ (Step 1), and shorter words are padded with extra factors $\lambda_i$. So if $Y\,\delta_{t}$ lies in
$P_mM_m$ for a labelling $t$ with $t_i=c_0$ and some $Y\in P_m$, then $\lambda_i^LY\,\delta_{t[i\mapsto c]}\in P_mM_m$ for every $c$
(\lean{slot_all}).

\emph{Step 3 (all slots).} Start from $\delta_{c_0}^{\otimes m}\in M_m$ and change the labels one slot at a time: for every
labelling $t$, $\bigl(\prod_i\lambda_i^L\bigr)\delta_t\in P_mM_m$ (\lean{all_slots}). Since $\prod_i\lambda_i=\prod_i\prod_{j\neq i}(z_i-z_j)=\Delta$,
$\Delta^L\delta_t\in P_mM_m$, and hence $\Delta^Lx\in P_mM_m$ for \emph{every} tensor polynomial $x=\sum_tx_t\,\delta_t$
(\lean{disc_pow_mem_MP}).

\emph{Step 4 (descend to $M_m$).} If $x\in W_m$ then $x$ is symmetric, and so is $\Delta^Lx$ because $\Delta\in B_m$.
Lemma~\ref{l24:lem:reynolds} gives $\Delta^Lx\in M_m$. \qed

\begin{physics}
When all particles are at distinct points, each can be excited independently, so every wave function is reachable,
up to the discriminant factor that vanishes when particles collide.
\end{physics}

\begin{status}
Proved, unconditional given the sector data; axioms standard.
\end{status}

\begin{theorem}[B6: freeness from U5]\label{l24:thm:B6}
For sector data $S$ and a family $F$ satisfying U5, $M_m$ is a finite free $B_m$-module for every $m$.
\end{theorem}
\noindent File \leanpath{VOA/\allowbreak DeligneAudit/\allowbreak DeligneSectorB6.lean}.
\begin{leancode}
structure MFreenessG : Prop where
  free : ∀ m, Module.Free (SymAlg (K := ℂ) m) (MspaceR S m)
  finite : ∀ m, Module.Finite (SymAlg (K := ℂ) m) (MspaceR S m)

theorem mFreeness (hU5 : FusionModuleBoundG S F) : MFreenessG S := by
\end{leancode}

\begin{dictionary}
\lean{MspaceR S m} is $M_m$ regarded as a $B_m$-submodule (it is $B_m$-stable: power sums act as
$h_\theta\otimes t^r$, \lean{smul_mem_Mspace}). Hypothesis: U5 only. The A$_2$ pilot is
\lean{VOA.DeligneAudit.A2.mFreeness} in \leanpath{A2SectorB6.lean}.
\end{dictionary}

\textbf{Proof on paper.} Grade $M_m$ by polynomial degree; $B_m$ is graded with Hilbert series $1/(q)_m$, and
$B_+$ is its ideal of positive degree. Five steps.

\emph{Step 1 (the fibre is a current module).} The loop operators commute with multiplication by symmetric polynomials, so
they act on the \emph{Nakayama fibre} $\bar M=M_m/B_+M_m$, the quotient by everything with a positive-degree symmetric factor.
Let $w$ be the image of $\delta_{c_0}^{\otimes m}$; it generates $\bar M$.

\emph{Step 2 (local Weyl and fusion relations for $w$).} Each relation of Definition~\ref{l24:def:U5} holds for $w$:
\begin{itemize}
\item raising vectors kill $\delta_{c_0}$ in every slot (\lean{raising_table}), so $(e\otimes t^r)\delta_{c_0}^{\otimes m}=0$;
\item a Cartan element acts on $\delta_{c_0}$ by $\theta(h)$ in each slot, so $(h\otimes t^r)\delta_{c_0}^{\otimes m}=\theta(h)\,p_r(z)\,\delta_{c_0}^{\otimes m}$
  with $p_r=\sum_iz_i^r$; for $r\ge1$ this lies in $B_+M_m$ (zero in $\bar M$), and for $r=0$ it is $m\theta(h)\,w$;
\item lowering: $(f_i\otimes1)^{m\theta(h_i)+1}\delta_{c_0}^{\otimes m}=0$ by the weight argument in the dictionary of
  Definition~\ref{l24:def:sectordata};
\item fusion (Vieta): each $z_k$ is a root of $\prod_i(T-z_i)=\sum_{j=0}^m(-1)^je_jT^{m-j}$, so
  $z_k^m=-\sum_{j=1}^m(-1)^je_jz_k^{m-j}$ with $e_j$ the elementary symmetric polynomials. Summing over slots,
  $(x\otimes t^m)=-\sum_{j=1}^m(-1)^je_j\,(x\otimes t^{m-j})$ as operators, and every term has a factor $e_j\in B_+$. Hence
  $(x\otimes t^m)w=0$ in $\bar M$, for every $x$ (\lean{loopOp_vieta}).
\end{itemize}

\emph{Step 3 (an upper bound on generators).} By U5 (its \lean{total} clause) the fibre has total dimension $\le n^m$, and in
particular vanishes in high degree. By graded Nakayama, lifting a homogeneous basis of $\bar M$ gives at most $n^m$ homogeneous
elements that generate $M_m$ over $B_m$.

\emph{Step 4 ($n^m$ independent elements).} Fix the point $z^0=(0,1,\dots,m-1)$, with distinct coordinates, and the Lagrange
polynomial $\ell(z)=\prod_i\prod_{j\neq i}(z_i-z^0_j)/(z^0_i-z^0_j)$, which is $1$ at $z^0$ and $0$ at every permutation
$g(z^0)\neq z^0$. For each labelling $t$ put
\begin{equation*}
v_t=\sum_{g\in S_m}g\cdot\bigl(\ell\,\delta_t\bigr),
\end{equation*}
a \emph{symmetrized} interpolated basis tensor: $v_t$ is symmetric by construction, and at $z^0$ its components are
$(v_t)_{t'}(z^0)=\delta_{t,t'}$ (only $g=1$ survives). By the generic comparison, applied to the symmetric tensor $v_t$
(Theorem~\ref{l24:thm:generic} and its dictionary: $\Delta^Lv_t\in P_mM_m$ by Step 3 of that proof, and then
Lemma~\ref{l24:lem:reynolds}), the $n^m$ elements $\Delta^Lv_t$ lie in $M_m$. They are $B_m$-linearly independent: in a relation
$\sum_tb_t\Delta^Lv_t=0$ with $b_t\in B_m$, the matrix of components $\bigl((\Delta^Lv_t)_{t'}\bigr)$ has determinant whose
value at $z^0$ is $\Delta(z^0)^{Ln^m}\neq0$; so the determinant is a nonzero element of the domain $P_m$, and Cramer's rule gives
every $b_t=0$.

\emph{Step 5 (rank count).} So $M_m$ has rank at least $n^m$ over $\operatorname{Frac}(B_m)$ and is generated by at most $n^m$
elements. A generating family of $k$ elements of a module that contains $k'\ge k$ independent elements is itself independent
over a domain (otherwise the rank would be $<k\le k'$; \lean{linearIndependent_of_span_card}). Hence the at most $n^m$
generators of Step 3 are exactly $n^m$ and form a $B_m$-basis. \qed

\begin{physics}
The particle wave functions form a free module over symmetric functions of the positions: the fusion product does
not depend (in dimension) on where the particles are, and at generic positions it is the full tensor product
$\ltwentyfourg^{\otimes m}$. Only an upper bound on the fused space is needed; the lower bound comes from generic points.
\end{physics}

\begin{status}
Proved, conditional on U5; axioms standard. Only U5's upper bounds are used: the lower bound on the rank comes from Step 4.
\end{status}

\begin{theorem}[B8: $W_m=M_m$ from freeness and the divisor comparison]\label{l24:thm:B8}
Suppose $M_m$ is finite free (B6) and the \emph{divisor comparison} holds: for every $x\in W_m$ and every height-one
prime $\mathfrak p\subset B_m$ containing $\Delta$, some $s\notin\mathfrak p$ has $sx\in M_m$. Then $W_m\subseteq M_m$.
\end{theorem}
\noindent File \leanpath{VOA/\allowbreak DeligneAudit/\allowbreak DeligneSectorChain.lean}.
\begin{leancode}
structure DivisorComparisonG : Prop where
  le : ∀ m, ∀ x ∈ currentCarrier S.fullMatrix m 0, ∀ p : HeightOnePrime (SymAlg (K := ℂ) m),
    discS m ∈ p.asIdeal → ∃ s : SymAlg (K := ℂ) m, s ∉ p.asIdeal ∧ s • x ∈ Mspace S.f m S.c₀

theorem carrierGenerationG (hF : MFreenessG S) (hD : DivisorComparisonG S) (m : ℕ) :
    currentCarrier S.fullMatrix m 0 ≤ Mspace S.f m S.c₀ := by
\end{leancode}
\noindent File \leanpath{VOA/\allowbreak Commutative/\allowbreak HeightOneGlobalization.lean}.
\begin{leancode}
structure HeightOnePrime (R : Type*) [CommRing R] where
  asIdeal : Ideal R
  isPrime : asIdeal.IsPrime
  height_eq_one : asIdeal.height = 1
\end{leancode}
\noindent File \leanpath{VOA/\allowbreak DeligneAudit/\allowbreak A2SectorGlobal.lean}.
\begin{leancode}
theorem mem_of_local_witnesses [UniqueFactorizationMonoid R] [IsNoetherianRing R] [Finite ι]
    (htf : ∀ (s : R) (x : N), s ≠ 0 → s • x = 0 → x = 0)
    (M : Submodule R N) (b : Module.Basis ι R M) (x : N)
    (h0 : ∃ s : R, s ≠ 0 ∧ s • x ∈ M)
    (hloc : ∀ p : HeightOnePrime R, ∃ s : R, s ∉ p.asIdeal ∧ s • x ∈ M) : x ∈ M := by
\end{leancode}

\begin{dictionary}
\lean{HeightOnePrime} is a prime of height exactly one, with no dimension assumption on $R$ (unlike Mathlib's
Dedekind \lean{HeightOneSpectrum}). $B_m\cong\C[e_1..e_m]$ is a noetherian UFD (instances in \leanpath{A2SectorB8.lean}).
\lean{DivisorComparisonG} is a named per-type hypothesis structure; Section~\ref{l24:sec:divisor} proves it for all
seven types. \textbf{Quantifiers.} The theorem is stated for a fixed $m$, but its hypotheses \lean{hF : MFreenessG S} and
\lean{hD : DivisorComparisonG S} are quantified over \emph{all} $m$; so the Lean theorem is formally weaker than the
fixed-$m$ statement above (it needs freeness and the divisor comparison in every arity). This is harmless for the
use: both are proved (respectively, derived from U5) for all $m$ at once.
\end{dictionary}

\textbf{Proof on paper.} Let $(b_i)$ be a $B_m$-basis of $M_m$ (B6) and $x\in W_m$. By the generic comparison
(Theorem~\ref{l24:thm:generic}), $\Delta^Lx\in M_m$; write $\Delta^Lx=\sum_ia_ib_i$ with $a_i\in B_m$. Then, in the fraction
field $K=\operatorname{Frac}B_m$, the coordinates of $x$ are $a_i/\Delta^L$; we must show that they lie in $B_m$. A noetherian
UFD is the intersection of its localizations at height-one primes, so it suffices to show $a_i/\Delta^L\in(B_m)_{\mathfrak p}$
for every height-one prime $\mathfrak p$, and for that it suffices to find one $s\notin\mathfrak p$ with $sx\in M_m$: writing
$sx=\sum_ia'_ib_i$ with $a'_i\in B_m$ and comparing coordinates in $K$ (the $b_i$ stay independent there, since $M_m$ is
torsion-free), $a_i/\Delta^L=a'_i/s\in(B_m)_{\mathfrak p}$.
\begin{itemize}
\item If $\Delta\notin\mathfrak p$, take $s=\Delta^L$ (the generic comparison).
\item If $\Delta\in\mathfrak p$, the divisor comparison supplies $s$.
\end{itemize}
In the Lean lemma \lean{mem_of_local_witnesses} these are the hypotheses \lean{hloc}; the hypothesis \lean{h0}, one global
nonzero $s$ with $sx\in M_m$, is $s=\Delta^L$ again; \lean{htf} is torsion-freeness of the ambient tensor polynomials. So
every coordinate lies in $B_m$ and $x\in M_m$. \qed

\begin{physics}
The statement is about membership, not about regularity of a single function: an element of the carrier $W_m$ (a
collection of wave functions obeying the collision rules) is shown to be generated, with symmetric polynomial
coefficients, from the ground state. A priori it is generated only after multiplying by a power of $\Delta$, i.e.\ away
from all collisions. Normality of $B_m$ says that a ``denominator'' can only sit on a divisor of $\operatorname{Spec}B_m$;
the divisors on which $\Delta$ vanishes are the simple two-particle collisions, so only the two-body problem at
$z_a=z_b$, with the other particles as spectators, has to be solved. Triple collisions and pairs of simultaneous collisions
have codimension two and impose nothing.
\end{physics}

\begin{status}
Proved; axioms standard.
\end{status}

\begin{theorem}[The sector fusion bound from U5 and the divisor comparison]\label{l24:thm:sfb}
For sector data $S$, U5 for $F$ and the divisor comparison imply \lean{SectorFusionBound S.toDeligneCarrierData F}.
\end{theorem}
\noindent File \leanpath{VOA/\allowbreak DeligneAudit/\allowbreak DeligneSectorChain.lean}.
\begin{leancode}
theorem sectorFusionBoundG (F : ℕ → PowerSeries ℚ) (hU5 : FusionModuleBoundG S F)
    (hD : DivisorComparisonG S) : SectorFusionBound S.toDeligneCarrierData F := by
\end{leancode}
\noindent File \leanpath{VOA/\allowbreak DeligneAudit/\allowbreak A2SectorNakayama.lean}.
\begin{leancode}
theorem finrank_le_sum_fiberDim (h : GradedAction Rdeg Mdeg) (a : ℕ) :
    finrank K (Mdeg a) ≤
      ∑ b ∈ Finset.range (a + 1), fiberDim Rdeg Mdeg b * finrank K (Rdeg (a - b)) := by
\end{leancode}
\noindent File \leanpath{VOA/\allowbreak DeligneAudit/\allowbreak A2SectorTensor.lean}.
\begin{leancode}
theorem finrank_currentCarrier_le (m d r : ℕ) :
    Module.finrank K (degreePiece (currentCarrier Q m d) r) ≤
      ∑ a ∈ Finset.range (r + 1),
        Module.finrank K (degreePiece (currentCarrier Q m 0) a) * partitionsLe d (r - a) := by
\end{leancode}

\begin{dictionary}
\lean{finrank_le_sum_fiberDim} is graded Nakayama for a graded action of a graded algebra $R$ on a graded space
$M$: $\dim M_a\le\sum_b\dim(\text{fibre})_b\dim R_{a-b}$. \lean{finrank_currentCarrier_le} is the $1/(q)_d$ split:
\lean{partitionsLe d k} is the number of partitions of $k$ into at most $d$ parts, the dimension of degree-$k$
symmetric polynomials in $d$ variables; \lean{partitionsLe_series} identifies its generating function with
$1/(q)_d$. Only inequalities are needed.
\end{dictionary}

\textbf{Proof on paper.} By B6 and B8, $W_m=M_m$, so the Nakayama fibre of $W_m$ is that of $M_m$, whose degree-$b$
piece is bounded by $\dim\mathrm{cycDeg}(w,b)\le[q^b]F_m$ (U5). Graded Nakayama over $B_m$ (Hilbert series $1/(q)_m$) gives
$\mathrm{currentDim}(m,0,a)\le[q^a]F_m/(q)_m$; the auxiliary split gives the extra $1/(q)_d$. \qed

\begin{physics}
The $m$-current sector is the fusion fibre times free bosonic oscillations of the centre of mass and relative
positions ($1/(q)_m$), and the $d$ Casimir quanta add $1/(q)_d$.
\end{physics}

\begin{status}
Proved, conditional on U5 and \lean{DivisorComparisonG S}; axioms standard. The divisor comparison is then
discharged per type.
\end{status}

%======================================================================
\section{The divisor comparison}\label{l24:sec:divisor}
%======================================================================

\begin{definition}[Pair comparison]\label{l24:def:pair}
Let $E_{ab}=\prod(z_i-z_j)$ over ordered pairs $i\neq j$ with $\{i,j\}\neq\{a,b\}$. \emph{Pair comparison} holds if for
every $x\in W_m$ and $a\neq b$ some power $E_{ab}^N x$ lies in $P_mM_m$.
\end{definition}
\noindent File \leanpath{VOA/\allowbreak DeligneAudit/\allowbreak DeligneDivisorReduction.lean}.
\begin{leancode}
structure PairComparisonG : Prop where
  le : ∀ m, ∀ x ∈ currentCarrier S.fullMatrix m 0, ∀ a b : Fin (m + 0), a ≠ b →
    ∃ N : ℕ, (pairE a b) ^ N • x ∈ MP S m

theorem divisorComparisonG_of_pair (h : PairComparisonG S) : DivisorComparisonG S := by
\end{leancode}
\noindent File \leanpath{VOA/\allowbreak DeligneAudit/\allowbreak A2DivisorAvoid.lean}.
\begin{leancode}
def pairE (a b : Fin (m + 0)) : MvPolynomial (Fin (m + 0)) ℂ :=
  ∏ i, ∏ j, if i = j ∨ (i = a ∧ j = b) ∨ (i = b ∧ j = a) then 1 else (X i - X j)
\end{leancode}

\begin{dictionary}
\lean{MP S m} is the $P_m$-span of $M_m$ (\lean{Submodule.span (Pm m) (Mspace …)}). $E_{ab}$ vanishes on every
collision except $z_a=z_b$: pair comparison is B7 localized at the divisor $z_a=z_b$.
\end{dictionary}

\textbf{Proof on paper (of the reduction).} Let $x\in W_m$ and let $\mathfrak p$ be a height-one prime of $B_m$ with
$\Delta\in\mathfrak p$. We need $s\notin\mathfrak p$ with $sx\in M_m$.

\emph{Step 1 ($\mathfrak p$ lives on one collision).} $B_m\cong\C[e_1,\dots,e_m]$ is a UFD, so the height-one prime $\mathfrak p$ is
principal, $\mathfrak p=(\varpi)$ with $\varpi$ prime in $B_m$. In the larger UFD $P_m$, the non-unit $\varpi$ has an irreducible
(hence prime) factor $\phi$. Since $\varpi\mid\Delta=\prod_{i\neq j}(z_i-z_j)$, the prime $\phi$ divides some linear factor
$z_i-z_j$, so it is associate to it, and $\varpi$ vanishes identically on the hyperplane $z_i=z_j$. (No finiteness of $P_m$
over $B_m$ is needed: the argument only factors $\varpi$ inside $P_m$.)

\emph{Step 2 (an element outside $\mathfrak p$).} Let $w=\sum_{a\neq b}E_{ab}$, a symmetric polynomial. On the hyperplane
$z_i=z_j$, every $E_{ab}$ with $\{a,b\}\neq\{i,j\}$ contains the factor $z_i-z_j$ and vanishes, while
$E_{ij}=E_{ji}=\prod_{\{k,l\}\neq\{i,j\}}(z_k-z_l)$ (over ordered pairs) does not vanish identically there. So $w$ restricted to
$z_i=z_j$ is $2E_{ij}|_{z_i=z_j}\neq0$, while $\varpi$ restricts to $0$. Hence $w\notin(\varpi)=\mathfrak p$ (\lean{wS_not_mem}),
and $w^K\notin\mathfrak p$ for every $K$ since $\mathfrak p$ is prime.

\emph{Step 3 ($w^Kx\in M_m$).} Pair comparison gives, for each ordered pair $(a,b)$, a power with $E_{ab}^Nx\in P_mM_m$; take
$N$ common to all pairs. Expanding $w^K$ for $K$ large ($K>N\cdot\#\text{pairs}$), each monomial in the $E_{ab}$ contains some
$E_{ab}^N$, so $w^Kx\in P_mM_m$. Since $w$ and $x$ are symmetric, so is $w^Kx$, and Lemma~\ref{l24:lem:reynolds} puts it
in $M_m$. Take $s=w^K$. \qed

\begin{physics}
Away from a single collision $z_a=z_b$ everything is generic; the only physics is the two-body problem at that
collision, with the other particles as spectators at distinct positions.
\end{physics}

\begin{status}
Proved; axioms standard.
\end{status}

\paragraph{Strategy of the certificate.} Pair comparison asks, for each $x\in W_m$ and each pair of slots $(a,b)$, that
$x$ be generated (after multiplying by a power of $E_{ab}$, which is invertible near a generic point of $z_a=z_b$) from
the ground state $\delta_{c_0}^{\otimes m}$. Near that divisor the spectators sit at distinct points and can be moved
freely, so the problem is about the two slots $a,b$, i.e.\ about $n\times n$ tables of polynomials (elements of
$\ltwentyfourg^*\otimes\ltwentyfourg^*$ with polynomial coefficients). Three ideas organize it.
\begin{enumerate}
\item \emph{Spectators.} Split $x$ by the labels of its spectators: $x=\sum_sx_s$, where $s$ runs over labellings with
  the convention $s_a=s_b=c_0$ (\emph{canonical} spectators; the labels at $a,b$ are carried by the table) and $x_s$ is the
  embedding of the \emph{pair table} $Y_{cd}=x_{s[a\mapsto c,b\mapsto d]}$ (\lean{tensor_eq_sum_pemb}). Moving each spectator
  slot $j$ from $c_0$ to its actual label $s_j$ costs the factor $\lambda_j^L$ (Step 2 of the generic comparison); this is the
  factor $\Lambda_{ab}=\prod_{j\notin\{a,b\}}\lambda_j^L$ below.
\item \emph{Two-slot moves.} Two operators act on slots $a,b$ only and preserve $P_mM_m$: the diagonal move
  $S_x=H\,(x^{(a)}+x^{(b)})$ and the antisymmetric move $A_x=H\,(z_a-z_b)(x^{(a)}-x^{(b)})$ (Theorem~\ref{l24:thm:pairops}).
  Each costs one factor $H$, and $A_x$ also a factor $z_a-z_b$.
\item \emph{A generating tree.} Starting from the root table $\delta_{c_0}\otimes\delta_{c_0}$, apply moves $S_x$, $A_x$
  with integer structure constants. A \emph{tree} is a list of $n^2$ integer tables, each obtained from an earlier one (its
  \emph{parent}) by one move; each node records its parent, whether the move is $A$, the letter $x$, its \emph{step count}
  $\nu$ (number of moves from the root) and its \emph{collision order} $\ell$ (number of $A$-moves). Diagonal moves alone
  stay in $U(\ltwentyfourg)\cdot(\delta_{c_0}\otimes\delta_{c_0})=V(2\theta)$, the symmetric channel that is regular at contact;
  the nodes of collision order $1$ fill the antisymmetric part $\Lambda^2\ltwentyfourg^*$, and those of order $2$ the rest of the
  symmetric part, a copy of $I_2^*$. In the certificates the numbers of nodes of collision order $0,1,2$ are therefore
  $\dim V(2\theta)$, $\dim\Lambda^2\ltwentyfourg$, $\dim I_2$, adding up to $n^2$ (for $A_2$: $27+28+9=64$); this count is a
  property of the generated data, which Lean does not need: it checks only the moves and the final combinations.
\end{enumerate}
What the kernel then checks is that each node is the claimed move of its parent, and that the target tables (antisymmetric
basis tensors, a complement of $V(2\theta)$, symmetric basis tensors corrected by their $I_2$-coordinates) are integer
combinations of nodes of the right collision order. The ``good'' and ``level-OK'' vocabulary below makes this precise.

\begin{definition}[Two-slot tables, the good submodule, the level predicate]\label{l24:def:levok}
Fix slots $a\neq b$ and spectator labels $s$ with $s_a=s_b=c_0$. A \emph{pair table} $Y$ is an $n\times n$ matrix of
polynomials, embedded as $\sum_{c,d}Y_{cd}\,\delta_{s[a\mapsto c,b\mapsto d]}$. It is \emph{good} if the embedding lies
in $P_mM_m$. With $H=G_aG_b$, $G_i=\prod_{k\notin\{a,b\}}(z_i-z_k)$, $\Lambda_{ab}=\prod_{j\notin\{a,b\}}\lambda_j^L$, an integer
table $T$ is \emph{level-OK at $(\nu,\ell)$} if $H^\nu(z_a-z_b)^\ell\Lambda_{ab}\cdot T$ is good for all $m$, $a\neq b$ and canonical $s$.
The factor $\Lambda_{ab}$ pays for moving the spectators (item 1 above), $H^\nu$ for $\nu$ two-slot moves, $(z_a-z_b)^\ell$ for
$\ell$ antisymmetric moves.
\end{definition}
\noindent File \leanpath{VOA/\allowbreak DeligneAudit/\allowbreak DeligneDivisorModel.lean}.
\begin{leancode}
def Good (a b : Fin m) (s : Fin m → Fin S.n) : Submodule (Pm m) (PairTable S m) :=
  (MP S m).comap (pembL a b s)

abbrev Hpair (a b : Fin m) : Pm m := Gpair a b a * Gpair a b b

abbrev zab (a b : Fin m) : Pm m := zr (K := ℂ) a - zr b

def Fspec (a b : Fin m) : Pm m :=
  ∏ j ∈ Finset.univ.filter (fun j => j ≠ a ∧ j ≠ b), Dslot j ^ S.depth
\end{leancode}
\noindent File \leanpath{VOA/\allowbreak DeligneAudit/\allowbreak DeligneDivisorAssembly.lean}.
\begin{leancode}
def LevOK (D ℓ : ℕ) (T : ZTab S) : Prop :=
  ∀ (m : ℕ) (a b : Fin m), a ≠ b → ∀ s : Fin m → Fin S.n, IsCanon a b s →
    (Hpair a b ^ D * zab a b ^ ℓ * Fspec S a b) • castT (m := m) T ∈ Good a b s
\end{leancode}

\begin{dictionary}
\lean{ZTab S} is $\mathrm{Fin}\,n\times\mathrm{Fin}\,n\to\Z$ and \lean{castT} views it as a constant pair table.
\lean{IsCanon a b s} is $s_a=s_b=c_0$. \lean{Gpair a b i} is $\prod_{k\notin\{a,b\}}(z_i-z_k)$ (evaluation of
\lean{Lpair}). \lean{Fspec} is $\Lambda_{ab}$, with \lean{Dslot j} $=\lambda_j$. The Lean variable \lean{D} of \lean{LevOK} is our
step count $\nu$ (not the carrier data), and \lean{ℓ} the collision order. Level-OK is monotone: it is preserved when $\nu$ or $\ell$
is increased (\lean{levOK_mono}), and it is closed under integer linear combinations.
\end{dictionary}

\begin{physics}
Bookkeeping for the two-body problem: which two-particle states are produced from the ground state, by how many
excitations, and with which power of the separation $z_a-z_b$ (the collision order $\ell$ is the order of vanishing of the
two-particle wave function at contact: $0$ in the $V(2\theta)$ channel, $1$ in $\Lambda^2\ltwentyfourg$, $2$ in $I_2$).
\end{physics}

\begin{status}
Definitions.
\end{status}

\begin{theorem}[The pair operators and the W-side divisibilities]\label{l24:thm:pairops}
(i) $P_mM_m$ is stable under $S_x=G_aG_b(x^{(a)}+x^{(b)})$ and $A_x=G_aG_b(z_a-z_b)(x^{(a)}-x^{(b)})$. (ii) For
$x\in W_m$, the antisymmetric part of its $(a,b)$ pair table is divisible by $z_a-z_b$, and every relation
coordinate $\sum_{c,d}q_r(c,d)(Y_{cd}+Y_{dc})$ of $I_2$ is divisible by $(z_a-z_b)^2$.
\end{theorem}
\noindent File \leanpath{VOA/\allowbreak DeligneAudit/\allowbreak DeligneDivisorReduction.lean}.
\begin{leancode}
theorem pairS_mem (x : Fin S.n) (a b : Fin m) (hab : a ≠ b) {v : TPs S m} (hv : v ∈ MP S m) :
    (Gpair a b a * Gpair a b b) • (slotOp x a v + slotOp x b v) ∈ MP S m := by

theorem pairA_mem (x : Fin S.n) (a b : Fin m) (hab : a ≠ b) {v : TPs S m} (hv : v ∈ MP S m) :
    (Gpair a b a * Gpair a b b * (zr (K := ℂ) a - zr b)) • (slotOp x a v - slotOp x b v) ∈
      MP S m := by
\end{leancode}
\noindent File \leanpath{VOA/\allowbreak DeligneAudit/\allowbreak DeligneDivisorModel.lean}.
\begin{leancode}
theorem anti_dvd {x : TPs S m} (hx : x ∈ currentCarrier S.fullMatrix m 0) (c d : Fin S.n) :
    zab a b ∣ pairTable x a b t c d - pairTable x a b t d c := by

def relSym (r : Option S.ρ) (Y : PairTable S m) : Pm m :=
  ∑ c, ∑ d, MvPolynomial.C (S.fullMatrix r (c, d)) * (Y c d + Y d c)

theorem relSym_sq_dvd {x : TPs S m} (hx : x ∈ currentCarrier S.fullMatrix m 0)
    (r : Option S.ρ) : zab a b ^ 2 ∣ relSym r (pairTable x a b t) := by
\end{leancode}

\begin{dictionary}
\lean{slotOp x i} is the coadjoint action of $x$ in slot $i$ only. In \lean{anti_dvd}, \lean{relSym_sq_dvd} the
slots $a\neq b$ (\lean{hab}) and the spectators $t$ are section variables.
\end{dictionary}

\textbf{Proof on paper.} (i) As in Step 1 of the generic comparison, for $\phi(t)=\sum_r\phi_rt^r\in P_m[t]$ the operator
$\sum_r\phi_r(x\otimes t^r)$ acts as $\sum_i\phi(z_i)x^{(i)}$ and preserves $P_mM_m$. Take
$\phi(t)=\prod_{k\notin\{a,b\}}(t-z_k)\cdot\psi(t)$; it kills every spectator slot and gives $G_a\psi(z_a)x^{(a)}+G_b\psi(z_b)x^{(b)}$.
For $S_x$ we need $\psi(z_a)=G_b$ and $\psi(z_b)=G_a$: take $\psi(t)=G_b-\kappa\,(t-z_a)$ with the divided difference
$\kappa=(G_b-G_a)/(z_b-z_a)$, a polynomial because $G_b-G_a$ vanishes on $z_a=z_b$. For $A_x$ multiply this $\psi$ by
$2t-z_a-z_b$, which equals $z_a-z_b$ at $t=z_a$ and $-(z_a-z_b)$ at $t=z_b$. (ii) Swap symmetry: the $(d,c)$ entry is the $(c,d)$ entry with
$z_a\leftrightarrow z_b$, so the antisymmetric part vanishes on $z_a=z_b$. The current equation says the relation
coordinate vanishes on $z_a=z_b$; it is swap-invariant, so the vanishing is to second order. \qed

\begin{physics}
At a two-particle collision the symmetric $V(2\theta)$ channel is regular, the antisymmetric $\Lambda^2\ltwentyfourg$
channel has a simple zero, and the $I_2$ channel a double zero: the exclusion rule of the fusion of two adjoint
representations, $\ltwentyfourch F_2=V(2\theta)+q\Lambda^2\ltwentyfourg+q^2I_2$.
\end{physics}

\begin{status}
Proved; axioms standard.
\end{status}

\begin{theorem}[Pair comparison from integer certificates]\label{l24:thm:cert}
Let $a_r$ ($r$ running over the relations of $I_2$, the Casimir included) be integer matrices equal to the coefficient
matrices of the sector data, $\nu$ a natural number, $N_A$, $N_S$ positive integers and $U_r$ integer tables.
Suppose (a) $N_A(e_c\otimes e_d-e_d\otimes e_c)$ is level-OK at $(\nu,1)$ for $c<d$; (b) each $U_r$ is level-OK at $(\nu,2)$;
(c) $N_S(e_c\otimes e_d+e_d\otimes e_c)-\sum_r(a_r(c,d)+a_r(d,c))U_r$ is level-OK at $(\nu,0)$ for $c\le d$. Then pair
comparison holds.
\end{theorem}
\noindent File \leanpath{VOA/\allowbreak DeligneAudit/\allowbreak DeligneDivisorAssembly.lean}.
\begin{leancode}
theorem pairComparisonG_of_cert (qI : Option S.ρ → Fin S.n → Fin S.n → ℤ)
    (hq : ∀ r c d, S.fullMatrix r (c, d) = (qI r c d : ℂ)) (D NA NS : ℕ) (hNA : NA ≠ 0)
    (hNS : NS ≠ 0) (U : Option S.ρ → ZTab S)
    (hA : ∀ c d : Fin S.n, c < d → LevOK D 1 ((NA : ℤ) • antiZ c d))
    (hU : ∀ r, LevOK D 2 (U r))
    (hV : ∀ c d : Fin S.n, c ≤ d →
      LevOK D 0 ((NS : ℤ) • symZ c d - ∑ r, (qI r c d + qI r d c) • U r)) :
    PairComparisonG S := by
\end{leancode}

\begin{dictionary}
\lean{qI} is $(a_r)$ and the Lean \lean{D} is $\nu$. \lean{antiZ c d}, \lean{symZ c d} are $e_c\otimes e_d\mp e_d\otimes e_c$.
\lean{hq} says the integer tables are the actual relations. $N_A$ and $N_S$ are common denominators: the tree's tables are a
basis of $\ltwentyfourg^*\otimes\ltwentyfourg^*$ over $\Q$, the antisymmetric and symmetric basis tensors are rational combinations of
them, and multiplying by $N_A$, $N_S$ makes the combinations integral. They depend on the basis and the tree, which is why
they vary erratically between types (for instance $N_S=559104$ for $F_4$ but $96$ for $E_8$); dividing by them again is
harmless over $\C$ (\lean{good_unsmul}). (c) says $e_c\otimes e_d+e_d\otimes e_c$, corrected by its $I_2$-coordinates against the complement tables
$U_r$, lies in the level-0 span ($V(2\theta)$). No projector or dual functional appears: the $U$-part cancels because
$\sum_{c,d}Y_{cd}\psi_r(c,d)$ is the relation coordinate, which is divisible by $(z_a-z_b)^2$.
\end{dictionary}

\textbf{Proof on paper.} Fix $x\in W_m$ and slots $a\neq b$; put $K=H^\nu\Lambda_{ab}$.

\emph{Step 1 (reduce to canonical spectators).} Write $x=\sum_sx_s$ over canonical spectator labellings $s$, $x_s$ the
embedding of the pair table $Y^{(s)}_{cd}=x_{s[a\mapsto c,b\mapsto d]}$ (\lean{tensor_eq_sum_pemb}). It suffices to show
that $K\,Y^{(s)}$ is good for every $s$: then $Kx\in P_mM_m$, and since $K=H^\nu\Lambda_{ab}$ divides
$E_{ab}^{2\nu+Lm}$ (\lean{factor_dvd}: each $G_a$, $G_b$ divides $E_{ab}$, and each $\lambda_j$, $j\notin\{a,b\}$, divides $E_{ab}$), we
get $E_{ab}^{2\nu+Lm}x=(\text{cofactor})\cdot Kx\in P_mM_m$, which is pair comparison with $N=2\nu+Lm$
(\lean{pairComparisonG_of_good}). Fix $s$ and write $Y=Y^{(s)}$.

\emph{Step 2 (antisymmetric part).} By Theorem~\ref{l24:thm:pairops}(ii), $Y_{cd}-Y_{dc}=(z_a-z_b)Z_{cd}$ with $Z$
antisymmetric. Then $K\cdot\frac12\sum_{c,d}(Y_{cd}-Y_{dc})e_c\otimes e_d$ is a $P_m$-combination of the tables
$K(z_a-z_b)\,(e_c\otimes e_d-e_d\otimes e_c)$, which are good by (a) after dividing by $N_A$.

\emph{Step 3 (symmetric part).} The symmetric part $\frac12\sum_{c,d}(Y_{cd}+Y_{dc})e_c\otimes e_d$ is a $P_m$-combination of the
$e_c\otimes e_d+e_d\otimes e_c$. By (c), each of these equals, up to the factor $N_S$, a level-0 table plus
$\sum_r(a_r(c,d)+a_r(d,c))U_r$. Collecting the second terms over $(c,d)$ with the coefficients of $Y$ gives
$\sum_r\rho_r\,U_r$ with $\rho_r=\sum_{c,d}a_r(c,d)(Y_{cd}+Y_{dc})$, the $r$-th relation coordinate of $Y$; by
Theorem~\ref{l24:thm:pairops}(ii) $\rho_r=(z_a-z_b)^2\varphi_r$. So $K\rho_rU_r=\varphi_r\cdot K(z_a-z_b)^2U_r$ is good by (b), and
the level-0 part is good by (c).

\emph{Step 4.} Adding, $4N_AN_S\,K\,Y$ is good, hence $KY$ is good. \qed

\begin{physics}
The two-body fusion is checked once and for all by finite linear algebra on $\ltwentyfourg^*\otimes\ltwentyfourg^*$: every
two-particle state with the right collision behaviour is produced from the ground state by at most $\nu$ pair
excitations.
\end{physics}

\begin{status}
Proved; axioms standard (re-checked by the author, Section~\ref{l24:sec:verification}).
\end{status}

\begin{theorem}[Soundness of the sparse and encoded checks]\label{l24:thm:sparse}
(i) If a tree of sparse integer tables, rooted at $\delta_{c_0}\otimes\delta_{c_0}$, passes the Boolean node check (each
node equals $S_x$ or $A_x$ of an earlier parent, with step and collision-order counters), then every node $i$ is level-OK at
(its step count, its collision order). (ii) If the base-$2^{64}$ encoding of a sparse table with entry sum $<2^{63}$ is $0$, the table is
$0$.
\end{theorem}
\noindent File \leanpath{VOA/\allowbreak DeligneAudit/\allowbreak DeligneDivisorSparse.lean}.
\begin{leancode}
theorem levOK_tree (hn : S.n = n) (hf : ∀ x c e, S.f x c e = (fZ br x c e : ℂ))
    (hroot : castQZ (S := S) (tab 0) = baseZ) (h0 : (par 0).2.2.2.1 = 0 ∧ (par 0).2.2.2.2 = 0)
    (N : ℕ) (hnode : ∀ i, 0 < i → i < N → nodeCheck br n tab par i = true) :
    ∀ i, i < N → LevOK (par i).2.2.2.1 (par i).2.2.2.2 (castQZ (S := S) (tab i)) := by
\end{leancode}
\noindent Files \leanpath{VOA/\allowbreak DeligneAudit/\allowbreak DeligneDivisorNumeric.lean} (the definitions \lean{bpow},
\lean{keyE}, \lean{encQ}, \lean{encCheck}, moved there from the next file) and
\leanpath{VOA/\allowbreak DeligneAudit/\allowbreak DeligneDivisorEnc.lean} (\lean{castQZ_eq_of_enc}; it imports the first).
\begin{leancode}
def bpow (e : ℕ) : ℤ := ((2 ^ (64 * e) : ℕ) : ℤ)

def keyE (kb : ℕ) (t : ℕ × ℕ × ℤ) : ℕ := t.1 * kb + t.2.1

def encQ (kb : ℕ) (L : QV) : ℤ := (L.map fun t => t.2.2 * bpow (keyE kb t)).sum

def encCheck (kb : ℕ) (L : QV) : Bool :=
  (L.all fun t => decide (t.1 < kb) && decide (t.2.1 < kb)) && decide (absSum L < 2 ^ 63) &&
    decide (encQ kb L = 0)

theorem castQZ_eq_of_enc {S : DeligneSectorData} {kb : ℕ} (hn : S.n ≤ kb) {L M : QV}
    (h : encCheck kb (L ++ smulQ (-1) M) = true) : castQZ (S := S) L = castQZ M := by
\end{leancode}

\begin{dictionary}
A sparse table \lean{QV} is a list of triples $(c,d,k)$ meaning $k\,e_c\otimes e_d$ (repeats add); \lean{castQZ} reads it
as a \lean{ZTab}. \lean{par i} records (parent, is-$A$?, letter $x$, step count $\nu$, collision order $\ell$); the Lean comments call the last two
``depth'' and ``level''. This tree depth is unrelated to the cyclicity depth $L$ of the sector data. \lean{fZ br} are the integer
structure constants of a sparse bracket table \lean{br}; \lean{hf} ties them to the sector data. In (ii),
\lean{absSum} is the sum of $|k|$; the bound $<2^{63}$ guarantees every ``digit'' (total coefficient of a key) has
absolute value $<2^{63}=B/2$, so $\sum_eg_eB^e=0$ forces $g_e=0$ (\lean{digits_zero}). \lean{encLin} checks
$\mathrm{target}=\sum_j\mu_j\mathrm{tab}_j$ by comparing encodings, with per-table encodings shared by the kernel cache.
\end{dictionary}

\textbf{Proof on paper.} (i) Strong induction along the tree, using that $S_x$ and $A_x$ preserve goodness with one more
factor $H$ (and $z_a-z_b$ for $A_x$) (\lean{levOK_S}, \lean{levOK_A}). (ii) Balanced base-$B$ digits are unique. \qed

\begin{physics}
These lemmas are part of solving the two-body problem of adjoint quanta: the generating tree is a list of two-particle
states, each obtained from an earlier one by exciting one or both particles with a flavour current, and the theorem says
that if each listed step is a legitimate excitation then every listed state really is produced from the ground state with
the recorded number of excitations and order of vanishing at contact. It is what allows a 61\,504-state two-body
calculation for $E_8$ to be accepted as part of a proof.
\end{physics}

\begin{status}
Proved; axioms standard.
\end{status}

\begin{warning}[title={Superseded checkers}]
\leanpath{DeligneDivisorSort.lean} (a merge-sort zero check) and \leanpath{DeligneDivisorFast.lean} (row-based pair
operators) are sound but were superseded as too slow in the kernel; see Section~\ref{l24:sec:certs}. The column-based
operators of \leanpath{DeligneDivisorCol.lean} (\lean{levOK_treeC}) and the encoded checks are what the $E_6$, $E_7$,
$E_8$ instances use.
\end{warning}

%======================================================================
\section{The seven instances}\label{l24:sec:instances}
%======================================================================

Each type supplies carrier data (relations, Casimir, U3), sector data (structure constants, stability, $c_0$, local
Weyl data, cyclicity depth) and a divisor certificate. The table records the data as they appear in Lean (indices
are positions in the certificate basis).

\begin{center}\small
\begin{tabular}{lrrrrrrrr}
\toprule
type & $n$ & $|\rho|$ & $c_0$ & cycl.\ depth $L$ & nodes $n^2$ & by order $\ell=0/1/2$ & $N_A$ & $N_S$\\
\midrule
$A_2$ & 8 & 8 & 5 & 3 & 64 & 27/28/9 & -- & -- \\
$G_2$ & 14 & 27 & 12 & 10 & 196 & 77/91/28 & 18 & 48384\\
$D_4$ & 28 & 105 & 27 & 10 & 784 & 300/378/106 & 8 & 384\\
$F_4$ & 52 & 324 & 45 & 22 & 2704 & 1053/1326/325 & 48 & 559104\\
$E_6$ & 78 & 650 & 77 & 3 & 6084 & 2430/3003/651 & 12 & 216\\
$E_7$ & 133 & 1539 & 132 & 2 & 17689 & 7371/8778/1540 & 24 & 672\\
$E_8$ & 248 & 3875 & 247 & 2 & 61504 & 27000/30628/3876 & 120 & 96\\
\bottomrule
\end{tabular}
\end{center}
Columns: $|\rho|=\dim\ltwentyfourIns$ is the number of listed relations; $c_0$ is the index of $\delta_{c_0}$; ``cycl.\ depth $L$'' is
the cyclicity depth of the sector data (every basis vector of $\ltwentyfourg^*$ is reached from $\delta_{c_0}$ by at most $L$
coadjoint steps; it depends on the basis, which is why it ranges from $2$ to $22$); ``nodes'' is the number of tables in the
generating tree of the divisor certificate; the next column splits them by collision order (see the strategy paragraph of
Section~\ref{l24:sec:divisor}). A reader can check the arithmetic: in every row the three counts are $\dim V(2\theta)$,
$\dim\Lambda^2\ltwentyfourg=n(n-1)/2$ and $\dim I_2=|\rho|+1$, and they add up to $n^2$ (for $E_8$: $27000+30628+3876=61504$).
The step count of every tree node is at most $4$ (a separate kernel check, e.g.\ \lean{d4_dep}), so all certificates are used
with $\nu=4$. This is a different ``depth'' from $L$; that it is at most $4$ for every type is a fact about the generated
trees, checked by the kernel, not a theorem about Lie algebras. $N_A$, $N_S$ are the denominators of Theorem~\ref{l24:thm:cert}. The $A_2$
certificate predates the generic assembly and uses a projector form (\leanpath{A2DivisorPair*.lean}).

\begin{theorem}[The divisor comparison and the sector fusion bound, all seven types]\label{l24:thm:seven}
For each $X\in\{A_2,G_2,D_4,F_4,E_6,E_7,E_8\}$ the divisor comparison holds for the type's sector data, and therefore,
for every family $F$, U5 for $F$ implies the sector fusion bound \lean{SectorFusionBound Xdata F}. For $A_2$, U5 is
taken in the $A_2$-specific form \lean{SectorWeyl.FusionModuleBound F} below, and the Lean theorem takes the divisor
comparison as a second hypothesis, discharged by the proved theorem \lean{A2.divisorComparison}.
\end{theorem}
\noindent File \leanpath{VOA/\allowbreak DeligneAudit/\allowbreak A2DivisorFinal.lean}.
\begin{leancode}
theorem divisorComparison : DivisorComparison := divisorComparison_of_pair pairComparison
\end{leancode}
\noindent File \leanpath{VOA/\allowbreak DeligneAudit/\allowbreak DeligneSectorA2.lean}.
\begin{leancode}
theorem A2_sectorFusionBound_generic (F : ℕ → PowerSeries ℚ) (hU5 : FusionModuleBound F)
    (hD : A2.DivisorComparison) : SectorFusionBound A2data F :=
  sectorFusionBoundG A2sector F (fusionModuleBoundG_of F hU5) (divisorComparisonG_of hD)
\end{leancode}
\noindent The two $A_2$-specific hypotheses. File \leanpath{VOA/\allowbreak DeligneAudit/\allowbreak A2SectorWeyl.lean}.
\begin{leancode}
structure LocalWeylRel {V : Type*} [AddCommGroup V] [Module ℂ V]
    (ρ : CurrentModule (bracketTable) V) (w : V) (m : ℕ) : Prop where
  raising : ∀ (e : Fin 8), e.val < 3 → ∀ r, ρ.op e r w = 0
  cartan_pos : ∀ (h : Fin 8), 6 ≤ h.val → ∀ r, 1 ≤ r → ρ.op h r w = 0
  cartan_zero : ∀ (h : Fin 8), 6 ≤ h.val → ρ.op h 0 w = (m : ℂ) • w
  lowering_f1 : (ρ.op 3 0)^[m + 1] w = 0
  lowering_f2 : (ρ.op 4 0)^[m + 1] w = 0
  fusion : ∀ (fα : Fin 8), 3 ≤ fα.val → fα.val < 6 → ρ.op fα m w = 0

structure FusionModuleBound (F : ℕ → PowerSeries ℚ) : Prop where
  le : ∀ (m : ℕ) (V : Type) [AddCommGroup V] [Module ℂ V] (ρ : CurrentModule (bracketTable) V)
    (w : V), LocalWeylRel ρ w m → ∀ b, (Module.finrank ℂ (cycDeg ρ w b) : ℚ) ≤ coeff b (F m)
  /-- The total-dimension form of U5 (Ravinder Prop. 1): `dim W_loc(mθ)/⟨(f_α ⊗ t^m) w⟩ ≤ 8^m`,
  read on the `t`-graded pieces of any module satisfying the relations. -/
  total : ∀ (m : ℕ) (V : Type) [AddCommGroup V] [Module ℂ V] (ρ : CurrentModule (bracketTable) V)
    (w : V), LocalWeylRel ρ w m → ∀ B : ℕ,
      ∑ b ∈ Finset.range B, Module.finrank ℂ (cycDeg ρ w b) ≤ 8 ^ m
\end{leancode}
\noindent File \leanpath{VOA/\allowbreak DeligneAudit/\allowbreak A2SectorB8.lean}.
\begin{leancode}
structure DivisorComparison : Prop where
  le : ∀ m, ∀ x ∈ currentCarrier fullMatrix m 0, ∀ p : HeightOnePrime (SymAlg (K := ℂ) m),
    discS m ∈ p.asIdeal → ∃ s : SymAlg (K := ℂ) m, s ∉ p.asIdeal ∧
      s • x ∈ Mspace bracketTable m (5 : Fin 8)
\end{leancode}
\noindent File \leanpath{VOA/\allowbreak DeligneAudit/\allowbreak DeligneSectorA2.lean}.
\begin{leancode}
theorem fusionModuleBoundG_of (F : ℕ → PowerSeries ℚ) (hU5 : FusionModuleBound F) :
    FusionModuleBoundG A2sector F where

theorem divisorComparisonG_of (hD : A2.DivisorComparison) : DivisorComparisonG A2sector where
\end{leancode}
\noindent File \leanpath{VOA/\allowbreak DeligneAudit/\allowbreak G2Divisor.lean}.
\begin{leancode}
theorem g2_pairComparison : PairComparisonG G2sector :=
  pairComparisonG_of_cert (S := G2sector) qI hqD 4 18 48384 (by norm_num) (by norm_num) UD hAD hUD hVD

theorem g2_divisorComparison : DivisorComparisonG G2sector :=
  divisorComparisonG_of_pair g2_pairComparison

theorem G2_sectorFusionBound_of_U5 (F : ℕ → PowerSeries ℚ) (hU5 : FusionModuleBoundG G2sector F) :
    SectorFusionBound G2data F :=
  sectorFusionBoundG G2sector F hU5 g2_divisorComparison
\end{leancode}
\noindent File \leanpath{VOA/\allowbreak DeligneAudit/\allowbreak DeligneDivisorD4.lean}.
\begin{leancode}
theorem d4_pairComparison : PairComparisonG D4sector :=
  pairComparisonG_of_cert (S := D4sector) qI hq4 4 8 384 (by norm_num) (by norm_num) U4 hA4 hU4 hV4

theorem d4_divisorComparison : DivisorComparisonG D4sector :=
  divisorComparisonG_of_pair d4_pairComparison

theorem D4_sectorFusionBound_of_U5 (F : ℕ → PowerSeries ℚ) (hU5 : FusionModuleBoundG D4sector F) :
    SectorFusionBound D4data F :=
  D4_sectorFusionBound F hU5 d4_divisorComparison

theorem d4_carrierDim_le_demazureSum_of_U5_U6 (F : ℕ → PowerSeries ℚ)
    (hU5 : FusionModuleBoundG D4sector F)
    (hU6 : MixedFusionRecursion F (hilbSeries d4DemazureHilb)) (N : ℕ) :
    carrierDim D4data N ≤ demazureSumCoeff d4DemazureHilb N :=
  d4_carrierDim_le_demazureSum_of_sector F hU5 d4_divisorComparison hU6 N
\end{leancode}
\noindent File \leanpath{VOA/\allowbreak DeligneAudit/\allowbreak F4Divisor.lean}.
\begin{leancode}
theorem f4_pairComparison : PairComparisonG F4sector :=
  pairComparisonG_of_cert (S := F4sector) qI hqD 4 48 559104 (by norm_num) (by norm_num) UD hAD hUD hVD

theorem f4_divisorComparison : DivisorComparisonG F4sector :=
  divisorComparisonG_of_pair f4_pairComparison

theorem F4_sectorFusionBound_of_U5 (F : ℕ → PowerSeries ℚ) (hU5 : FusionModuleBoundG F4sector F) :
    SectorFusionBound F4data F :=
  sectorFusionBoundG F4sector F hU5 f4_divisorComparison
\end{leancode}
\noindent File \leanpath{VOA/\allowbreak DeligneAudit/\allowbreak E6Divisor.lean}.
\begin{leancode}
theorem e6_pairComparison : PairComparisonG E6sector :=
  pairComparisonG_of_cert (S := E6sector) qI hq6 4 12 216 (by norm_num) (by norm_num) U6 hA6 hU6 hV6

theorem e6_divisorComparison : DivisorComparisonG E6sector :=
  divisorComparisonG_of_pair e6_pairComparison

theorem E6_sectorFusionBound_of_U5 (F : ℕ → PowerSeries ℚ) (hU5 : FusionModuleBoundG E6sector F) :
    SectorFusionBound E6data F :=
  sectorFusionBoundG E6sector F hU5 e6_divisorComparison
\end{leancode}
\noindent File \leanpath{VOA/\allowbreak DeligneAudit/\allowbreak E7Divisor.lean}.
\begin{leancode}
theorem e7_pairComparison : PairComparisonG E7sector :=
  pairComparisonG_of_cert (S := E7sector) qI hq7 4 24 672 (by norm_num) (by norm_num) U7 hA7 hU7 hV7

theorem e7_divisorComparison : DivisorComparisonG E7sector :=
  divisorComparisonG_of_pair e7_pairComparison

theorem E7_sectorFusionBound_of_U5 (F : ℕ → PowerSeries ℚ) (hU5 : FusionModuleBoundG E7sector F) :
    SectorFusionBound E7data F :=
  sectorFusionBoundG E7sector F hU5 e7_divisorComparison
\end{leancode}
\noindent File \leanpath{VOA/\allowbreak DeligneAudit/\allowbreak E8Divisor.lean}.
\begin{leancode}
theorem e8_pairComparison : PairComparisonG E8sector :=
  pairComparisonG_of_cert (S := E8sector) qI hq8 4 120 96 (by norm_num) (by norm_num) U8 hA8 hU8 hV8

theorem e8_divisorComparison : DivisorComparisonG E8sector :=
  divisorComparisonG_of_pair e8_pairComparison

theorem E8_sectorFusionBound_of_U5 (F : ℕ → PowerSeries ℚ) (hU5 : FusionModuleBoundG E8sector F) :
    SectorFusionBound E8data F :=
  sectorFusionBoundG E8sector F hU5 e8_divisorComparison
\end{leancode}

\begin{dictionary}
\textbf{Full names.} \lean{VOA.DeligneAudit.DeligneDivisor.X.x_divisorComparison} for
$X\in\{\mathrm{G2},\mathrm{D4},\mathrm{F4},\mathrm{E6},\mathrm{E7},\mathrm{E8}\}$; for $A_2$,
\lean{VOA.DeligneAudit.A2.divisorComparison : A2.DivisorComparison} (the A$_2$-specific structure, same statement
for $A_2$'s data), bridged by \lean{VOA.DeligneAudit.DeligneSector.divisorComparisonG_of} (there is no declaration named
\lean{a2_divisorComparison}). The sector data are \lean{G2sector}, \lean{D4sector},
\lean{F4sector}, \lean{E6sector}, \lean{E7sector}, \lean{E8sector}, \lean{A2sector}, each with
\lean{Xsector.toDeligneCarrierData = Xdata} by \lean{rfl}. In each certificate call, $D=4$ and the two following
numerals are $N_A$ and $N_S$; \lean{qI} are the integer relation tables, \lean{hq*} their identification with the
carrier data, \lean{U*} the complement tables, \lean{hA*}, \lean{hU*}, \lean{hV*} the three certificate families,
obtained from the kernel checks by \lean{levOK_tree}/\lean{levOK_treeC} and \lean{encLin}.

\textbf{The $A_2$ hypotheses.} \lean{SectorWeyl.FusionModuleBound F} is U5 for $\mathfrak{sl}_3$ written in the
bespoke $A_2$ indexing of \lean{AffineSl3.bracketTable} (indices $0,1,2$ the raising vectors $e_1,e_2,e_{12}$; $3,4$ the
simple lowering vectors; $5=c_0$ the lowest root vector $f_\theta$; $6,7$ the Cartan elements, with $\theta(h_1)=\theta(h_2)=1$):
\lean{SectorWeyl.LocalWeylRel} is the local Weyl relations of $W_{\mathrm{loc}}(m\theta)$ plus fusion relations for the three
negative root vectors, and the two clauses \lean{le}, \lean{total} are those of \lean{FusionModuleBoundG}
(Definition~\ref{l24:def:U5}). \lean{fusionModuleBoundG_of} shows \lean{FusionModuleBound F} implies
\lean{FusionModuleBoundG A2sector F} (every module satisfying \lean{LocalWeylRelG A2sector} satisfies
\lean{LocalWeylRel}, by \lean{localWeylRel_of_G}); we do not claim the converse. \lean{A2.DivisorComparison} is the
divisor comparison of Theorem~\ref{l24:thm:B8} for $A_2$'s data (\lean{fullMatrix}, \lean{bracketTable}, $c_0=5$),
and \lean{divisorComparisonG_of} converts it to \lean{DivisorComparisonG A2sector}. It is proved, with no hypothesis, by
\lean{A2.divisorComparison} (\leanpath{A2DivisorFinal.lean}, printed above).

\textbf{Hypotheses.} The divisor comparisons are unconditional. For the six types $G_2,\dots,E_8$ the sector fusion
bounds \lean{X_sectorFusionBound_of_U5} assume exactly U5 (\lean{FusionModuleBoundG Xsector F}) for $F$. For $A_2$ the
printed theorem \lean{A2_sectorFusionBound_generic} has two hypotheses, \lean{hU5 : FusionModuleBound F} ($A_2$ form of
U5) and \lean{hD : A2.DivisorComparison}; the bound from U5 alone is the term
\lean{A2_sectorFusionBound_generic F hU5 A2.divisorComparison}. The library has no declaration of that statement under
its own name (there is no \lean{A2_sectorFusionBound_of_U5}).
The D$_4$ corollary adds U6 and gives the carrier side of the chain against the Demazure sum.
\end{dictionary}

\textbf{Proof on paper.} Theorems~\ref{l24:thm:cert}, \ref{l24:thm:sparse} with the type's certificates, then the
reduction of Definition~\ref{l24:def:pair} and Theorem~\ref{l24:thm:sfb}. \qed

\begin{physics}
For each flavour symmetry in the series, the two-body problem of adjoint quanta is solved exactly; together with
the fusion bound U5 this bounds the number of operators of the chiral algebra in every sector by the fusion-product
count.
\end{physics}

\begin{status}
Proved; every name above is \texttt{\#print axioms}-standard. \lean{x_divisorComparison}: \textbf{unconditional}.
\lean{X_sectorFusionBound_of_U5} ($G_2,\dots,E_8$): \texttt{\#print axioms}-clean, \textbf{conditional on U5}
(\lean{FusionModuleBoundG}). $A_2$: \lean{A2_sectorFusionBound_generic} is \texttt{\#print axioms}-clean,
\textbf{conditional on the $A_2$ form of U5 and on} \lean{A2.DivisorComparison}, which is discharged by the proved
\lean{A2.divisorComparison} (see the dictionary: the resulting ``U5 alone'' statement has no name of its own). Build
evidence (campaign record): per-type builds of the divisor modules green ($E_8$ 9919 jobs, $E_7$ 9317, $E_6$ 8927, $F_4$ 9002,
$G_2$ 8835, $D_4$ 8795), all within the earlier umbrella build (the $A_2$ divisor certificate was later split; its record is the
targeted build of Section~\ref{l24:sec:verification}). The author re-ran \texttt{\#print axioms} for \lean{d4_divisorComparison} and
\lean{D4_sectorFusionBound_of_U5}.
\end{status}

%======================================================================
\section{The certificate library, described collectively}\label{l24:sec:certs}
%======================================================================

About 1020 of the chapter's 1195 files are certificate data or certificate checks. None contains mathematics beyond
the soundness lemmas of the previous sections; each family exists to discharge one finite hypothesis.

\subsection{Trust model: why a file of integers is a proof}\label{l24:sec:trust}

A reader who has never used Lean will ask why 10.8\,MB of integers should count as a proof, and what must be trusted.
The answer has four parts.
\begin{enumerate}
\item \textbf{The generators are untrusted.} The certificates were produced by Python scripts (exact rational
  arithmetic, then integers). Nothing about those scripts is assumed: a wrong script produces a certificate that fails to
  check, not a false theorem.
\item \textbf{The predicates are defined in Lean and proved sound once.} Each family is checked by a Boolean function
  written in Lean (for instance \lean{nodeCheck}, \lean{encCheck}, \lean{symP}), and a generic theorem says what a
  \lean{true} value means mathematically (Theorems~\ref{l24:thm:cert} and \ref{l24:thm:sparse}: ``if every node passes the
  node check, every node is level-OK''). These soundness theorems are ordinary proofs, valid for all data.
\item \textbf{The kernel evaluates the predicates on the literal data, and that evaluation is a proof.} A statement
  \lean{rangeAll pred lo hi = true} is proved by \lean{decide +kernel}: Lean's kernel, the small trusted program that checks
  every proof, reduces the closed Boolean term to \lean{true} by computation, and the equation then holds by definition
  (\lean{rfl}). This is a proof in the same sense as any other: the kernel is the only arbiter. It differs from
  \lean{native_decide}, which runs compiled code and adds the compiler to the trusted base; \lean{native_decide} is
  forbidden in the library and used nowhere here. Likewise a theorem of the form \lean{a = b := rfl} (the ``bridges'' below)
  is a proof that both sides compute to the same term, checked by the kernel.
\item \textbf{What a referee must audit.} Only the soundness theorems (Sections~\ref{l24:sec:divisor}--\ref{l24:sec:instances})
  and the identification of the literal tables with the Lie data, i.e.\ the hypotheses \lean{hq*} (the integer relation
  tables are the relations of the sector data) and \lean{hf} (the integer structure constants are those of the sector
  data), which are themselves kernel-checked equalities. The numbers in the data files need no inspection.
\end{enumerate}
The three axioms that \texttt{\#print axioms} reports for these theorems, \lean{propext}, \lean{Classical.choice} and
\lean{Quot.sound}, are the standard axioms of Lean's logic (Chapter~\ref{ch:01}, Section~\ref{ch01:sec:trust}); no project
axiom enters.

\subsection{Families}

\begin{center}\small
\begin{longtable}{p{0.27\textwidth}p{0.36\textwidth}p{0.29\textwidth}}
\toprule
family (files) & what the kernel checks & size\\
\midrule\endhead
U3 certificates: \leanpath{DeligneCarrierD4}, \leanpath{E6CarrierU*}, \leanpath{E7CarrierU*}, \leanpath{E8CarrierU*},
F$_4$/G$_2$ in \leanpath{*SectorTables} & $L_a\,X_aQ=\sum c\,X_x P_b$ in $\operatorname{Sym}^3$ (a cubic identity per current) &
$D_4$: 28 \lean{ring} proofs (18--54 terms); $E_7$: 23 files; $E_8$: 62 files, 248 certificates (mean 590, max 993 terms;
expansions up to 14\,646 terms)\\
Presentation steps: \leanpath{D4LiePresentationSteps}, \leanpath{E6PresChecks}, \leanpath{E7LieSteps*}, \leanpath{F4LieSteps*},
\leanpath{G2LieSteps} ($E_8$: the vacuum-step checks of Chapter~\ref{ch:L25}) & $q\operatorname{ad}_bP_j=pP_k$, sparse quadratic identities & 102 / 649 / 1538 (8 files of 193) / 323 / 26 steps\\
Singular-vector checks: \leanpath{E6SingChecks}, \leanpath{F4LieSing*}, \leanpath{G2LieSing*}, \leanpath{E8LieSingular*}, in-file for $D_4$, $E_7$ &
$J^c(1)$, $J^c(2)$ of the highest-weight vector vanish (integer sums) & small; $F_4$ split into 9 files\\
Sector: \leanpath{DeligneSectorD4*}, \leanpath{E6Sector*}, \leanpath{E7Sector*}, \leanpath{E8Sector*}, \leanpath{F4Sector*},
\leanpath{G2Sector*} & Jacobi (or inherited from the lattice Lie algebra), ad-stability of $I_2$ under generators,
weights, raising/Cartan tables, coadjoint cyclicity words, links between literal tables and the Lie-data tables &
$D_4$ 1380 stability certificates; $E_6$ 651, $E_7$ 1540, $E_8$ 3876 per generator; $E_8$ 37 files\\
Divisor: \leanpath{DeligneDivisorD4*}, \leanpath{G2Divisor*}, \leanpath{F4Divisor*}, \leanpath{E6Divisor*},
\leanpath{E7Divisor*}, \leanpath{E8Divisor*} & tree nodes, column tables, depth $\le4$, $\psi$ completeness, complement
tables $U$, antisymmetric and symmetric certificates & files: 3 / 12 / 83 / 7 / 77 / 619\\
\bottomrule
\end{longtable}
\end{center}

\subsection{Format and checking mechanism}

\begin{itemize}
\item \textbf{Generation.} Python scripts compute certificates in
  exact rational arithmetic ($E_8$: mod-$(2^{61}-1)$ pivots per weight block, then exact verification over $\Z$) and
  emit Lean. The emitters reproduce the committed files byte for byte (campaign record). Constants
  ($N_A$, $N_S$, the $U$ scale) are computed per type, not copied: the $E_6$ values failed for $E_7$.
\item \textbf{Data.} Literal \lean{List (List ℕ)} blocks, integers zig-zag encoded, stored two-level (rows of $n$) so
  that \lean{getD} lookups stay shallow; each data file $\le200$\,KB of source, \lean{set_option maxRecDepth 100000}.
  Data files import only \lean{Mathlib.Data.Nat.Notation}.
\item \textbf{Predicates.} Per-index Boolean predicates (\lean{nodeP}, \lean{colP}, \lean{depthP}, \lean{psiP},
  \lean{UCheck}, \lean{antiP}, \lean{symP}), with \lean{rangeAll f lo hi} checking an index range.
\item \textbf{Kernel checks.} Each chunk theorem is \lean{rangeAll pred lo hi = true} proved by \lean{decide +kernel}, sized
  by a calibrated cost model to about 1\,s; chunks are glued by \lean{rangeAll_split} (a balanced tree for $E_8$) and
  read back by \lean{rangeAll_sound}. No \lean{native_decide} anywhere.
\item \textbf{Light imports bridged by \lean{rfl}.} The $E_8$ chunk files import a file \leanpath{E8DivisorLight.lean}
  containing verbatim copies of the Boolean checkers; \leanpath{E8Divisor.lean} proves each copy equal to its original
  by \lean{rfl} before applying the generic soundness theorems.
\end{itemize}
A typical chunk and the bridges:
\noindent File \leanpath{VOA/\allowbreak DeligneAudit/\allowbreak E8DivisorChk450.lean}.
\begin{leancode}
theorem e8_sym_8144 : rangeAll symP 8144 8147 = true := by decide +kernel

theorem e8_sym_8147 : rangeAll symP 8147 8156 = true := by decide +kernel
\end{leancode}
\noindent File \leanpath{VOA/\allowbreak DeligneAudit/\allowbreak E8DivisorDefs.lean}.
\begin{leancode}
def rangeAll (f : ℕ → Bool) (lo hi : ℕ) : Bool := (List.range' lo (hi - lo)).all f

def symP (p : ℕ) : Bool := !decide (p / 248 ≤ p % 248) || symCheck (p / 248) (p % 248)
\end{leancode}
\noindent File \leanpath{VOA/\allowbreak DeligneAudit/\allowbreak E8Divisor.lean}.
\begin{leancode}
theorem nodeCheckCL_eq : @nodeCheckCL = @nodeCheckC := rfl
theorem colCheckL_eq : @colCheckL = @colCheck := rfl
theorem encLinL_eq : @encLinL = @encLin := rfl
theorem smulQL_eq : @smulQL = @smulQ := rfl
theorem qgetL_eq : @qgetL = @qget := rfl
\end{leancode}

\begin{dictionary}
The pair index $p=248c+d$ encodes $(c,d)$; \lean{symP p} is vacuously true unless $c\le d$. As explained in
Section~\ref{l24:sec:trust}, the data need no trust. The \lean{rfl} bridges are checked by the kernel too, so the light copies
cannot differ from the originals.
\end{dictionary}

\subsection{Sizes and measured costs}

The table collects the size and measured peak memory of each certificate family. The campaign worked under a per-file
memory rule of about 7\,GB, tightened during the campaign; the last column says whether the measured files respect it.
\begin{center}\small
\begin{tabular}{p{0.18\textwidth}p{0.36\textwidth}p{0.20\textwidth}p{0.14\textwidth}}
\toprule
family & files & measured peak, time & rule respected\\
\midrule
$E_8$ divisor & 619 (556 chunk files, 9291 theorems; 59 data files, 10.8\,MB) & 3.5--7.2\,GB, 11--20\,s (13 sampled chunks) & yes\\
$E_7$ divisor & 77 (643 chunk theorems in 54 files) & 7.0--8.0\,GB, 23--53\,s & no (earlier)\\
$E_6$ divisor & 7 & 8.8--10.9\,GB, 45--110\,s & no (earlier)\\
$F_4$ divisor & 83 & $\le6.9$\,GB & yes\\
$G_2$ divisor & 12 & $\le6.4$\,GB, $\le8$\,s & yes\\
$D_4$ divisor & 3 (checks file: 33 theorems) & 10.9\,GB, 172\,s & no (earlier)\\
$E_8$ U3 (carrier) & 62 & 6.9\,GB max; 6254\,s CPU in total & yes\\
$E_8$ sector & 37 & 5.1--6.6\,GB, up to 94\,s & yes\\
\bottomrule
\end{tabular}
\end{center}
The figures come from the campaign's logs, measured with \texttt{/usr/bin/time -l lake env lean -DmaxSynthPendingDepth=3};
not every file was measured. They are evidence of feasibility, not part of any proof.

\begin{warning}[title={Dead ends in kernel checking (do not retry)}]
Only the lessons of architectural interest are kept here.
\begin{itemize}
\item \textbf{One small declaration per check.} The kernel's cache grows within a declaration and is released only between
  declarations, and memory is also retained across theorems of one file. Whole-family declarations reached 18--38\,GB
  (an $E_7$ antisymmetric family, a depth check over all 17\,689 $E_7$ nodes, a whole-range $E_8$ sector check). The fix
  is many theorems of a few thousand entries each, and per-file as well as per-theorem budgets.
\item \textbf{Literal lists, not vector literals.} Kernel lookups in \lean{![…]} tables are slow at high indices; literal
  \lean{List} tables read with \lean{getD} are about ten times faster.
\item \textbf{Light imports.} Importing the sector and VOA stack costs about 6\,GB before any work. Check files import
  verbatim copies of the checkers instead, bridged to the originals by \lean{rfl} (floor about 2.2\,GB).
\item \textbf{No per-index re-walking, no kernel sorting.} Statements quantified over an index that re-walk a table entry
  for each index, kernel merge sorts, and dense zero checks are all too slow. What works: big-integer encodings of sparse
  tables (Theorem~\ref{l24:thm:sparse}(ii)) whose per-table encodings are shared by the kernel cache, column tables, and
  bucketed checks of long identities.
\item \textbf{Mathematical choices.} A kernel \lean{decide} of the Jacobi identity for the $E_8$ bracket table (about
  $1.5\cdot10^7$ triples) is infeasible without \lean{native_decide}; for $E_7$ and $E_8$, Jacobi is inherited from the
  lattice Lie algebra instead. Solving U3 current by current for $E_7$ did not finish; one solve at the highest weight plus
  propagation by lowering derivations works. A breadth-first search of all coadjoint words is hopeless for $E_8$
  ($248^4$ words at depth 4); extending only unit-vector states finds every basis vector at depth $\le2$.
\end{itemize}
\end{warning}

%======================================================================
\section{Verification record}\label{l24:sec:verification}
%======================================================================

\textbf{What the machine confirmed (campaign records).} The last complete umbrella build is at an earlier commit:
\begin{plaincode}
bash scripts/campaign_build.sh VOA.Deligne     # earlier: green, 11099 jobs, no sorryAx
#print axioms (all campaign files)             # subset of [propext, Classical.choice, Quot.sound]
lake build VOA (aggregate)                      # last green on 26 Oct (13868 jobs)
bash scripts/campaign_build.sh VOA              # 28 Oct (pre-release): green, 18901 jobs,
                                                #   then 18907 jobs after a merge
\end{plaincode}
No \lean{sorry}, \lean{admit}, \lean{native_decide} or \lean{axiom} in the campaign files (campaign greps). The
Lean sources changed after that build: between it and the state this chapter was written against, the diff of \texttt{VOA} touches 1390 files. The
changes split certificate files to fit the per-file memory budget (the $A_2$ divisor certificate
\lean{A2DivisorPairCert} into tree identities, matrix-pair checks and fixed-row bridges; the $D_4$ divisor checks;
the $E_7$ Lie steps), move declarations into lighter files (for example \lean{SingP} into
\leanpath{VOA/DeligneAudit/D4LieSingularGeneric.lean}), and add the new $A_2$ files of Chapters~\ref{ch:L26}
and~\ref{ch:L27}. According to the project's internal build record of 27 October 2026 (not re-run for this chapter), the exported statements of the split
files are unchanged and the repaired targets it lists were built green in targeted builds with standard axioms (some
splits were still in progress when it was written); the targeted
build of \lean{VOA.DeligneAudit.A2Unconditional} (9755 jobs), whose import closure contains the $A_2$ certificates of
this chapter, is green. On 2026-10-28 the aggregate build \texttt{bash scripts/campaign\_build.sh VOA} (the full \texttt{lake build VOA} under the project's memory-capped wrapper) was green on 2026-10-28 (a pre-release build): 18{,}901 jobs, and again 18{,}907 jobs after merging the then-current \texttt{origin/main}; the aggregate target contains the umbrella \lean{VOA.Deligne}. \texttt{scripts/check\_repository.py} passed (10{,}431 library files, all reachable, 0 project axioms) and \texttt{scripts/audit/run\_axiom\_audits.py} passed (5{,}958 reports across 24 probes). ``Axioms standard'' in the status boxes remains the campaign record of the \texttt{\#print axioms} sections and targeted
builds, except where a claim is marked re-checked.

\textbf{What this chapter's author checked.} (1) Every \texttt{leancode} block occurs verbatim in the sources of
the state this chapter was written against (\texttt{sync.py check}; one quotation was re-pointed to the file its declarations moved to). (2) On a
scratch file importing the single module \lean{VOA.DeligneAudit.DeligneDivisorD4}, with the checkout's oleans, at an
earlier commit (not re-run for this edition):
\begin{plaincode}
lake env lean -DmaxSynthPendingDepth=3 AxCheck.lean     # 3.98 s, 5.6 GB peak RSS
'VOA.DeligneAudit.DeligneDivisor.D4.d4_divisorComparison' depends on axioms: [propext, Classical.choice, Quot.sound]
'VOA.DeligneAudit.DeligneDivisor.D4.D4_sectorFusionBound_of_U5' depends on axioms: [propext, Classical.choice, Quot.sound]
'VOA.DeligneAudit.DeligneDivisor.pairComparisonG_of_cert' depends on axioms: [propext, Classical.choice, Quot.sound]
'VOA.DeligneAudit.DeligneSector.sectorFusionBoundG' depends on axioms: [propext, Classical.choice, Quot.sound]
'VOA.DeligneAudit.DeligneCarrier.casimirLayerBound' depends on axioms: [propext, Classical.choice, Quot.sound]
'VOA.ker_eq_of_liPresentationClassicallyFree' depends on axioms: [propext, Classical.choice, Quot.sound]
'VOA.DeligneHigherRankConditionalCapstone.caseClassicallyFree_of_explicit_inputs' depends on axioms: [propext, Classical.choice, Quot.sound]
\end{plaincode}
A second scratch file with the same single import, run during the revision of this chapter (4.6\,s):
\begin{plaincode}
'VOA.DeligneAudit.DeligneSector.mem_Mspace_of_symmetric' depends on axioms: [propext, Classical.choice, Quot.sound]
'VOA.DeligneAudit.DeligneSector.genericComparisonG' depends on axioms: [propext, Classical.choice, Quot.sound]
'VOA.DeligneAudit.DeligneSector.mFreeness' depends on axioms: [propext, Classical.choice, Quot.sound]
'VOA.DeligneAudit.DeligneDivisor.divisorComparisonG_of_pair' depends on axioms: [propext, Classical.choice, Quot.sound]
'VOA.DeligneAudit.DeligneDivisor.pairComparisonG_of_cert' depends on axioms: [propext, Classical.choice, Quot.sound]
\end{plaincode}
(3) No \texttt{\#print axioms} was re-run at the state this chapter was written against: the checkout has no compiled objects for the new
$A_2$ modules, and building them is outside the scope of a chapter revision. The author did not run \texttt{lake build}.

%======================================================================
\section{Findings, oddities and Lean defects}\label{l24:sec:defects}
%======================================================================

No statement-level defect was found: the headline statements say what the mathematics says, about genuine simple
affine VOAs, with every hypothesis visible. The documentation defects and scope remarks are collected here, each once;
the places in the chapter where they matter point back to this section.

\begin{warning}[title={Documentation defects (Lean docstrings and project notes)}]
Severity S4 means that documentation (a docstring or a note) is wrong or stale; no statement is affected.
\begin{enumerate}
\item \textbf{(Docstring, S4)} \leanpath{DeligneHigherRankConditionalCapstone.lean}: the module docstring says ``There is no
  universal affine VOA for a general simple Lie algebra yet, so this file does not pretend that the remaining higher-rank
  affine VOAs have been instantiated''. This is false since \leanpath{VOA/AffineGeneral} (Chapter~\ref{ch:L23}); all seven
  cases are instantiated in Section~\ref{l24:sec:presentation}. The docstring of
  \lean{caseClassicallyFree_of_explicit_inputs} speaks of ``the exact current $C_2$ presentation'' as a hypothesis, while only
  containment is assumed.
\item \textbf{(Docstring, S4)} \leanpath{DeligneCarrierChain.lean}: the docstring of \lean{CasimirAuxiliaryBound} says ``Not a
  literature input: a formalization target'' and refers to $\mathrm{Jet}(\mathrm{Sym}\,\mathfrak{sl}_3/I_2^{ns})$ in a type-generic
  file; it is proved (\lean{casimirAuxiliaryBound}, Theorem~\ref{l24:thm:B1}). The docstring of \lean{SectorFusionBound}
  says it uses ``U4 (Chari--Fourier--Khandai)''; the Lean derivation (\lean{sectorFusionBoundG}, Theorem~\ref{l24:thm:sfb})
  uses U5 only.
\item \textbf{(Project notes, S4)} The project's internal records (entry of 27 October 2026) name \lean{a2_divisorComparison}; no such declaration
  exists (the $A_2$ theorem is \lean{VOA.DeligneAudit.A2.divisorComparison}, with the bridge
  \lean{DeligneSector.divisorComparisonG_of}).
\end{enumerate}
The statements themselves are correct in all three cases.
\end{warning}

Further remarks for the referee (not defects):
\begin{enumerate}
\item \textbf{(Scope, S4)} $D_4$, $E_7$, $E_8$ results are about \lean{simpleVOA} of the certificate-basis Lie data; their
  identification with the \lean{Module.finBasis} constructions (and the three $\mathfrak e_6$ models) is not proved (Warning
  ``Which VOA?'' in Section~\ref{l24:sec:presentation}).
\item \textbf{(Formulation)} \lean{FusionModuleBoundG} uses \lean{Module.finrank}, which is $0$ on infinite-dimensional
  spaces; this weakens the named input harmlessly. The fusion set is free data; the two readings in use (Ravinder's set, all
  negative root vectors) impose the same relations for every type, short roots of $G_2$ and $F_4$ included: the extra
  vectors have $\langle\alpha,\theta^\vee\rangle=0$ and are killed by the bracket step of Lemma~\ref{l27:lem:truncweyl}
  (dictionary of Definition~\ref{l24:def:U5}).
\item \textbf{(Formulation)} The relation lists are not proved in Lean to span exactly $\ltwentyfourIns$ (dictionary of
  Definition~\ref{l24:def:A2rel}).
\item \textbf{(Naming)} Short names \lean{I2ns_le_ker}, \lean{explicitPresentation}, \lean{mFreeness},
  \lean{divisorComparison}, \lean{geometricCarrierEmbedding} denote different theorems in different namespaces; use the
  full names (dictionary of Theorem~\ref{l24:thm:sevenpres}).
\item \textbf{(Design history)} An earlier version of \lean{CurrentC2Presentation} had a field asserting that the ideal equals
  the kernel; it was replaced by the containment field \lean{relations} and the theorem
  \lean{exact_kernel_of_caseClassicallyFree}, so exactness is now a conclusion.
\end{enumerate}

%======================================================================
\section{Referee's checklist}\label{l24:sec:checklist}
%======================================================================

\begingroup\small
\begin{longtable}{p{0.30\textwidth}p{0.30\textwidth}p{0.22\textwidth}p{0.10\textwidth}}
\toprule
Lean declaration & Intended meaning & Hypotheses / status & Fidelity check\\
\midrule\endhead
\lean{LiPresentationClassicallyFree}, \lean{liInfinitePresentedJetMap} & classical freeness of a presentation & definitions & Def.~\ref{l24:def:lpcf}\\
\lean{HigherRankDeligneCase}, \lean{level} & the seven cases, $k=-1-h^\vee/6$ & definitions & Def.~\ref{l24:def:case}\\
\lean{CurrentC2Presentation}, \lean{CurrentJetSurjectivity}, \lean{GeometricCarrierEmbedding}, \lean{OrdinaryCarrierCharacter}, \lean{CaseClassicallyFree} & the four packages & definitions; containment only & Def.~\ref{l24:def:packages}\\
\lean{classicallyFree_}\allowbreak\lean{of_}\allowbreak\lean{carrier_}\allowbreak\lean{weightCharacter} & weightwise squeeze & packages; proved & Thm.~\ref{l24:thm:squeeze}\\
\lean{VOA.}\allowbreak\lean{ker_}\allowbreak\lean{eq_}\allowbreak\lean{of_}\allowbreak\lean{liPresentationClassicallyFree} & exactness is a conclusion & unconditional & Thm.~\ref{l24:thm:kereq}\\
\lean{caseClassicallyFree_}\allowbreak\lean{of_}\allowbreak\lean{explicit_}\allowbreak\lean{inputs} & capstone & three packages; reduction & Thm.~\ref{l24:thm:capstone}\\
\lean{D4Lie.}\allowbreak\lean{mem_}\allowbreak\lean{maxIdeal_}\allowbreak\lean{of_}\allowbreak\lean{SingP} & singular $\Rightarrow$ in max ideal & any Lie data; unconditional & Thm.~\ref{l24:thm:singp}\\
\lean{A2.}\allowbreak\lean{simpleProjection_}\allowbreak\lean{singVec} & Perše's vector dies & unconditional & Thm.~\ref{l24:thm:A2sing}\\
\lean{A2.}\allowbreak\lean{clsG_}\allowbreak\lean{adPolyT} & kernel is ad-stable & table hypothesis; proved & Thm.~\ref{l24:thm:adprop}\\
\lean{I2ns_le_ker*}, \lean{explicitPresentation*} (seven types) & $\ltwentyfourIns\subseteq\ker$ & \textbf{unconditional} & Thm.~\ref{l24:thm:sevenpres}\\
\lean{currentJetSurjectivity_}\allowbreak\lean{of_}\allowbreak\lean{genWordSpan}, \lean{Deligne.}\allowbreak\lean{genWordsSpan} & currents generate & unconditional & Thm.~\ref{l24:thm:jetsurj}\\
\lean{currentCarrier}, \lean{auxiliaryCarrier}, \lean{degreePiece}, \lean{conductor} & $W_m$, auxiliary carrier & definitions & Def.~\ref{l24:def:carrier}\\
\lean{restrictedDual}\allowbreak\lean{EquivCurrentCarrier} & B4: dual $\cong$ carrier & unconditional & Thm.~\ref{l24:thm:B4}\\
\lean{geometricCarrierEmbedding_}\allowbreak\lean{of_}\allowbreak\lean{quadraticIdeal} & package 3 for quadrics & unconditional (weak alone) & Thm.~\ref{l24:thm:geom}\\
\lean{caseClassicallyFree_}\allowbreak\lean{of_}\allowbreak\lean{quadraticCarrierUpperBound} & one-inequality capstone & \lean{QuadraticCarrierUpperBound} & Thm.~\ref{l24:thm:upper}\\
\lean{DeligneCarrierData} & type data with U3 & U3 field & Def.~\ref{l24:def:carrierdata}\\
\lean{casimirLayerBound}, \lean{casimirAuxiliaryBound} & B1 & \textbf{unconditional} & Thm.~\ref{l24:thm:B1}\\
\lean{sectorCoeff_}\allowbreak\lean{eq_}\allowbreak\lean{of_}\allowbreak\lean{mixed} & D3 resummation & U6 form; proved & Thm.~\ref{l24:thm:D3}\\
\lean{carrierDim_}\allowbreak\lean{le_}\allowbreak\lean{demazureSum} & carrier $\le$ Demazure sum & sector bound, U6 & Thm.~\ref{l24:thm:carrierside}\\
\lean{DeligneSectorData}, \lean{LocalWeylRelG}, \lean{FusionModuleBoundG} & sector data; U5 & definitions; U5 named input & Defs.~\ref{l24:def:sectordata}, \ref{l24:def:U5}\\
\lean{Mspace_le} & B5 & sector data & Thm.~\ref{l24:thm:B5}\\
\lean{mem_}\allowbreak\lean{Mspace_}\allowbreak\lean{of_}\allowbreak\lean{symmetric} & symmetrization & sector data & Lem.~\ref{l24:lem:reynolds}\\
\lean{DeligneSector.}\allowbreak\lean{mFreeness} & B6 & U5 & Thm.~\ref{l24:thm:B6}\\
\lean{genericComparisonG} & $\Delta^LW_m\subseteq M_m$ & sector data & Thm.~\ref{l24:thm:generic}\\
\lean{carrierGenerationG} & B8 & B6, divisor comparison & Thm.~\ref{l24:thm:B8}\\
\lean{sectorFusionBoundG} & sector fusion bound & U5, divisor comparison & Thm.~\ref{l24:thm:sfb}\\
\lean{divisorComparisonG_}\allowbreak\lean{of_}\allowbreak\lean{pair}, \lean{pairComparisonG_}\allowbreak\lean{of_}\allowbreak\lean{cert} & B7 at the divisor, from certificates & certificate hypotheses & Thms.~\ref{l24:thm:cert}, Def.~\ref{l24:def:pair}\\
\lean{levOK_tree}, \lean{castQZ_}\allowbreak\lean{eq_}\allowbreak\lean{of_}\allowbreak\lean{enc} & checker soundness & unconditional & Thm.~\ref{l24:thm:sparse}\\
\lean{x_divisorComparison} (G2, D4, F4, E6, E7, E8), \lean{A2.divisorComparison} & divisor comparison per type & \textbf{unconditional} & Thm.~\ref{l24:thm:seven}\\
\lean{X_sectorFusionBound_}\allowbreak\lean{of_}\allowbreak\lean{U5} (G2, D4, F4, E6, E7, E8) & sector bound per type & U5 & Thm.~\ref{l24:thm:seven}\\
\lean{A2_}\allowbreak\lean{sectorFusionBound_}\allowbreak\lean{generic} & $A_2$ sector bound & $A_2$-form U5 (\lean{SectorWeyl.}\allowbreak\lean{FusionModuleBound}) plus \lean{hD}, discharged by the proved \lean{A2.}\allowbreak\lean{divisorComparison} & Thm.~\ref{l24:thm:seven}\\
\lean{d4_}\allowbreak\lean{carrierDim_}\allowbreak\lean{le_}\allowbreak\lean{demazureSum_}\allowbreak\lean{of_}\allowbreak\lean{U5_}\allowbreak\lean{U6} & $D_4$ carrier $\le$ Demazure sum & U5, U6 & Thm.~\ref{l24:thm:seven}\\
\bottomrule
\end{longtable}
\endgroup

%======================================================================
\section{Notes and references}\label{l24:sec:notes}
%======================================================================

\textbf{The Deligne series and the associated varieties.} P.~Deligne, ``La s\'erie exceptionnelle de groupes de Lie'',
\emph{C.\ R.\ Acad.\ Sci.\ Paris} 322 (1996), 321--326; P.~Deligne and R.~de~Man, ``La s\'erie exceptionnelle de groupes de
Lie II'', \emph{C.\ R.\ Acad.\ Sci.\ Paris} 323 (1996), 577--582; T.~Arakawa and A.~Moreau, ``Joseph ideals and lisse minimal
$W$-algebras'', \emph{J.\ Inst.\ Math.\ Jussieu} 17 (2018), 397--417, arXiv:1506.00710 (weight-two singular vectors,
quasi-lisse); A.~Joseph, ``The minimal orbit in a simple Lie algebra and its associated maximal ideal'', \emph{Ann.\ Sci.\
\'Ecole Norm.\ Sup.} 9 (1976), 1--29; D.~Garfinkle, \emph{A new construction of the Joseph ideal}, PhD thesis, MIT (1982)
(the generators of the Joseph ideal, whose symbols are $I_2$); O.~Per\v{s}e, ``Vertex operator algebras associated to certain admissible representations of affine Lie algebras of
type $A$'', \emph{Glas.\ Mat.} 43 (2008), 41--57, arXiv:0707.4129 (the $A_2$ vector); B.~Kostant and W.~T.~Gan--G.~Savin on the
Joseph ideal (Gan--Savin, ``On minimal representations'', \emph{Represent.\ Theory} 9 (2005), 46--93); T.~Arakawa and
K.~Kawasetsu, ``Quasi-lisse vertex algebras and modular linear differential equations'', in \emph{Lie Groups, Geometry, and
Representation Theory}, Progr.\ Math.\ 326 (2018), arXiv:1610.05865.

\textbf{Li filtration and classical freeness.} H.~Li, ``Abelianizing vertex algebras'', \emph{Comm.\ Math.\ Phys.} 259 (2005),
391--411; T.~Arakawa, ``A remark on the $C_2$-cofiniteness condition on vertex algebras'', \emph{Math.\ Z.} 270 (2012),
559--575; J.~van~Ekeren and R.~Heluani, ``Chiral homology of elliptic curves and the Zhu algebra'', \emph{Comm.\ Math.\ Phys.}
386 (2021), 495--550.

\textbf{Fusion, Weyl and Demazure modules.} V.~Chari and A.~Pressley, ``Weyl modules for classical and quantum affine
algebras'', \emph{Represent.\ Theory} 5 (2001), 191--223; V.~Chari, G.~Fourier and T.~Khandai, ``A categorical approach to Weyl
modules'', \emph{Transform.\ Groups} 15 (2010), 517--549, arXiv:0906.2014; E.~Feigin, ``The PBW filtration, Demazure modules and
toroidal current algebras'', \emph{SIGMA} 4 (2008), 070, arXiv:0806.4851, and ``Infinite fusion products and
$\widehat{\mathfrak{sl}}_2$ cosets''; B.~Ravinder, ``Demazure modules, Chari--Venkatesh modules and fusion products'', \emph{SIGMA} 10 (2014), 110, arXiv:1409.0274; G.~Fourier and
P.~Littelmann, ``Weyl modules, Demazure modules, KR-modules, crystals, fusion products and limit constructions'',
\emph{Adv.\ Math.} 211 (2007), 566--593; B.~Feigin and A.~Stoyanovsky, arXiv:hep-th/9308079 (functional models).

\textbf{Physics.} C.~Beem, M.~Lemos, P.~Liendo, W.~Peelaers, L.~Rastelli and B.~C.~van~Rees, ``Infinite chiral symmetry in four
dimensions'', \emph{Comm.\ Math.\ Phys.} 336 (2015), 1359--1433, arXiv:1312.5344; C.~Beem and L.~Rastelli, ``Vertex operator
algebras, Higgs branches, and modular differential equations'', \emph{JHEP} 2018(8), 114, arXiv:1707.07679; J.~A.~Minahan and
D.~Nemeschansky, ``An $N=2$ superconformal fixed point with $E_6$ global symmetry'', \emph{Nucl.\ Phys.\ B} 482 (1996), 142--152.

\textbf{Project records.} \texttt{Notes/deligne/audit\_2026\_10\_25/DeligneCampaignAudit.tex} and
\texttt{B\_D\_uniform\_review.md} (chain B1--B8, input ledger U1--U7); \texttt{Notes/deligne/campaign/LEAN\_CAMPAIGN.md} and the
logs \texttt{WS-REL.md}, \texttt{WS-CAR.md}, \texttt{WS-ZHU.md}, \texttt{WS-AFF.md}.

\endgroup
